\chapter{Geodesics; Convex Neighborhoods}
\label{chap:dc-geodesics-convex-neighborhoods}

\section{The geodesic flow}

Throughout this chapter, $M$ is a Riemannian manifold with its Riemannian
connection.

\begin{definition}[geodesic]
  \label{def:dc-ch3-2-1}
  \lean{Riemannian.Geodesic.HasGeodesicEquationAt,
    Riemannian.Geodesic.hasGeodesicEquationAt_iff_covDerivAlong_velocity_eq_zero,
    Riemannian.Geodesic.continuousAt_and_hasGeodesicEquationAt_iff_leviCivita_covDerivAlong_velocity_eq_zero,
    Riemannian.Geodesic.IsGeodesicCurve,
    Riemannian.Geodesic.isGeodesicCurve_iff_covDerivAlong_velocity_eq_zero,
    Riemannian.Geodesic.isGeodesicCurve_iff_leviCivita_covDerivAlong_velocity_eq_zero,
    Riemannian.RiemannianMetric.leviCivita_covDerivAlong_eq}
  \leanok
  \dcref{ch3:2.1}
  A parametrized curve $\gamma:I\to M$ is a \emph{geodesic at} $t_0\in I$ if
  $\frac{D}{dt}\left(\frac{d\gamma}{dt}\right)=0$ at the point $t_0$; if $\gamma$ is
  a geodesic at $t$ for all $t\in I$, then $\gamma$ is a \emph{geodesic}. If
  $[a,b]\subset I$, the restriction of $\gamma$ to $[a,b]$ is a \emph{geodesic segment joining $\gamma(a)$ to $\gamma(b)$}. By abuse of language, the image $\gamma(I)$
  is also called a geodesic.
  If $\gamma$ is a geodesic then
  $$
  \frac{d}{dt}\Big\langle\frac{d\gamma}{dt},\frac{d\gamma}{dt}\Big\rangle=2\Big\langle\frac{D}{dt}\frac{d\gamma}{dt},\frac{d\gamma}{dt}\Big\rangle=0,
  $$
  so $\left|\frac{d\gamma}{dt}\right|=c\neq 0$ is constant (we exclude geodesics
  reducing to points). Arc length $s(t)=\int_{t_0}^t\left|\frac{d\gamma}{dt}\right|dt=c(t-t_0)$
  is then proportional to the parameter; when $c=1$ the geodesic is \emph{normalized}.
  In a coordinate system $(U,\mathbf{x})$ about $\gamma(t_0)$, a curve
  $\gamma(t)=(x_1(t),\dots,x_n(t))$ is a geodesic iff
  $$
  0=\frac{D}{dt}\Big(\frac{d\gamma}{dt}\Big)=\sum_k\Big(\frac{d^2x_k}{dt^2}+\sum_{i,j}\Gamma_{ij}^k\frac{dx_i}{dt}\frac{dx_j}{dt}\Big)\frac{\partial}{\partial x_k},
  $$
  i.e. iff the second-order system
  $$
  \frac{d^2x_k}{dt^2}+\sum_{i,j}\Gamma_{ij}^k\frac{dx_i}{dt}\frac{dx_j}{dt}=0,\qquad k=1,\dots,n,\qquad\text{(1)}
  $$
  holds.

\end{definition}

\begin{lemma}[geodesics have constant speed]
  \label{lem:dc-ch3-2-1-constspeed}
  \lean{Riemannian.hasDerivAt_chartMetricInner_geodesic_speed_zero, Riemannian.Geodesic.speedSq, Riemannian.Geodesic.IsGeodesicOn.hasDerivAt_speedSq_zero, Riemannian.Geodesic.IsGeodesicOn.speedSq_eq}
  \uses{def:dc-ch3-2-1, lem:dc-ch2-3-2-coord}
  \leanok
  \dcref{ch3:2.1}
  If $\gamma$ is a geodesic at $t$, then
  $\frac{d}{ds}\big\langle\gamma'(s),\gamma'(s)\big\rangle=0$ at $s=t$. Read in
  the chart at the foot $\gamma(t)$, with $u=\varphi_{\gamma(t)}\circ\gamma$ the
  chart representative and $\langle\cdot,\cdot\rangle$ the chart Gram inner
  product $g_{ij}(u)\,\cdot^i\,\cdot^j$, the squared speed
  $s\mapsto\langle u'(s),u'(s)\rangle$ has vanishing derivative at $s=t$.

  Intrinsically: the squared speed
  $t\mapsto\langle\gamma'(t),\gamma'(t)\rangle$ of the manifold velocity
  $\gamma'(t)\in T_{\gamma(t)}M$ has vanishing derivative at every time of an
  open set on which the (continuous) curve $\gamma$ is a geodesic, and is
  therefore constant on every such connected open set. Hence $|\gamma'|=c$ is
  constant along a geodesic (excluding geodesics reducing to points), and the
  arc length $s(t)=\int_{t_0}^t|\gamma'|=c(t-t_0)$ is proportional to the
  parameter.
\end{lemma}
\begin{proof}
  \sketch
  \leanok
  By metric compatibility in coordinate form
  (\cref{lem:dc-ch2-3-2-coord}),
  $$
  \frac{d}{ds}\big\langle u'(s),u'(s)\big\rangle
    =\Big\langle\tfrac{Du'}{ds},u'\Big\rangle+\Big\langle u',\tfrac{Du'}{ds}\Big\rangle,
  $$
  where $\frac{Du'}{ds}=u''+\Gamma(u',u')(u)$ is the covariant derivative of the
  velocity. The geodesic equation at $t$ is exactly $\frac{Du'}{ds}=0$ at $s=t$,
  so both inner products vanish and the derivative is $0$.

  For the intrinsic statement, fix a base time $t$ and read everything in the
  chart at $\gamma(t)$: for $\sigma$ near $t$ the chart Gram inner product of
  the chart velocity at $u(\sigma)$ equals
  $\langle\gamma'(\sigma),\gamma'(\sigma)\rangle$ at $\gamma(\sigma)$, because
  the Gram entries are by definition the inner products of the chart frame
  vectors, and the chart velocity taken at base time $\sigma$ transforms into
  the one taken at base time $t$ by the derivative of the smooth chart
  transition — exactly the change of frame between the two chart
  presentations of $T_{\gamma(\sigma)}M$. Hence the intrinsic squared speed
  agrees near $t$ with the chart expression above, so its derivative vanishes
  at $t$; a real function on a connected open set whose derivative vanishes
  everywhere is constant.
\end{proof}

\subsection{Tangent bundle}
$TM$ is the set of pairs $(q,v)$, $q\in M$, $v\in T_qM$. If $(U,\mathbf{x})$ is a
coordinate system on $M$, any $v\in T_qM$, $q\in\mathbf{x}(U)$, is
$\sum_i y_i\frac{\partial}{\partial x_i}$, and $(x_1,\dots,x_n,y_1,\dots,y_n)$ are
coordinates of $(q,v)$ on $TU$, giving $TM$ a differentiable structure (cf.
\cref{ex:dc-ch0-4-1} of Chap. 0). Locally $TU=U\times\mathbb{R}^n$, and the canonical
projection $\pi:TM\to M$, $\pi(q,v)=q$, is differentiable. A curve $t\mapsto\gamma(t)$
in $M$ determines $t\mapsto(\gamma(t),\frac{d\gamma}{dt}(t))$ in $TM$; if $\gamma$
is a geodesic, this curve satisfies the first-order system

$$
\begin{cases}\dfrac{dx_k}{dt}=y_k,\\[4pt]\dfrac{dy_k}{dt}=-\sum_{i,j}\Gamma_{ij}^k y_i y_j,\end{cases}\qquad k=1,\dots,n,\qquad\text{(1')}
$$

on $TU$. Thus the second-order system (1) on $U$ is equivalent to the first-order
system (1') on $TU$.

\begin{lemma}[trajectories of a vector field: existence and local uniqueness]
  \label{lem:dc-ch3-2-2-pointwise}
  \lean{exists_isMIntegralCurveAt_of_contMDiffAt_boundaryless, isMIntegralCurveAt_eventuallyEq_of_contMDiffAt_boundaryless}
  \dcref{ch3:2.2}
  \mathlibok
  Let $X$ be a $C^1$ vector field on a boundaryless manifold $M$ modelled on a
  complete normed space, and $p\in M$. Then there is a trajectory of $X$ through
  $p$ at $t=0$ (a curve $\gamma$ with $\gamma'=X(\gamma)$ defined near $0$), and
  any two trajectories of $X$ that agree at a time $t_0$ agree on a
  neighbourhood of $t_0$. (Picard--Lindel\"of and Gr\"onwall, transported to the
  manifold through a chart.)
\end{lemma}

\begin{lemma}[local flow with Lipschitz dependence on the initial condition]
  \label{lem:dc-ch3-2-2-flow}
  \lean{IsPicardLindelof.exists_forall_mem_closedBall_eq_hasDerivWithinAt_lipschitzOnWith, IsPicardLindelof.exists_forall_mem_closedBall_eq_hasDerivWithinAt_continuousOn, IsPicardLindelof.of_contDiffAt_one}
  \uses{lem:dc-ch3-2-2-pointwise}
  \dcref{ch3:2.2}
  \mathlibok
  Let $f$ be a $C^1$ vector field on a Banach space. Around any point there are
  $r > 0$, $\varepsilon > 0$ and a \emph{local flow}: a single family
  $Z(x, \cdot)$ of solutions of $z' = f(z)$, one for each initial condition $x$
  in the closed ball of radius $r$, all defined on the fixed time interval
  $[-\varepsilon, \varepsilon]$, and Lipschitz in $x$ uniformly in time (hence
  jointly continuous in $(x, t)$). This is the uniform-time and
  continuous-dependence content of \cref{thm:dc-ch3-2-2}; beyond it, that
  theorem asserts the joint $C^\infty$ regularity of the flow, whose $C^1$-in-%
  the-initial-condition part is \cref{lem:dc-ch3-2-2-c1dep}.
\end{lemma}

\begin{lemma}[$C^1$ dependence on the initial condition along a base solution]
  \label{lem:dc-ch3-2-2-c1dep}
  \lean{Riemannian.FlowDependence.hasStrictFDerivAt_of_picardResidual_curve, Riemannian.FlowDependence.exists_hasStrictFDerivAt_of_picardResidual_curve, Riemannian.FlowDependence.hasStrictFDerivAt_eval_of_picardResidual_curve, Riemannian.FlowDependence.hasStrictFDerivAt_superposition}
  \uses{lem:dc-ch3-2-2-flow}
  \leanok
  \dcref{ch3:2.2}
  Let $f$ be a vector field on a Banach space $E$, differentiable on an open
  set $u$ with derivative $f'$, and let $\alpha_0 : [0,T] \to u$ be a solution
  of $\alpha' = f(\alpha)$ with initial value $x_0$, with $f'$ continuous at
  every point of $\alpha_0$. Write $A_0(t) = f'(\alpha_0(t))$ for the operator
  curve along the solution and assume $T \sup_t \|A_0(t)\| < 1$. If $\sigma$
  assigns to every initial value $x$ near $x_0$ a curve
  $\sigma(x) \in C([0,T], E)$ solving the integral equation
  $\sigma(x)(t) = x + \int_0^t f(\sigma(x)(s))\,ds$, with
  $\sigma(x_0) = \alpha_0$ and $\sigma$ continuous at $x_0$, then $\sigma$ is
  strictly Fr\'echet differentiable at $x_0$ as a map $E \to C([0,T], E)$. Its
  derivative $D$ exists and is characterized by the linearized (variational)
  integral equation
  $$
  (Dv)(t) - \int_0^t A_0(s)\,\big((Dv)(s)\big)\,ds = v, \qquad v \in E,
  $$
  and for each fixed time $t$ the time-$t$ map $x \mapsto \sigma(x)(t)$ is
  strictly differentiable at $x_0$ with derivative $v \mapsto (Dv)(t)$. This is
  the $C^1$ dependence of the flow on its initial condition at an arbitrary
  (non-equilibrium) base solution, the pointwise-$C^1$ core of
  \cref{thm:dc-ch3-2-2}.
\end{lemma}
\begin{proof}
  \sketch
  \leanok
  A curve $\alpha$ solves the ODE with initial value $x$ iff the Picard
  residual $G(x, \alpha) = \alpha - x - \int_0^t f(\alpha)$ vanishes, viewed as
  a map $E \times C([0,T], E) \to C([0,T], E)$ of Banach spaces. The only
  nonlinearity of $G$ is the superposition operator
  $\alpha \mapsto f \circ \alpha$; its strict derivative at the base curve
  $\alpha_0$ is the postcomposition operator
  $\beta \mapsto (t \mapsto A_0(t)\,\beta(t))$. Indeed, for curves
  $\alpha, \beta$ uniformly within $\delta$ of $\alpha_0$, the pointwise error
  $f(\alpha(t)) - f(\beta(t)) - f'(\alpha_0(t))(\alpha(t) - \beta(t))$ is at
  most $\varepsilon\,\|\alpha(t) - \beta(t)\|$ by the mean value inequality
  applied to $z \mapsto f(z) - f'(\alpha_0(t))z$ on the ball
  $B_\delta(\alpha_0(t))$; the estimate holds simultaneously for all $t$
  because the range of $\alpha_0$ is compact, so a uniform $\delta$-thickening
  of the range stays inside $u$ and $f'$ is uniformly continuous at the range
  (Heine: continuity at each point of a compact set is uniform against nearby
  points, even ones outside the set). Hence $G$ is strictly differentiable at
  $(x_0, \alpha_0)$, with partial derivative in the curve variable
  $1 - J \circ M$, where $J$ is the Volterra primitive
  ($\|J\| \le T$) and $M$ the postcomposition by $A_0$
  ($\|M\| \le \sup_t \|A_0(t)\|$). Since $T \sup_t \|A_0(t)\| < 1$, the
  operator $1 - J \circ M$ is invertible by the Neumann series, and the Banach
  implicit function theorem applies to $G$ at $(x_0, \alpha_0)$: the local zero
  family is strictly differentiable at $x_0$, and it coincides with $\sigma$
  near $x_0$ because $\sigma$ is continuous at $x_0$ and solves
  $G(x, \sigma(x)) = 0$. The derivative equation
  $(1 - J \circ M) \circ D = \mathrm{const}$ is exactly the variational
  integral equation, solved by $D = (1 - J \circ M)^{-1} \circ \mathrm{const}$
  and determining $D$ uniquely; postcomposing with the bounded evaluation at a
  fixed time $t$ gives the time-$t$ statement.
\end{proof}

\begin{lemma}[the variational flow computes the derivative of the flow]
  \label{lem:dc-ch3-2-2-variational}
  \lean{Riemannian.FlowDependence.flowDeriv_eq_applyCurve_opFlow, Riemannian.FlowDependence.variationalField, Riemannian.FlowDependence.picardResidual_variationalField_compRight, Riemannian.FlowDependence.applyCurve_sub_intervalPrimitive_of_picardResidual_variationalField}
  \uses{lem:dc-ch3-2-2-c1dep}
  \leanok
  \dcref{ch3:2.2}
  Let $f$ be a vector field on a Banach space $E$, twice differentiable on an
  open set $u$, and let $\alpha_0 : [0,T] \to u$ solve $\alpha' = f(\alpha)$
  with $T \sup_t \|f'(\alpha_0(t))\| < 1$. Consider the \emph{extended
  (variational) field} on $E \times L(E,E)$,
  $$\hat f(z, W) = \big(f(z),\; f'(z) \circ W\big),$$
  whose trajectories pair a trajectory of $f$ with the solution of the
  linearized equation $W' = f'(\alpha(t)) \circ W$ along it. If $(\alpha_0, W)$
  solves the integral equation of $\hat f$ with initial condition $(x_0, 1)$,
  then the derivative $D$ of the solution family of $f$ at $x_0$
  (\cref{lem:dc-ch3-2-2-c1dep}) is computed by the operator flow:
  $Dv = (t \mapsto W(t)\,v)$ for every $v \in E$. Moreover right composition
  scales the operator slot: $(\alpha_0, W \circ W_0)$ solves the extended
  equation with initial condition $(x_0, W_0)$, for every $W_0 \in L(E,E)$.
\end{lemma}
\begin{proof}
  \sketch
  \leanok
  Applying the second component of the extended integral equation to a fixed
  vector $v$, and commuting the (bounded linear) evaluation with the Volterra
  primitive, shows that the curve $t \mapsto W(t)v$ satisfies
  $$W(t)v - \int_0^t f'(\alpha_0(s))\big(W(s)v\big)\,ds = v ,$$
  which is exactly the variational integral equation
  $(1 - J \circ M)\beta = \mathrm{const}\, v$ characterizing $Dv$
  (\cref{lem:dc-ch3-2-2-c1dep}). Since $T \sup_t \|f'(\alpha_0(t))\| < 1$, the
  operator $1 - J \circ M$ is invertible by the Neumann series, so the two
  solutions coincide: $Dv = (t \mapsto W(t)v)$. For the scaling clause, the map
  $\psi(z, W) = (z, W \circ W_0)$ is linear and bounded, commutes with the
  extended field ($\hat f \circ \psi = \psi \circ \hat f$, by associativity of
  composition), with the constant embedding ($\psi(x_0, 1) = (x_0, W_0)$), and
  with the primitive; hence it maps zeros of the extended residual with initial
  condition $(x_0, 1)$ to zeros with initial condition $(x_0, W_0)$.
\end{proof}

\begin{lemma}[second-order dependence on the initial condition]
  \label{lem:dc-ch3-2-2-c2dep}
  \lean{Riemannian.FlowDependence.exists_hasStrictFDerivAt_opFlow_of_picardResidual_curve, Riemannian.FlowDependence.hasFDerivAt_variationalField, Riemannian.FlowDependence.variationalFieldDeriv}
  \uses{lem:dc-ch3-2-2-c1dep, lem:dc-ch3-2-2-variational}
  \leanok
  \dcref{ch3:2.2}
  In the setting of \cref{lem:dc-ch3-2-2-variational}, suppose additionally
  that for every initial value $x$ near $x_0$ the operator flow
  $\tau(x) \in C([0,T], L(E,E))$ solves the extended integral equation along
  $\sigma(x)$ with initial condition $(x, 1)$, that $\tau$ is continuous at
  $x_0$, and that the linearization $\hat A_0$ of the extended field along the
  base pair $(\alpha_0, \tau(x_0))$ satisfies $T \sup_t \|\hat A_0(t)\| < 1$.
  Then the family of variational flows $x \mapsto \tau(x)$ is strictly
  Fr\'echet differentiable at $x_0$. Combined with
  \cref{lem:dc-ch3-2-2-variational} — the first derivative of the flow is
  computed by $\tau$ — this is second-order differentiability of the flow in
  its initial condition, the $C^2$ step of the joint regularity asserted by
  \cref{thm:dc-ch3-2-2}.
\end{lemma}
\begin{proof}
  \sketch
  \leanok
  The extended field $\hat f(z, W) = (f(z), f'(z) \circ W)$ is differentiable
  on $u \times L(E,E)$ with derivative
  $$D\hat f_{(z,W)}(u', V) = \big(f'(z)u',\; f'(z) \circ V
    + (f''(z)u') \circ W\big),$$
  continuous wherever $f'$ and $f''$ are continuous: the first slot is the
  chain rule, and the second is the derivative of the bounded bilinear
  composition. The pairs $\hat\sigma(x, W_0) = (\sigma(x), \tau(x) \circ W_0)$
  form a solution family of $\hat f$ for \emph{all} initial conditions
  $(x, W_0)$ near $(x_0, 1)$ — right composition scales the operator slot
  (\cref{lem:dc-ch3-2-2-variational}) — and $\hat\sigma$ is continuous at
  $(x_0, 1)$, since $(W_0, \gamma) \mapsto \gamma \circ W_0$ is bounded
  bilinear. The trajectories $\sigma(x)$ stay in $u$ for $x$ near $x_0$: the
  compact range of $\alpha_0$ has a uniform thickening inside $u$, and
  $\sigma$ is continuous at $x_0$ in the sup norm. Since
  $T \sup_t \|\hat A_0(t)\| < 1$, the $C^1$-dependence theorem
  (\cref{lem:dc-ch3-2-2-c1dep}, applied on the Banach space
  $E \times L(E,E)$) makes $\hat\sigma$ strictly differentiable at $(x_0, 1)$;
  composing with the affine embedding $x \mapsto (x, 1)$ and the bounded
  projection to the operator component yields the strict differentiability of
  $\tau$ at $x_0$.
\end{proof}

\begin{theorem}[flow of a vector field]
  \label{thm:dc-ch3-2-2}
  \uses{lem:dc-ch3-2-2-pointwise, lem:dc-ch3-2-2-flow, lem:dc-ch3-2-2-c1dep, lem:dc-ch3-2-2-variational, lem:dc-ch3-2-2-c2dep, lem:dc-ch3-2-2-cinfty}
  \notready
  \dcref{ch3:2.2}
  \notready
  If $X$ is a $C^\infty$ vector field on the open set $V$ in the manifold $M$ and
  $p\in V$, then there exist an open set $V_0\subset V$ with $p\in V_0$, a number
  $\delta>0$, and a $C^\infty$ mapping $\varphi:(-\delta,\delta)\times V_0\to V$ such
  that $t\mapsto\varphi(t,q)$, $t\in(-\delta,\delta)$, is the unique trajectory of
  $X$ passing through $q$ at $t=0$, for every $q\in V_0$. The map
  $\varphi_t:V_0\to V$, $\varphi_t(q)=\varphi(t,q)$, is the \emph{flow} of $X$ on $V$.
  (Beyond \cref{lem:dc-ch3-2-2-pointwise}, this asserts a uniform time of
  existence and the joint $C^\infty$ regularity of $\varphi$ in $(t,q)$ — the
  smooth dependence of solutions of an ODE on their initial conditions. The
  uniform time and Lipschitz dependence is \cref{lem:dc-ch3-2-2-flow}; the
  pointwise $C^1$ dependence on the initial condition, at an arbitrary base
  solution, is \cref{lem:dc-ch3-2-2-c1dep}; the second-order dependence, given
  the variational flow along the base solution, is \cref{lem:dc-ch3-2-2-c2dep}
  via \cref{lem:dc-ch3-2-2-variational}; the construction of the variational
  flow for the geodesic spray and the joint $C^k$ regularity for $k \ge 3$
  remain open.)
\end{theorem}

\begin{lemma}[the geodesic field]
  \label{lem:dc-ch3-2-3}
  \lean{Riemannian.Geodesic.existsUnique_geodesicField}
  \uses{def:dc-ch3-2-1}
  \leanok
  \dcref{ch3:2.3}
  \notready
  There exists a unique vector field $G$ on $TM$ whose trajectories are of the form
  $t\mapsto(\gamma(t),\gamma'(t))$, where $\gamma$ is a geodesic on $M$.
\end{lemma}
\begin{proof}
  \leanok
  \sketch
  Uniqueness: in a coordinate system $(U,\mathbf{x})$ the trajectories of
  $G$ on $TU$ are $t\mapsto(\gamma(t),\gamma'(t))$ with $\gamma$ a geodesic, hence
  solutions of (1'); by uniqueness of trajectories of such a system, $G$ is unique
  if it exists. Existence: define $G$ locally by (1'); uniqueness shows $G$ is
  well-defined on all of $TM$.
\end{proof}

\begin{definition}[geodesic field / geodesic flow]
  \label{def:dc-ch3-2-4}
  \lean{Riemannian.Geodesic.geodesicVectorField, Riemannian.Geodesic.geodesicVectorFieldChart, Riemannian.Geodesic.geodesicVectorFieldChart_contMDiffOn}
  \uses{thm:dc-ch3-2-2, lem:dc-ch3-2-3}
  \leanok
  \dcref{ch3:2.4}
  The vector field $G$ above is the \emph{geodesic field} on $TM$, and its flow is the
  \emph{geodesic flow} on $TM$. Applying \cref{thm:dc-ch3-2-2} to $G$ at $(p,0)\in TM$: for each
  $p\in M$ there exist an open set $\mathcal{U}\subset TU$ with $(p,0)\in\mathcal{U}$,
  a number $\delta>0$, and a $C^\infty$ map $\varphi:(-\delta,\delta)\times\mathcal{U}\to TU$
  such that $t\mapsto\varphi(t,q,v)$ is the unique trajectory of $G$ with
  $\varphi(0,q,v)=(q,v)$. One may take
  $$
  \mathcal{U}=\{(q,v)\in TU;\ q\in V,\ v\in T_qM,\ |v|<\varepsilon_1\},
  $$
  $V\subset U$ a neighborhood of $p$. Putting $\gamma=\pi\circ\varphi$, where
  $\pi:TM\to M$ is the canonical projection, we can describe the previous result as
  follows.
\end{definition}

\begin{lemma}[pointwise existence of a geodesic with prescribed initial velocity]
  \label{lem:dc-ch3-2-5-exist}
  \lean{Riemannian.Geodesic.exists_geodesic_with_initial_velocity_at}
  \uses{def:dc-ch3-2-1, lem:dc-ch3-2-2-pointwise}
  \leanok
  \dcref{ch3:2.5}
  For every $p\in M$ and every $v\in T_pM$ there is a curve $\gamma:\mathbb{R}\to M$
  with $\gamma(0)=p$ that is a local geodesic at $t=0$ and whose canonical tangent
  lift $t\mapsto(\gamma(t),\gamma'(t))\in TM$ takes the value $(p,v)$ at $t=0$; i.e.
  $\gamma$ is a geodesic through $p$ with initial velocity $v$.
\end{lemma}
\begin{proof}
  \sketch
  \leanok
  Read the geodesic equation on $TM$ as the first-order system (1') for the
  geodesic spray $G$. The chart-fixed spray is $C^\infty$ at $(p,v)$, so the
  integral-curve existence theorem (Picard--Lindel\"of,
  \cref{lem:dc-ch3-2-2-pointwise}) produces a local integral curve $f$ of $G$ with
  $f(0)=(p,v)$; its base projection $\gamma=\pi\circ f$ is the desired geodesic.
\end{proof}

\begin{lemma}[pointwise uniqueness of geodesics]
  \label{lem:dc-ch3-2-5-unique}
  \lean{Riemannian.Geodesic.isGeodesicAt_eventuallyEq_of_lift_eq}
  \uses{def:dc-ch3-2-1, lem:dc-ch3-2-2-pointwise}
  \leanok
  \dcref{ch3:2.5}
  If two curves $\gamma_1,\gamma_2$ are base projections of integral curves of the
  geodesic spray whose lifts agree at $t_0$ (equal initial point \emph{and} initial
  velocity), then $\gamma_1$ and $\gamma_2$ agree on a neighbourhood of $t_0$.
\end{lemma}
\begin{proof}
  \sketch
  \leanok
  The lifts are integral curves of the same $C^\infty$ vector field $G$ on $TM$
  with the same value at $t_0$, so by uniqueness of integral curves
  (\cref{lem:dc-ch3-2-2-pointwise}, Gr\"onwall) they agree near $t_0$; projecting to $M$
  gives $\gamma_1=\gamma_2$ near $t_0$.
\end{proof}

\begin{lemma}[a geodesic leaves its initial point with its initial velocity]
  \label{lem:dc-ch3-2-5-velocity}
  \lean{Riemannian.Geodesic.IsGeodesicOnWithInitial.hasDerivAt_extChartAt_zero, Riemannian.Geodesic.hasDerivAt_extChartAt_maximalGeodesic}
  \uses{lem:dc-ch3-2-5-exist, lem:dc-ch3-2-5-unique}
  \leanok
  \dcref{ch3:2.5}
  A geodesic with initial data $(p,v)$ really has velocity $v$ at time $0$: read
  in the chart at $p$, the curve $s\mapsto\varphi_p(\gamma(s))$ is differentiable
  at $s=0$ with derivative $v$. The same holds for the canonical maximal
  geodesic $\gamma(\cdot,p,v)$.
\end{lemma}
\begin{proof}
  \sketch
  \leanok
  The velocity lift $f$ of the geodesic is an integral curve of the geodesic
  field with $f(0)=(p,v)$. Reading the integral-curve equation in the chart of
  $TM$ at $(p,v)$ — which is based at the foot $p$ — the pair
  $(x,w)=(\varphi_p\circ\gamma,\ \text{fibre coordinate of }f)$ satisfies the
  first-order system $(x',w')=(w,\,-\Gamma_p(w,w)(x))$ near $0$; in particular
  $x'(0)=w(0)=v$, since the trivialisation at $p$ is the identity on the fibre
  over $p$. For the canonical maximal geodesic, local uniqueness
  (\cref{lem:dc-ch3-2-5-unique}) propagated along the connected witness interval
  identifies it with a single witness near $0$.
\end{proof}

\begin{lemma}[a geodesic satisfies the geodesic equation at the initial time]
  \label{lem:dc-ch3-2-5-equation}
  \lean{Riemannian.Geodesic.IsGeodesicOnWithInitial.hasGeodesicEquationAt_zero, Riemannian.Geodesic.hasGeodesicEquationAt_maximalGeodesic_zero}
  \uses{def:dc-ch3-2-1, lem:dc-ch3-2-5-exist, lem:dc-ch3-2-5-velocity}
  \leanok
  \dcref{ch3:2.5}
  A geodesic with initial data $(p,v)$ — the base projection of an integral
  curve of the geodesic field — satisfies the second-order geodesic equation of
  \cref{def:dc-ch3-2-1} at the initial time: in the chart at the foot
  $\gamma(0)=p$, the chart curve $x=\varphi_p\circ\gamma$ has a second
  derivative at $0$ and $x''(0)+\Gamma_p(v,v)(x(0))=0$. The same holds for the
  canonical maximal geodesic. This bridges the first-order (geodesic-field)
  formulation back to the second-order equation at the initial time.
\end{lemma}
\begin{proof}
  \sketch
  \leanok
  With $(x,w)$ as in \cref{lem:dc-ch3-2-5-velocity}, the first-order system
  gives $x'=w$ near $0$ and $w'(0)=-\Gamma_p(w(0),w(0))(x(0))$; hence $x'$ is
  differentiable at $0$ with
  $x''(0)=w'(0)=-\Gamma_p(v,v)(x(0))$, which is the geodesic equation at $0$.

\end{proof}

\begin{proposition}[local existence/uniqueness of geodesics]
  \label{prop:dc-ch3-2-5}
  \lean{Riemannian.GeodesicLocal.geodesic_local_existence_metricBall}
  \uses{lem:dc-ch3-2-5-exist, lem:dc-ch3-2-5-unique, lem:dc-ch3-2-5-chart-family, lem:dc-ch3-2-5-metric-family}
  \dcref{ch3:2.5}
  Given $p\in M$, there exist an open set $V\subset M$ with $p\in V$, numbers
  $\delta>0$, $\varepsilon_1>0$, and a $C^\infty$ map
  $$
  \gamma:(-\delta,\delta)\times\mathcal{U}\to M,\quad\mathcal{U}=\{(q,v);\ q\in V,\ v\in T_qM,\ |v|<\varepsilon_1\},
  $$
  such that $t\mapsto\gamma(t,q,v)$, $t\in(-\delta,\delta)$, is the unique geodesic
  passing through $q$ with velocity $v$ at $t=0$, for each $q\in V$ and each
  $v\in T_qM$ with $|v|<\varepsilon_1$.
\end{proposition}

\begin{lemma}[affine rescaling of a geodesic]
  \label{lem:dc-ch3-2-6-reparam}
  \lean{Riemannian.Geodesic.isGeodesic_comp_mul_left}
  \leanok
  \dcref{ch3:2.6}
  If $\gamma:\mathbb{R}\to M$ is a geodesic, then for every $a\in\mathbb{R}$ the
  affine reparametrisation $h(t)=\gamma(at)$ is again a geodesic, with velocity
  $h'(t)=a\,\gamma'(at)$.
\end{lemma}
\begin{proof}
  \sketch
  \leanok
  Read the geodesic equation in a chart at the moving foot $\gamma(at)$: with
  $u(s)=\mathbf{x}(\gamma(s))$ the chart representative, $h$ reads as $s\mapsto
  u(as)$, so $h'(t)=a\,u'(at)$ and $h''(t)=a^2\,u''(at)$. Since the Christoffel
  contraction $\Gamma$ scales as a $(2,0)$-tensor,
  $\Gamma(au'(at),au'(at))=a^2\,\Gamma(u'(at),u'(at))$, hence
  $$
  h''(t)+\Gamma\big(h'(t),h'(t)\big)
    =a^2\Big(u''(at)+\Gamma\big(u'(at),u'(at)\big)\Big)=a^2\cdot 0=0,
  $$
  so $h$ satisfies the geodesic equation at every $t$.
\end{proof}

\begin{lemma}[homogeneity of a geodesic]
  \label{lem:dc-ch3-2-6}
  \lean{Riemannian.Geodesic.IsGeodesicOnWithInitial.fiberScale, Riemannian.Geodesic.maximalGeodesicInterval_fiberScale, Riemannian.Geodesic.maximalGeodesic_fiberScale}
  \uses{lem:dc-ch3-2-6-reparam, lem:dc-ch3-2-5-unique}
  \leanok
  \dcref{ch3:2.6}
  If the geodesic $\gamma(t,q,v)$ is defined on $(-\delta,\delta)$, then for
  $a\in\mathbb{R}$, $a>0$, the geodesic $\gamma(t,q,av)$ is defined on
  $(-\frac{\delta}{a},\frac{\delta}{a})$ and

  $$
  \gamma(t,q,av)=\gamma(at,q,v).
  $$
\end{lemma}
\begin{proof}
  \sketch
  \leanok
  Let $h:(-\frac{\delta}{a},\frac{\delta}{a})\to M$, $h(t)=\gamma(at,q,v)$.
  Then $h(0)=q$, $\frac{dh}{dt}(0)=av$, and by \cref{lem:dc-ch3-2-6-reparam}
  $h$ is a geodesic (with $h'(t)=a\gamma'(at,q,v)$), so

  $$
  \frac{D}{dt}\Big(\frac{dh}{dt}\Big)=\nabla_{h'(t)}h'(t)=a^2\nabla_{\gamma'(at,q,v)}\gamma'(at,q,v)=0.
  $$

  So $h$ is a geodesic through $q$ with velocity $av$ at $t=0$; by uniqueness
  $h(t)=\gamma(at,q,v)=\gamma(t,q,av)$.
\end{proof}

\begin{lemma}[pointwise uniform interval of definition]
  \label{lem:dc-ch3-2-7-pointwise}
  \lean{Riemannian.Geodesic.exists_pos_smul_maximalGeodesicInterval}
  \uses{lem:dc-ch3-2-6}
  \leanok
  \dcref{ch3:2.7}
  For every $p\in M$, $v\in T_pM$ and every $T>0$ there is a $c>0$ such that
  the geodesic with initial data $(p,\,c\,v)$ is defined on $(-T,T)$. This is
  the pointwise-in-$(p,v)$ content of \cref{prop:dc-ch3-2-7}: shrinking the
  velocity extends the interval of definition. (The uniformity in the velocity
  over a ball $|v|\le r$ is now closed by \cref{lem:dc-ch3-2-7-vuniform}; the
  uniformity in the base point $q \in V$ together with the joint $C^\infty$
  regularity of the family — the smooth-family clause of
  \cref{prop:dc-ch3-2-7} — remains open, pending smooth dependence of ODE
  solutions on initial data.)
\end{lemma}
\begin{proof}
  \sketch
  \leanok
  The maximal interval $I_{\max}(p,v)$ is open and contains $0$
  (\cref{lem:dc-ch3-2-5-exist}), so it contains some $(-\delta,\delta)$,
  $\delta>0$. By the homogeneity of geodesics (\cref{lem:dc-ch3-2-6}),
  $I_{\max}(p,\,c\,v)=\{t:\ c\,t\in I_{\max}(p,v)\}$; taking
  $c=\frac{\delta}{2T}$ sends $(-T,T)$ inside $(-\delta/2,\delta/2)\subset
  I_{\max}(p,v)$, so $(-T,T)\subset I_{\max}(p,\,c\,v)$.
\end{proof}

\begin{lemma}[uniform interval in the velocity]
  \label{lem:dc-ch3-2-7-vuniform}
  \lean{Riemannian.Geodesic.exists_uniform_geodesic_flow, Riemannian.Geodesic.isGeodesicOnWithInitial_of_hasDerivAt_sprayCoord, Riemannian.Geodesic.maximalGeodesic_eq_witness_of_mem_chart, Riemannian.Geodesic.contDiffOn_geodesicSprayCoord_prod}
  \uses{lem:dc-ch3-2-2-flow, def:dc-ch3-2-4, lem:dc-ch3-2-5-exist, lem:dc-ch3-2-5-unique}
  \leanok
  \dcref{ch3:2.7}
  For each $p \in M$ there are $r > 0$, $\varepsilon > 0$ and a local flow $Z$
  of the coordinate geodesic spray at $p$ such that for \emph{every} velocity
  $v \in T_pM$ with $|v| \le r$: the geodesic $\gamma(\cdot, p, v)$ is defined
  on $(-\varepsilon, \varepsilon)$ — one interval for all such $v$ — and its
  chart reading is computed by the flow,
  $\varphi_p(\gamma(s, p, v)) = Z((\varphi_p(p), v), s)_1$; the flow is
  Lipschitz in $v$ uniformly in time, and the zero-velocity flow line is the
  constant $p$. This is the uniform-in-velocity clause of
  \cref{prop:dc-ch3-2-7} at a fixed base point.
\end{lemma}
\begin{proof}
  \sketch
  \leanok
  Apply the Picard--Lindelöf local flow (\cref{lem:dc-ch3-2-2-flow}) to the
  coordinate spray $F(x, w) = (w, -\Gamma_p(w, w)(x))$, which is $C^\infty$ on
  the chart target; trajectories are confined to the chart by Lipschitz
  dependence on the initial condition together with continuity of the centre
  trajectory (no compactness is needed). Each flow line, read through the
  tangent-bundle chart at $(p, 0)$, descends to a geodesic witness with foot in
  the chart at $p$ — the reverse of the chart reading of the integral-curve
  equation. Finally, the canonical maximal geodesic agrees with \emph{any}
  witness whose own foot stays in the chart: the clopen uniqueness propagation
  constrains only the given witness, so no chart-validity hypothesis about all
  witnesses is required.
\end{proof}

\begin{lemma}[the geodesic local flow is $C^1$ in the initial condition]
  \label{lem:dc-ch3-2-7-c1flow}
  \lean{Riemannian.Geodesic.exists_uniform_geodesic_flow_hasStrictFDerivAt}
  \uses{lem:dc-ch3-2-2-c1dep, lem:dc-ch3-2-7-vuniform, def:dc-ch3-2-4}
  \leanok
  \dcref{ch3:2.7}
  For each $p \in M$ there are $r, \varepsilon > 0$, a local flow $Z$ of the
  coordinate geodesic spray at $p$ with the properties of
  \cref{lem:dc-ch3-2-7-vuniform} (one time interval $[-\varepsilon,
  \varepsilon]$ for all initial conditions in the flow ball of radius $r$,
  Lipschitz dependence on the initial condition, chart confinement, and
  computation of the canonical maximal geodesic by the flow), and a time
  $0 < T < \varepsilon$, such that the flow — read as the family
  $\sigma : x \mapsto Z(x, \cdot)|_{[0,T]} \in C([0,T], \mathbb{R}^{2n})$ of
  its solution curves — is strictly Fr\'echet differentiable in the initial
  condition at \emph{every} $x_0$ in the open flow ball, equilibrium or not,
  with derivative characterized by the linearized (variational) integral
  equation of \cref{lem:dc-ch3-2-2-c1dep} along the base trajectory
  $\sigma(x_0)$. This is the pointwise-$C^1$ part, in the initial condition,
  of the $C^\infty$ family clause of \cref{prop:dc-ch3-2-7}.
\end{lemma}
\begin{proof}
  \sketch
  \leanok
  Trajectories starting in the flow ball are confined — by Lipschitz
  dependence on the initial condition together with constancy of the centre
  trajectory — to a ball whose closure is a compact subset of the chart
  region. The spray derivative $dF$ is continuous on the chart region, hence
  bounded by some $C \ge 0$ on the compact confinement ball; consequently the
  operator curve $A_0(t) = dF(Z(x_0, t))$ of \emph{every} base trajectory has
  sup norm at most $C$. Choose $T < \min(\varepsilon, 1/C)$; then
  $T \sup_t \|A_0(t)\| \le T C < 1$ simultaneously for all initial conditions
  in the flow ball. At each $x_0$ the family $\sigma$ satisfies the
  hypotheses of \cref{lem:dc-ch3-2-2-c1dep}: each flow line solves the Picard
  integral equation of the spray (its derivative is the spray value, and the
  values stay in the chart region where the spray is continuous), and
  $\sigma$ is Lipschitz — hence continuous — in the initial condition. The
  conclusion of that lemma is the claimed strict differentiability with the
  variational characterization.
\end{proof}

\begin{lemma}[the geodesic local flow is $C^2$ in the initial condition]
  \label{lem:dc-ch3-2-7-c2flow}
  \lean{Riemannian.Geodesic.exists_uniform_geodesic_flow_hasStrictFDerivAt_opFlow, Riemannian.FlowDependence.picardResidual_variationalField_eq_zero, Riemannian.FlowDependence.norm_variationalFieldDeriv_le}
  \uses{lem:dc-ch3-2-7-c1flow, lem:dc-ch3-2-2-variational, lem:dc-ch3-2-2-c2dep, def:dc-ch3-2-4}
  \leanok
  \dcref{ch3:2.7}
  For each $p \in M$ there are $r, \varepsilon > 0$, a local flow $Z$ of the
  coordinate geodesic spray at $p$ with the properties of
  \cref{lem:dc-ch3-2-7-vuniform}, a time $0 < T < \varepsilon$, the flow read
  as the family $\sigma : x \mapsto Z(x, \cdot)|_{[0,T]}$ of its solution
  curves, and the family $\tau : x \mapsto W_x \in C([0,T],
  L(\mathbb{R}^{2n}))$ of \emph{variational (operator) flows} along the
  trajectories — $W_x$ solves $W' = dF(\sigma(x)(t)) \circ W$, $W(0) = 1$ —
  such that at \emph{every} $x_0$ in the open flow ball:
  \begin{itemize}
  \item $\sigma$ is strictly Fr\'echet differentiable at $x_0$ and its
    derivative \emph{is computed by the variational flow}:
    $D\sigma(x_0)\,v = \big(t \mapsto \tau(x_0)(t)\,v\big)$
    (\cref{lem:dc-ch3-2-2-variational});
  \item the variational-flow family $\tau$ is itself strictly Fr\'echet
    differentiable at $x_0$ (\cref{lem:dc-ch3-2-2-c2dep}).
  \end{itemize}
  Together these say the local geodesic flow is $C^2$ in its initial
  condition; this is the $C^2$ part, in the initial condition, of the
  $C^\infty$ family clause of \cref{prop:dc-ch3-2-7}.
\end{lemma}
\begin{proof}
  \sketch
  \leanok
  As in \cref{lem:dc-ch3-2-7-c1flow}, trajectories from the flow ball are
  confined to a compact ball inside the chart region, on which the first and
  second derivatives of the spray are bounded, say $\|dF\| \le C$ and
  $\|d^2F\| \le C_2$. Choose the Picard time
  $T < \min\big(\varepsilon,\ \tfrac{1}{2(C + 2C_2 + 1)}\big)$, so that the
  contraction estimates below hold simultaneously for every initial condition
  in the flow ball.

  \emph{Construction of $\tau$.} For each $x$, the variational equation is
  the linear Volterra fixed-point equation $W = 1 + J_x W$ on $C([0,T],
  L(\mathbb{R}^{2n}))$, where $(J_x W)(t) = \int_0^t dF(\sigma(x)(s)) \circ
  W(s)\, ds$. Since $\|J_x\| \le T\,C \le \tfrac12 < 1$, the Neumann series
  inverts $1 - J_x$ and we set $\tau(x) = (1 - J_x)^{-1}(\mathbf{1})$. This
  single formula yields existence, the fixed-point identity
  $\tau(x) = \mathbf{1} + J_x \tau(x)$ (hence the uniform bound
  $\|\tau(x)\| \le (1 - \tfrac12)^{-1} = 2$), \emph{and} continuity of
  $x \mapsto \tau(x)$: the coefficient curve $t \mapsto dF(\sigma(x)(t))$
  depends continuously on $x$ (uniform continuity of $dF$ on the compact
  confinement ball together with Lipschitz dependence of the flow on its
  initial condition), $J_x$ depends continuously — indeed
  $\|J_x - J_{x'}\| \le T \sup_t \|dF(\sigma(x)(t)) - dF(\sigma(x')(t))\|$ —
  and inversion is continuous at units of a Banach algebra.

  \emph{The extended system.} The pair $(\sigma(x), \tau(x))$ is then a
  solution of the Picard integral equation of the extended (variational)
  field $\hat F(z, W) = (F(z),\ dF(z) \circ W)$ with initial condition
  $(x, 1)$: the residual splits into its two components, the first being the
  Picard residual of the base flow and the second the Volterra fixed-point
  identity above.

  \emph{Conclusion.} At each $x_0$ in the open flow ball, the derivative
  characterization of \cref{lem:dc-ch3-2-7-c1flow} and the uniqueness of the
  linearized fixed point identify $D\sigma(x_0)\,v$ with
  $t \mapsto \tau(x_0)(t)\,v$ (\cref{lem:dc-ch3-2-2-variational}). For the
  second derivative, apply the abstract second-order dependence theorem
  (\cref{lem:dc-ch3-2-2-c2dep}) to the extended system: its linearization
  along the base pair $(\sigma(x_0), \tau(x_0))$ is bounded by
  $C + C_2 \|\tau(x_0)\| \le C + 2C_2$ (the block estimate
  $\|D\hat F(z, W)\| \le \|dF(z)\| + \|d^2F(z)\|\,\|W\|$), so
  $T\,(C + 2C_2) < 1$ makes its hypotheses hold, and the family
  $x \mapsto \tau(x)$ is strictly differentiable at $x_0$.
\end{proof}

\begin{proposition}[uniform interval of definition]
  \label{prop:dc-ch3-2-7}
  \lean{Riemannian.GeodesicLocal.geodesic_local_existence_fixedInterval_metricBall}
  \uses{prop:dc-ch3-2-5, lem:dc-ch3-2-6, lem:dc-ch3-2-7-vuniform, lem:dc-ch3-2-7-c1flow, lem:dc-ch3-2-7-c2flow, lem:dc-ch3-2-5-unique, lem:dc-ch3-2-5-metric-ball}
  \leanok
  \dcref{ch3:2.7}
  \notready
  Given $p\in M$, there exist a neighborhood $V$ of $p$ in $M$, a number
  $\varepsilon>0$, and a $C^\infty$ map $\gamma:(-2,2)\times\mathcal{U}\to M$,
  $\mathcal{U}=\{(q,w)\in TM;\ q\in V,\ w\in T_qM,\ |w|<\varepsilon\}$, such that
  $t\mapsto\gamma(t,q,w)$, $t\in(-2,2)$, is the unique geodesic of $M$ passing
  through $q$ with velocity $w$ at $t=0$, for every $q\in V$ and $w\in T_qM$ with
  $|w|<\varepsilon$.
\end{proposition}
\begin{proof}
  \leanok
  \sketch
  The geodesic $\gamma(t,q,v)$ of \cref{prop:dc-ch3-2-5} is defined for
  $|t|<\delta$ and $|v|<\varepsilon_1$. By the lemma of homogeneity (2.6),
  $\gamma(t,q,\frac{\delta v}{2})$ is defined for $|t|<2$. Taking
  $\varepsilon<\frac{\delta\varepsilon_1}{2}$, the geodesic $\gamma(t,q,w)$ is
  defined for $|t|<2$ and $|w|<\varepsilon$.
\end{proof}

\begin{remark}[large velocity]
  \label{rem:dc-ch3-2-8}
  \lean{Riemannian.Geodesic.maximalGeodesicInterval_fiberScale}
  \uses{lem:dc-ch3-2-6}
  \leanok
  \dcref{ch3:2.8}
  By an analogous argument, one can make the velocity of a geodesic uniformly large
  in a neighborhood of $p$: the exact rescaling
  $I_{\max}(p,\,a\,v)=\{t:\ a\,t\in I_{\max}(p,v)\}$ of the maximal interval
  trades interval length against velocity in both directions.
\end{remark}

\subsection{Exponential map}
\cref{prop:dc-ch3-2-7} lets us define the \emph{exponential map}. Let $p\in M$ and let
$\mathcal{U}\subset TM$ be as in \cref{prop:dc-ch3-2-7}. The map $\exp:\mathcal{U}\to M$,

$$
\exp(q,v)=\gamma(1,q,v)=\gamma\Big(|v|,q,\frac{v}{|v|}\Big),\quad(q,v)\in\mathcal{U},
$$

is the \emph{exponential map} on $\mathcal{U}$; it is differentiable. In applications
we use the restriction to an open set of $T_qM$:

$$
\exp_q:B_\varepsilon(0)\subset T_qM\to M,\qquad\exp_q(v)=\exp(q,v),
$$

where $B_\varepsilon(0)$ is the open ball of center $0$ and radius $\varepsilon$
in $T_qM$. Then $\exp_q$ is differentiable and $\exp_q(0)=q$. Geometrically,
$\exp_q(v)$ is obtained by going a length $|v|$ from $q$ along the geodesic with
velocity $\frac{v}{|v|}$.

\begin{definition}[exponential map]
  \label{def:dc-ch3-expmap}
  \lean{Riemannian.Exponential.expMapIntrinsic, Riemannian.Exponential.expDomainIntrinsic, Riemannian.Exponential.zero_mem_expDomainIntrinsic, Riemannian.Exponential.expMapIntrinsic_zero}
  \uses{lem:dc-ch3-2-5-exist, lem:dc-ch3-2-5-unique}
  \leanok
  \dcref{ch3:after-2.8}
  The \emph{exponential map} at $p\in M$ is $\exp_p(v)=\gamma(1,p,v)$, the value at
  time $1$ of the (maximal) geodesic with initial data $(p,v)$; its natural
  domain is the set of $v\in T_pM$ whose maximal geodesic is defined at time
  $1$, which always contains $0$, and $\exp_p(0)=p$.

\end{definition}

\begin{lemma}[the exponential map traces geodesic rays]
  \label{lem:dc-ch3-expmap-ray}
  \lean{Riemannian.Exponential.expMapIntrinsic_smul_eq_intrinsicMaximalGeodesic_of_mem}
  \uses{def:dc-ch3-expmap, lem:dc-ch3-2-6}
  \leanok
  \dcref{ch3:after-2.8}
  For $v\in T_pM$ and $t$ in the maximal interval of the geodesic with initial
  data $(p,v)$,
  $$\exp_p(t\,v)=\gamma(t,p,v).$$
  In particular $t\mapsto\exp_p(t\,v)$ \emph{is} the geodesic through $p$ with
  initial velocity $v$; this is the computational heart of
  $d(\exp_p)_0=\mathrm{id}$ in \cref{prop:dc-ch3-2-9}.
\end{lemma}
\begin{proof}
  \sketch
  \leanok
  For $t=0$ both sides are $p$. For $t\neq0$ this is the homogeneity
  $\gamma(1,p,t\,v)=\gamma(t\cdot 1,p,v)$ of \cref{lem:dc-ch3-2-6} evaluated at
  time $1$.
\end{proof}

\begin{lemma}[$d(\exp_p)_0=\mathrm{id}$ along rays]
  \label{lem:dc-ch3-2-9-dexp}
  \lean{Riemannian.Exponential.hasDerivAt_extChartAt_expMap_smul}
  \uses{lem:dc-ch3-expmap-ray, lem:dc-ch3-2-5-velocity}
  \leanok
  \dcref{ch3:2.9}
  For every $v\in T_pM$, read in the chart at $p$, the curve
  $t\mapsto\exp_p(t\,v)$ is differentiable at $t=0$ with derivative $v$:
  $$
  d(\exp_p)_0(v)=\frac{d}{dt}\big(\exp_p(tv)\big)\Big|_{t=0}
    =\frac{d}{dt}\big(\gamma(t,p,v)\big)\Big|_{t=0}=v.
  $$
  Together with smooth dependence on initial conditions and the inverse function
  theorem, this derivative identity gives the local regularity and local
  invertibility of $\exp_p$ near $0$.
\end{lemma}
\begin{proof}
  \sketch
  \leanok
  By \cref{lem:dc-ch3-expmap-ray}, $\exp_p(t\,v)=\gamma(t,p,v)$ for $t$ in the
  maximal interval, hence for all $t$ near $0$; by
  \cref{lem:dc-ch3-2-5-velocity} the geodesic $\gamma(\cdot,p,v)$ has velocity
  $v$ at $0$.
\end{proof}

\begin{lemma}[$\exp_p$ is strictly differentiable at $0$ with derivative the identity]
  \label{lem:dc-ch3-2-9-strict}
  \lean{Riemannian.Exponential.exists_hasStrictFDerivAt_extChartAt_expMap, Riemannian.Exponential.map_expMap_nhds, Riemannian.Exponential.exists_injOn_expMap, Riemannian.FlowDependence.hasStrictFDerivAt_of_picardResidual}
  \uses{def:dc-ch3-expmap, lem:dc-ch3-2-6, lem:dc-ch3-2-7-vuniform, lem:dc-ch3-2-9-dexp}
  \leanok
  \dcref{ch3:2.9}
  There is $\rho > 0$ such that the ball $B_\rho(0) \subset T_pM$ lies in the
  domain of $\exp_p$, and, read in the chart at $p$, the map $\exp_p$ is
  \emph{strictly} Fréchet differentiable at $0$ with derivative the identity.
  Consequently (Banach inverse function theorem) $\exp_p$ maps the
  neighbourhood filter of $0$ exactly onto the neighbourhood filter of $p$ —
  the image of any neighbourhood of $0$ is a neighbourhood of $p$ — and
  $\exp_p$ is injective on a ball around $0$. This is
  \cref{prop:dc-ch3-2-9} minus the $C^\infty$ regularity of $\exp_p$ away from
  the strict-derivative point.
\end{lemma}
\begin{proof}
  \sketch
  \leanok
  $C^1$ dependence on initial conditions at an equilibrium, on the Banach
  space $C([0,T], E \times E)$: the Picard integral equation
  $G(x, \alpha) = \alpha - x - \int_0^t F(\alpha)$ is strictly differentiable
  at the zero-section equilibrium — the superposition operator
  $\alpha \mapsto F \circ \alpha$ has a strict derivative at a constant curve
  as soon as $DF$ is continuous at the point — and its curve-partial
  derivative $1 - J \circ M_{DF}$ is invertible by a Neumann series once
  $T\,\|DF\| < 1$; the Banach implicit function theorem then makes the
  solution family strictly differentiable in its initial condition. The
  linearization of the geodesic spray at the zero section is
  $A(u, w) = (w, 0)$, quadratically nilpotent ($A^2 = 0$, the Christoffel term
  being quadratic in the velocity), so the Neumann series terminates and the
  derivative is computed exactly; the homogeneity
  $\gamma(1, p, w) = \gamma(T, p, w/T)$ of \cref{lem:dc-ch3-2-6} trades the
  unit evaluation time for the short Picard time $T$. Evaluating at time $T$
  and projecting to the base gives the strict derivative
  $d(\exp_p)_0 = \mathrm{id}$ of the chart reading, and the inverse function
  theorem for strictly differentiable maps yields the neighbourhood and
  injectivity statements.
\end{proof}

\begin{lemma}[$\exp_p$ is strictly differentiable on a ball]
  \label{lem:dc-ch3-2-9-c1ball}
  \lean{Riemannian.Exponential.exists_hasStrictFDerivAt_extChartAt_expMap_ball}
  \uses{def:dc-ch3-expmap, lem:dc-ch3-2-6, lem:dc-ch3-2-7-c1flow}
  \leanok
  \dcref{ch3:2.9}
  There is $\rho > 0$ such that the ball $B_\rho(0) \subset T_pM$ lies in the
  domain of $\exp_p$, its image under $\exp_p$ stays in the chart at $p$, and
  at \emph{every} $v_0$ with $\|v_0\| < \rho$ — zero or not — the chart
  reading $w \mapsto \varphi_p(\exp_p(w))$ has a strict Fr\'echet derivative
  at $v_0$. In particular $\exp_p$ is differentiable at every point of a
  neighbourhood of the origin, which is the regularity clause of
  \cref{prop:dc-ch3-2-9} at the $C^1$ level; at $v_0 \ne 0$ the derivative is
  the solution of the variational equation along the geodesic
  $t \mapsto \exp_p(t\,v_0)$, not the identity.
\end{lemma}
\begin{proof}
  \sketch
  \leanok
  Let $Z$, $r$, $\varepsilon$, $T$ be as in \cref{lem:dc-ch3-2-7-c1flow} and
  put $\rho = rT$. For $\|w\| < \rho$, the homogeneity of geodesics
  (\cref{lem:dc-ch3-2-6}, realized at the level of the flow witness by fibre
  scaling) trades the fixed evaluation time $1$ for the short Picard time
  $T$: $\exp_p(w) = \gamma(1, p, w) = \gamma(T, p, w/T)$, and the chart
  reading of the right-hand side is computed by the flow,
  $\varphi_p(\exp_p(w)) = \big(Z((\varphi_p(p), w/T))(T)\big)_1$. In
  particular $B_\rho(0)$ lies in the exponential domain and its image stays
  in the chart. The map $w \mapsto (\varphi_p(p), w/T)$ is continuous affine
  with linear part $w \mapsto (0, w/T)$, and it sends $B_\rho(0)$ into the
  open flow ball; composing its strict derivative with the strict derivative
  in the initial condition of the flow family (\cref{lem:dc-ch3-2-7-c1flow})
  evaluated at the fixed time $T$, then with the (linear, bounded) projection
  to the base coordinate, yields the strict derivative of the chart reading
  of $\exp_p$ at every $v_0 \in B_\rho(0)$.
\end{proof}

\begin{lemma}[$\exp_p$ is $C^1$ near the origin]
  \label{lem:dc-ch3-2-9-c1}
  \lean{Riemannian.Exponential.exists_contDiffOn_extChartAt_expMap_ball, Riemannian.FlowDependence.contDiffOn_one_of_forall_hasStrictFDerivAt, Riemannian.FlowDependence.continuousOn_of_forall_hasStrictFDerivAt}
  \uses{lem:dc-ch3-2-9-c1ball}
  \leanok
  \dcref{ch3:2.9}
  There is $\rho > 0$ such that the ball $B_\rho(0) \subset T_pM$ lies in the
  domain of $\exp_p$, its image under $\exp_p$ stays in the chart at $p$, and
  the chart reading $w \mapsto \varphi_p(\exp_p(w))$ is $C^1$ on $B_\rho(0)$.
  This upgrades the pointwise strict differentiability of
  \cref{lem:dc-ch3-2-9-c1ball} to genuine $C^1$ regularity, the differentiable
  half of the local-diffeomorphism claim of \cref{prop:dc-ch3-2-9} at the
  $C^1$ level.
\end{lemma}
\begin{proof}
  \sketch
  \leanok
  A map that is \emph{strictly} differentiable at every point of an open set
  automatically has a continuous derivative there: if $x, x'$ are close, the
  strict estimates
  $\|f(y) - f(z) - df_x(y - z)\| \le \varepsilon \|y - z\|$ at $x$ and the
  corresponding estimate at $x'$ both hold for all pairs $y, z$ in a common
  small ball around $x'$; applying them to the pairs $(x' + w,\, x')$ for all
  small $w$ and subtracting gives
  $\|(df_{x'} - df_x)(w)\| \le 2\varepsilon \|w\|$, whence
  $\|df_{x'} - df_x\| \le 2\varepsilon$ by homogeneity. Thus the derivative
  map is continuous on the open set, and on open sets continuity of the
  derivative together with differentiability is exactly the $C^1$
  characterization. Applying this to the chart reading of $\exp_p$ on
  $B_\rho(0)$, which is strictly differentiable at every point by
  \cref{lem:dc-ch3-2-9-c1ball}, gives the claim.
\end{proof}

\begin{lemma}[the derivative of $\exp_p$ is invertible near the origin]
  \label{lem:dc-ch3-2-9-invertible}
  \lean{Riemannian.Exponential.exists_hasStrictFDerivAt_equiv_extChartAt_expMap_ball}
  \uses{lem:dc-ch3-2-9-c1ball, lem:dc-ch3-2-9-strict, lem:dc-ch3-2-9-c1}
  \leanok
  \dcref{ch3:2.9}
  There is $\rho > 0$ such that the ball $B_\rho(0) \subset T_pM$ lies in the
  domain of $\exp_p$, its image under $\exp_p$ stays in the chart at $p$, and
  at every $v_0$ with $\|v_0\| < \rho$ the strict Fr\'echet derivative of the
  chart reading $w \mapsto \varphi_p(\exp_p(w))$ is a continuous linear
  isomorphism of $T_pM$.
\end{lemma}
\begin{proof}
  \sketch
  \leanok
  By \cref{lem:dc-ch3-2-9-c1ball} the chart reading $f$ of $\exp_p$ is strictly
  differentiable at every point of a ball $B_{\rho_0}(0)$, so its derivative
  map $x \mapsto df_x$ is continuous there (\cref{lem:dc-ch3-2-9-c1}), and
  $df_0 = \mathrm{id}$ (\cref{lem:dc-ch3-2-9-strict}). Choose
  $\rho \le \rho_0$ such that $\|df_{v_0} - \mathrm{id}\| < 1$ for
  $\|v_0\| < \rho$. Writing $df_{v_0} = \mathrm{id} - t$ with
  $t = \mathrm{id} - df_{v_0}$, $\|t\| < 1$, the Neumann series
  $\sum_{n \ge 0} t^n$ converges in the Banach algebra of bounded endomorphisms
  of $T_pM$ and inverts $\mathrm{id} - t$; hence $df_{v_0}$ is a continuous
  linear isomorphism.
\end{proof}

\begin{lemma}[$\exp_p$ is a $C^1$ diffeomorphism near the origin]
  \label{lem:dc-ch3-2-9-c1diffeo}
  \lean{Riemannian.Exponential.exists_c1_local_diffeomorphism_expMap}
  \uses{def:dc-ch3-expmap, lem:dc-ch3-2-9-c1, lem:dc-ch3-2-9-invertible, lem:dc-ch3-2-9-strict}
  \leanok
  \dcref{ch3:2.9}
  There is $\varepsilon > 0$ such that the ball
  $B_\varepsilon(0) \subset T_pM$ lies in the domain of $\exp_p$ and
  $\exp_p : B_\varepsilon(0) \to M$ is a $C^1$ diffeomorphism onto an open
  subset of $M$: it is injective on $B_\varepsilon(0)$, its image is open in
  $M$, its chart reading is $C^1$ on $B_\varepsilon(0)$, and it admits a left
  inverse whose chart reading is $C^1$ on the image. This is
  \cref{prop:dc-ch3-2-9} with ``diffeomorphism'' read at the $C^1$ regularity
  level.
\end{lemma}
\begin{proof}
  \sketch
  \leanok
  Let $f = \varphi_p \circ \exp_p$ be the chart reading of $\exp_p$, which is
  $C^1$ on a ball around the origin by \cref{lem:dc-ch3-2-9-c1}, with
  $df_0 = \mathrm{id}$ (\cref{lem:dc-ch3-2-9-strict}). By the inverse function
  theorem for strictly differentiable maps, $f$ admits a local inverse
  $f^{-1}$ near $0$; since $f$ is $C^1$ at $0$ with invertible derivative,
  $f^{-1}$ is $C^1$ on a neighbourhood $V$ of $f(0)$. Moreover $\exp_p$ is
  injective on a small ball (\cref{lem:dc-ch3-2-9-strict}), and at every point
  of the ball of \cref{lem:dc-ch3-2-9-invertible} the strict derivative of $f$
  is invertible, so the inverse function theorem applied at each point makes
  $f$ an open map there: $f(B_\varepsilon(0))$ is open, and the image
  $\exp_p(B_\varepsilon(0))$ — the intersection of the chart domain with the
  $\varphi_p$-preimage of $f(B_\varepsilon(0))$ — is open in $M$. Shrinking
  $\varepsilon$ so that $f(B_\varepsilon(0)) \subset V$ and so that the
  left-inverse identity $f^{-1}(f(w)) = w$ holds on $B_\varepsilon(0)$ gives
  all the claims.
\end{proof}

\begin{lemma}[$\exp_p$ is $C^2$ on a ball]
  \label{lem:dc-ch3-2-9-c2ball}
  \lean{Riemannian.Exponential.exists_contDiffOn_two_extChartAt_expMap_ball}
  \uses{def:dc-ch3-expmap, lem:dc-ch3-2-6, lem:dc-ch3-2-7-c2flow, lem:dc-ch3-2-9-c1}
  \leanok
  \dcref{ch3:2.9}
  There is $\rho > 0$ such that the ball $B_\rho(0) \subset T_pM$ lies in the
  domain of $\exp_p$, its image under $\exp_p$ stays in the chart at $p$, and
  the chart reading $w \mapsto \varphi_p(\exp_p(w))$ is $C^2$ on $B_\rho(0)$.
  This upgrades the $C^1$ regularity of \cref{lem:dc-ch3-2-9-c1} one order:
  the derivative of the chart reading at $w$ is the explicit
  continuous-linear image $\Phi(\tau(\iota(w)))$ of the variational flow of
  \cref{lem:dc-ch3-2-7-c2flow}, and is itself $C^1$ in $w$. It is the
  regularity input for the Gauss lemma (\cref{lem:dc-ch3-3-5}): the mixed
  second derivatives of the chart reading of $\exp_p$ exist and are symmetric
  on the ball.
\end{lemma}
\begin{proof}
  \sketch
  \leanok
  Let $Z$, $r$, $\varepsilon$, $T$, $\sigma$, $\tau$ be as in
  \cref{lem:dc-ch3-2-7-c2flow} and put $\rho = rT$. As in
  \cref{lem:dc-ch3-2-9-c1ball}, the homogeneity of geodesics realized on the
  flow witness by fibre scaling (\cref{lem:dc-ch3-2-6}) identifies, for
  $\|w\| < \rho$,
  $\varphi_p(\exp_p(w)) = \big(\sigma(\iota(w))(T)\big)_1$ where
  $\iota(w) = (\varphi_p(p), w/T)$ is affine into the open flow ball. By the
  chain rule and the computation of $D\sigma$ by the variational flow
  (\cref{lem:dc-ch3-2-7-c2flow}), the strict derivative of the chart reading
  at $w$ is
  $$
  v \;\longmapsto\; \Big(\tau(\iota(w))(T)\,\big(d\iota(v)\big)\Big)_1
  \;=\; \Phi\big(\tau(\iota(w))\big)\,v,
  $$
  where $\Phi$ — evaluate the operator curve at time $T$, compose with the
  fixed linear map $d\iota$, project to the base component — is a fixed
  bounded linear map from $C([0,T], L(\mathbb{R}^{2n}))$ to
  $L(\mathbb{R}^n)$. Since $\tau$ is strictly differentiable at every point
  of the flow ball (\cref{lem:dc-ch3-2-7-c2flow}) and $\Phi$, $\iota$ are
  (affine) continuous linear, the derivative map
  $w \mapsto \Phi(\tau(\iota(w)))$ is strictly differentiable, hence $C^1$,
  on $B_\rho(0)$ (strict differentiability at every point of an open set
  self-improves to $C^1$, as in \cref{lem:dc-ch3-2-9-c1}). A differentiable
  map whose derivative map is $C^1$ on an open set is $C^2$ there.
\end{proof}

\begin{lemma}[$\exp_p$ is a $C^2$ diffeomorphism near the origin]
  \label{lem:dc-ch3-2-9-c2diffeo}
  \lean{Riemannian.Exponential.exists_c2_local_diffeomorphism_expMap}
  \uses{lem:dc-ch3-2-9-c2ball, lem:dc-ch3-2-9-invertible, lem:dc-ch3-2-9-c1diffeo}
  \leanok
  \dcref{ch3:2.9}
  There is $\varepsilon > 0$ such that $\exp_p$ is injective on
  $B_\varepsilon(0) \subset T_pM$ with open image in $M$, the chart reading
  $f = \varphi_p \circ \exp_p$ is $C^2$ on $B_\varepsilon(0)$ with open image
  in the chart, and the local inverse $f^{-1}$ is $C^2$ on the open chart
  image $f(B_\varepsilon(0))$. This upgrades the inverse-regularity of
  \cref{lem:dc-ch3-2-9-c1diffeo} one order; it is the regularity of
  $\exp_p^{-1}$ consumed by the second-derivative computation of
  \cref{lem:dc-ch3-4-1}.
\end{lemma}
\begin{proof}
  \sketch
  \leanok
  \uses{lem:dc-ch3-2-9-c2ball, lem:dc-ch3-2-9-invertible, lem:dc-ch3-2-9-c1diffeo}
  Take $\varepsilon$ below the radii of \cref{lem:dc-ch3-2-9-c2ball} ($f$ is
  $C^2$ on the ball), \cref{lem:dc-ch3-2-9-invertible} ($df$ is invertible at
  every point of the ball), and of the injectivity ball of
  \cref{lem:dc-ch3-2-9-c1diffeo}. Injectivity of $f$ and openness of both
  images follow as in \cref{lem:dc-ch3-2-9-c1diffeo} (an everywhere-invertible
  derivative makes $f$ an open map on the ball). Define $f^{-1}$ globally on
  the image of the injective $f$. At each image point $f(v_0)$ the inverse
  function theorem for $C^2$ maps — applicable since $f$ is $C^2$ at $v_0$
  with invertible derivative — provides a local inverse which is $C^2$ at
  $f(v_0)$; it agrees with $f^{-1}$ near $f(v_0)$, because both are left
  inverses of $f$ on a neighborhood of $v_0$ and $f$ maps neighborhoods of
  $v_0$ onto neighborhoods of $f(v_0)$. Hence $f^{-1}$ is $C^2$ at every
  point of the open image.
\end{proof}

\begin{proposition}[exp_q is a local diffeomorphism]
  \label{prop:dc-ch3-2-9}
  \lean{Riemannian.Exponential.exists_infty_local_diffeomorphism_expMap}
  \uses{def:dc-ch3-expmap, lem:dc-ch3-expmap-ray, lem:dc-ch3-2-9-dexp, lem:dc-ch3-2-9-strict, lem:dc-ch3-2-9-c1ball, lem:dc-ch3-2-9-c1, lem:dc-ch3-2-9-invertible, lem:dc-ch3-2-9-c1diffeo, lem:dc-ch3-2-9-c2ball, lem:dc-ch3-2-9-c2diffeo, lem:dc-ch3-2-9-cinfty}
  \leanok
  \dcref{ch3:2.9}
  \notready
  Given $q\in M$, there exists $\varepsilon>0$ such that
  $\exp_q:B_\varepsilon(0)\subset T_qM\to M$ is a diffeomorphism of
  $B_\varepsilon(0)$ onto an open subset of $M$.
\end{proposition}
\begin{proof}
  \sketch
  Compute $d(\exp_q)_0$:
  
  $$
  d(\exp_q)_0(v)=\frac{d}{dt}\big(\exp_q(tv)\big)\Big|_{t=0}=\frac{d}{dt}\big(\gamma(1,q,tv)\big)\Big|_{t=0}=\frac{d}{dt}\big(\gamma(t,q,v)\big)\Big|_{t=0}=v.
  $$
  
  Hence $d(\exp_q)_0=\mathrm{Id}$ on $T_qM$, and by the inverse function theorem
  $\exp_q$ is a local diffeomorphism on a neighborhood of $0$.
\end{proof}

\begin{example}[Euclidean space]
  \label{ex:dc-ch3-2-10}
  \lean{Riemannian.euclidean_chartBasisVecFiber, Riemannian.euclidean_chartGramMatrix, Riemannian.euclidean_chartChristoffel, Riemannian.euclidean_chartChristoffelContraction, Riemannian.euclideanGeodesic, Riemannian.isGeodesic_euclideanGeodesic, Riemannian.expMapIntrinsic_euclidean}
  \leanok
  \dcref{ch3:2.10}
  \notready
  Let $M=\mathbb{R}^n$. The covariant derivative is the usual derivative, so the
  geodesics are straight lines parametrized proportionally to arc length; the
  exponential is the identity (under the usual identification $T_p\mathbb{R}^n\cong\mathbb{R}^n$).
\end{example}

\begin{example}[the sphere]
  \label{ex:dc-ch3-2-11}
  \uses{prop:dc-ch3-2-5}
  \notready
  \dcref{ch3:2.11}
  \notready
  Let $M=S^n\subset\mathbb{R}^{n+1}$ be the unit sphere. Great circles
  parametrized by arc length are geodesics. All
  geodesics of $S^n$ are great circles parametrized proportionally to arc length:
  given $p\in S^n$ and a unit $v\in T_pS^n$, the intersection of $S^n$ with the
  plane through the origin, $p$, and $v$ is a great circle parametrizable as the
  geodesic through $p$ with velocity $v$; by uniqueness of \cref{prop:dc-ch3-2-5} the claim
  follows. Here $\exp_p$ is defined on all of $T_p S^n$: it maps $B_\pi(0)$
  injectively onto $S^n-\{q\}$ ($q$ the antipode of $p$), collapses
  $\partial B_\pi(0)$ to $q$, maps the annulus $B_{2\pi}(0)-\overline{B_\pi(0)}$
  injectively onto $S^n-\{p,q\}$, and collapses $\partial B_{2\pi}(0)$ to $p$.
  For the manifold $S^n-\{q\}$, however, $\exp_p$ is defined only on
  $B_\pi(0)\subset T_p(S^n-\{q\})$.
\end{example}

\section{Minimizing properties of geodesics}

\begin{definition}[piecewise differentiable curve]
  \label{def:dc-ch3-3-1}
  \lean{Riemannian.Geodesic.IsPiecewiseDifferentiableCurve, Riemannian.Geodesic.IsPiecewiseDifferentiableCurveOfOrder, Riemannian.Geodesic.IsPiecewiseSmoothCurve, Riemannian.Geodesic.SubdividedCurveOfOrder, Riemannian.Geodesic.SubdividedCurveOfOrder.vertex}
  \dcref{ch3:3.1}
  A \emph{piecewise differentiable curve} is a continuous map $c:[a,b]\to M$ for which
  there is a partition $a=t_0<t_1<\dots<t_k=b$ such that each restriction
  $c|_{[t_i,t_{i+1}]}$ is differentiable; $c$ \emph{joins} $c(a)$ and $c(b)$. $c(t_i)$ is
  a \emph{vertex} of $c$, and the \emph{vertex angle at $c(t_i)$} is the angle between
  $\lim_{t\to t_i^+}c'(t)$ and $\lim_{t\to t_i^-}c'(t)$. Parallel transport extends
  to such curves by transporting successively over each differentiable segment.
\end{definition}

\begin{definition}[minimizing geodesic]
  \label{def:dc-ch3-3-2}
  \lean{Riemannian.Geodesic.IsMinimizingGeodesicSegment}
  \uses{def:dc-ch3-3-1}
  \leanok
  \dcref{ch3:3.2}
  A segment of a geodesic $\gamma:[a,b]\to M$ is \emph{minimizing} if
  $\ell(\gamma)\le\ell(c)$ for every piecewise differentiable curve $c$ joining
  $\gamma(a)$ to $\gamma(b)$.
\end{definition}

\begin{definition}[parametrized surface]
  \label{def:dc-ch3-3-3}
  \lean{Riemannian.Geodesic.IntrinsicParametrizedSurface,
    Riemannian.Geodesic.IntrinsicParametrizedSurface.velocityU,
    Riemannian.Geodesic.IntrinsicParametrizedSurface.velocityV,
    Riemannian.Geodesic.IsSmoothParametrizedSurfaceWithBoundary, Riemannian.Geodesic.IsNonPiAngle, Riemannian.Geodesic.IsNonPiVertexAngle, Riemannian.Geodesic.isNonPiVertexAngle_iff_vertexAngle_ne_pi, Riemannian.Geodesic.vertexAngle_eq_pi_of_reversed, Riemannian.Geodesic.isNonPiAngle_iff_angle_ne_pi, Riemannian.Geodesic.PiecewiseC1BoundaryParametrization.interior_vertex_angle_ne_pi, Riemannian.Geodesic.PiecewiseC1BoundaryParametrization.closing_vertex_angle_ne_pi, Riemannian.Geodesic.IsParametrizedSurfaceWithBoundary.toExtended, Riemannian.Geodesic.IsParametrizedSurfaceWithBoundary.surfaceDomain, Riemannian.Geodesic.IsParametrizedSurfaceWithBoundary.regularOn, Riemannian.Geodesic.IsParametrizedSurfaceWithBoundary.isParametrizedSurface, Riemannian.Geodesic.PiecewiseC1BoundaryParametrization.interior_vertex_nonPi, Riemannian.Geodesic.PiecewiseC1BoundaryParametrization.closing_vertex_nonPi}
  \leanok
  \dcref{ch3:3.3}
  Let $A$ be a connected set in $\mathbb{R}^2$, $U\subset A\subset\overline{U}$, $U$
  open, such that $\partial A$ is a piecewise differentiable curve with vertex
  angles different from $\pi$. A \emph{parametrized surface} in $M$ is a differentiable
  map $s:A\subset\mathbb{R}^2\to M$ (differentiable on $A$ meaning extendable to an
  open $U\supset A$; the vertex-angle condition ensures $ds$ is independent of the
  extension). A vector field $V$ along $s$ assigns to each $q\in A$ a vector
  $V(q)\in T_{s(q)}M$; it is differentiable if $q\mapsto V(q)f$ is differentiable
  for every differentiable function $f$ on $M$. For cartesian $(u,v)$ on
  $\mathbb{R}^2$, $\frac{\partial s}{\partial u}(u_0,v_0)$ is the velocity of
  $u\mapsto s(u,v_0)$ at $u_0$, and $\frac{\partial s}{\partial v}$ is analogous.
  If $V$ is a vector field along $s$, $\frac{DV}{\partial u}(u_0,v_0)$ is the
  covariant derivative along $u\mapsto s(u,v_0)$ of the restriction of $V$, and
  $\frac{DV}{\partial v}$ is defined analogously.
\end{definition}

\begin{lemma}[symmetry]
  \label{lem:dc-ch3-3-4}
  \lean{Riemannian.Geodesic.IntrinsicParametrizedSurface.mixedCovariantDerivatives_commute,
    Riemannian.Geodesic.ParametrizedSurface.covariantDerivativeUOfVelocityV_eq_covariantDerivativeVOfVelocityU,
    Riemannian.Geodesic.covariant_sndFDeriv_symm, Riemannian.Geodesic.covariant_sndFDeriv_symm_of_eventually, Riemannian.Geodesic.covariant_sndFDeriv_symm_of_surface}
  \leanok
  \dcref{ch3:3.4}
  If $M$ has a symmetric connection and $s:A\to M$ is a parametrized surface, then
  $$
  \frac{D}{\partial v}\frac{\partial s}{\partial u}=\frac{D}{\partial u}\frac{\partial s}{\partial v}.
  $$
\end{lemma}
\begin{proof}
  \sketch
  \leanok
  In a coordinate system $\mathbf{x}$ with
  $\mathbf{x}^{-1}\circ s(u,v)=(x^1(u,v),\dots,x^n(u,v))$,
  
  $$
  \frac{D}{\partial v}\Big(\frac{\partial s}{\partial u}\Big)=\sum_i\frac{\partial^2 x^i}{\partial v\,\partial u}\frac{\partial}{\partial x^i}+\sum_{i,j}\frac{\partial x^i}{\partial u}\frac{\partial x^j}{\partial v}\nabla_{\partial/\partial x^j}\frac{\partial}{\partial x^i}.
  $$
  
  By symmetry of the connection,
  $\nabla_{\partial/\partial x^j}\frac{\partial}{\partial x^i}=\nabla_{\partial/\partial x^i}\frac{\partial}{\partial x^j}$,
  so computing $\frac{D}{\partial u}\frac{\partial s}{\partial v}$ gives the same
  expression.
\end{proof}

\begin{lemma}[rays of the exponential chart reading are coordinate geodesics]
  \label{lem:dc-ch3-3-5-rayode}
  \lean{Riemannian.Exponential.exists_expMap_ray_ode_ball}
  \uses{def:dc-ch3-expmap, lem:dc-ch3-2-6, lem:dc-ch3-2-7-vuniform, lem:dc-ch3-2-7-c2flow, lem:dc-ch3-2-9-c2ball}
  \leanok
  \dcref{ch3:3.5}
  There are $\rho > 0$ and $b > 1$ such that, writing
  $f = \varphi_p \circ \exp_p$ for the chart reading of the exponential map:
  (i) for $\|u\| < \rho$ and $|a| < b$ the vector $a\,u$ lies in the domain of
  $\exp_p$ and $\exp_p(a\,u)$ stays in the chart at $p$;
  (ii) $f$ is $C^2$ on $B_\rho(0)$;
  (iii) $(df)_0 = \mathrm{id}$;
  (iv) for $\|u\| < \rho$, $|t| < b$, $\|t\,u\| < \rho$, the ray velocity
  $V(t) = (df)_{tu}(u)$ satisfies the chart-coordinate geodesic equation
  $$V'(t) = -\,\Gamma\big(V(t), V(t)\big)\big(f(tu)\big),$$
  so each ray $t \mapsto f(t\,u)$ is a coordinate geodesic.
\end{lemma}
\begin{proof}
  \sketch
  \leanok
  Let $Z$, $r$, $\varepsilon$, $T$ be the uniform local flow of the coordinate
  spray (\cref{lem:dc-ch3-2-7-c2flow}); put
  $\rho = \min(\rho_2, rT)$, with $\rho_2$ the $C^2$ radius of
  \cref{lem:dc-ch3-2-9-c2ball}, and $b = \varepsilon/T > 1$. For
  $\|u\| \le rT$ and $|a| < b$, the fibre- and time-rescaled trajectory
  $$\zeta(s) = \big(Z_1(aTs),\; aT\,Z_2(aTs)\big),$$
  computed along the flow line of the initial condition
  $(\varphi_p(p), u/T)$, again solves the spray ODE — the first component
  because $\dot Z_1 = Z_2$, the second because the Christoffel contraction is
  quadratic in the velocity — now with initial data $(\varphi_p(p), a\,u)$,
  on the interval $\{s : |aTs| < \varepsilon\} \supseteq [0,1]$. By the
  witness identification (\cref{lem:dc-ch3-2-7-vuniform}) the canonical
  geodesic through $(p, a u)$ is defined at time $1$ with
  $\varphi_p(\exp_p(a u)) = Z_1(aT)$; this realizes the homogeneity
  $\exp_p(a u) = \gamma(a, p, u)$ of \cref{lem:dc-ch3-2-6} at the level of
  the flow, proving (i). Clause (ii) is \cref{lem:dc-ch3-2-9-c2ball}.
  Differentiating the identification in $a$ and comparing with the chain
  rule at points of the $C^2$ ball gives $(df)_{tu}(u) = T\,Z_2(tT)$; at
  $t = 0$ this is $u$, which proves (iii) after extending by homogeneity of
  both sides, and differentiating once more in $t$ through the flow ODE
  $\dot Z_2 = -\Gamma(Z_2, Z_2)(Z_1)$ gives (iv).
\end{proof}

\begin{lemma}[Gauss surface computation]
  \label{lem:dc-ch3-3-5-surface}
  \lean{Riemannian.Exponential.gauss_surface_computation}
  \uses{def:dc-ch3-3-3, lem:dc-ch3-3-4, lem:dc-ch2-3-2-coord}
  \leanok
  \dcref{ch3:3.5}
  Let $\rho > 0$, $b > 1$, and let $f$ be a map of $B_\rho(0)$ into the chart
  target at $p$ which is $C^2$ on $B_\rho(0)$ with $f(0) = \varphi_p(p)$ and
  $(df)_0 = \mathrm{id}$, and whose ray velocities satisfy the
  chart-coordinate geodesic equation as in \cref{lem:dc-ch3-3-5-rayode}(iv).
  Then for $\|v\| < \rho$ and every $w$,
  $$\big\langle (df)_v(v),\, (df)_v(w) \big\rangle_{f(v)}
    = \langle v, w\rangle_{\varphi_p(p)},$$
  the inner products being the chart Gram inner products at the indicated
  base points: $f$ is a radial isometry.
\end{lemma}
\begin{proof}
  \sketch
  \leanok
  Consider the parametrized surface (\cref{def:dc-ch3-3-3})
  $c(t,s) = f\big(t\,(v + s w)\big)$ on the open parameter set
  $\Omega = \{(t,s) : \|t(v+sw)\| < \rho,\ \|v+sw\| < \rho,\ |t| < b\}$,
  which contains $[0,1] \times \{0\}$; its partial velocities are
  $\partial_t c = (df)_{t(v+sw)}(v+sw)$ and $\partial_s c = (df)_{t(v+sw)}(tw)$.

  \emph{Constant speed.} For fixed $s$ the curve $t \mapsto c(t,s)$ is a
  coordinate geodesic, so by metric compatibility
  (\cref{lem:dc-ch2-3-2-coord})
  $\partial_t \langle \partial_t c, \partial_t c\rangle
  = 2\,\langle \tfrac{D}{\partial t}\partial_t c, \partial_t c\rangle = 0$:
  the squared speed $\psi(t,s)$ is constant in $t$ on the (connected)
  $t$-sections of $\Omega$, and $\psi(0,s) = \langle v+sw, v+sw\rangle$
  because $(df)_0 = \mathrm{id}$.

  \emph{The mixed term.} Let
  $h(t) = \langle \partial_s c, \partial_t c\rangle(t,0)$. By compatibility
  along the $t$-direction and the geodesic equation
  $\tfrac{D}{\partial t}\partial_t c = 0$,
  $h'(t) = \langle \tfrac{D}{\partial t}\partial_s c, \partial_t c\rangle$;
  by the symmetry lemma (\cref{lem:dc-ch3-3-4}),
  $\tfrac{D}{\partial t}\partial_s c = \tfrac{D}{\partial s}\partial_t c$,
  and by compatibility along the $s$-direction
  $$2\,\big\langle \tfrac{D}{\partial s}\partial_t c, \partial_t c\big\rangle
    = \partial_s\, \psi(t, \cdot)\big|_{s=0}
    = \tfrac{d}{ds}\,\langle v+sw, v+sw\rangle\big|_{s=0}
    = 2\,\langle v, w\rangle.$$
  Hence $h'(t) = \langle v, w\rangle$ for every $t \in [0,1]$.

  \emph{Endpoints.} $h(0) = 0$ since
  $\partial_s c(0, \cdot) = (df)_0(0 \cdot w) = 0$, so
  $h(1) = \langle v, w\rangle$. As
  $h(1) = \langle (df)_v(w), (df)_v(v)\rangle_{f(v)}$, symmetry of the
  metric gives the claim.
\end{proof}

\begin{lemma}[Gauss lemma on a ball]
  \label{lem:dc-ch3-3-5-ball}
  \lean{Riemannian.Exponential.exists_gauss_lemma_ball}
  \uses{def:dc-ch3-expmap, lem:dc-ch3-3-5-rayode, lem:dc-ch3-3-5-surface}
  \leanok
  \dcref{ch3:3.5}
  There is $\rho > 0$ such that the ball $B_\rho(0) \subset T_pM$ lies in the
  domain of $\exp_p$, its image under $\exp_p$ stays in the chart at $p$, and
  for every $v$ with $\|v\| < \rho$ and every
  $w \in T_pM \approx T_v(T_pM)$, read in the chart at $p$,
  $$\big\langle (d\exp_p)_v(v),\, (d\exp_p)_v(w) \big\rangle_{\exp_p(v)}
    = \langle v, w\rangle_p.$$
  In particular $\exp_p$ preserves the radial component of the metric: the
  geodesic spheres $\exp_p(\partial B_r(0))$, $r < \rho$, are orthogonal to
  the radial geodesics.
\end{lemma}
\begin{proof}
  \sketch
  \leanok
  Apply the surface computation (\cref{lem:dc-ch3-3-5-surface}) to the chart
  reading $f = \varphi_p \circ \exp_p$, whose ball domain, $C^2$ regularity,
  value $f(0)=\varphi_p(p)$, derivative $(df)_0 = \mathrm{id}$, and
  ray-geodesic property are supplied by \cref{lem:dc-ch3-3-5-rayode}.

\end{proof}

\begin{lemma}[radial lower bound]
  \label{lem:dc-ch3-3-5-radial}
  \lean{Riemannian.Exponential.exists_gauss_radial_lower_bound_ball}
  \uses{lem:dc-ch3-3-5-ball}
  \leanok
  \dcref{ch3:3.5}
  On the Gauss ball of \cref{lem:dc-ch3-3-5-ball}: for every $v$ with
  $\|v\| < \rho$ and every $\xi \in T_pM$, read in the chart at $p$,
  $$\langle v, \xi\rangle_p^2
    \;\le\; \langle v, v\rangle_p \cdot
      \big\langle (d\exp_p)_v(\xi),\, (d\exp_p)_v(\xi)\big\rangle_{\exp_p(v)}.$$
  The exponential map does not shrink the radial component of any vector;
  this is the pointwise inequality behind the minimizing property of radial
  geodesics (\cref{prop:dc-ch3-3-6}).
\end{lemma}
\begin{proof}
  \sketch
  \leanok
  For $v = 0$ both sides vanish (the right-hand side is nonnegative, the
  metric being positive semidefinite). For $v \ne 0$ put
  $\lambda = \langle v,\xi\rangle_p / \langle v,v\rangle_p$ and decompose
  $\xi = \lambda v + \xi_N$ with $\langle v, \xi_N\rangle_p = 0$. Expanding by
  bilinearity and applying the Gauss identity (\cref{lem:dc-ch3-3-5-ball})
  twice,
  $$\big|(d\exp_p)_v(\xi)\big|^2
    = \lambda^2 \langle v,v\rangle_p
      + \big|(d\exp_p)_v(\xi_N)\big|^2
    \;\ge\; \lambda^2 \langle v,v\rangle_p
    = \frac{\langle v,\xi\rangle_p^2}{\langle v,v\rangle_p},$$
  the cross term $2\lambda\,\langle (d\exp_p)_v(v), (d\exp_p)_v(\xi_N)\rangle
  = 2\lambda\,\langle v, \xi_N\rangle_p$ vanishing by orthogonality.
  Multiplying through by $\langle v,v\rangle_p > 0$ gives the claim.

\end{proof}

\begin{lemma}[Gauss]
  \label{lem:dc-ch3-3-5}
  \uses{def:dc-ch3-expmap, lem:dc-ch3-3-5-ball}
  \notready
  \dcref{ch3:3.5}
  \notready
  We identify $T_pM$ with $T_v(T_pM)$ at $v\in T_pM$, writing $T_pM\approx T_v(T_pM)$.
  Let $p\in M$ and $v\in T_pM$ with $\exp_p v$ defined, and $w\in T_pM\approx T_v(T_pM)$.
  Then
  $$
  \langle(d\exp_p)_v(v),\,(d\exp_p)_v(w)\rangle=\langle v,w\rangle.\qquad\text{(2)}
  $$
\end{lemma}
\begin{proof}
  \sketch
  Write $w=w_T+w_N$, $w_T$ parallel to $v$, $w_N$ normal to $v$. Since
  $d\exp_p$ is linear and
  $\langle(d\exp_p)_v(v),(d\exp_p)_v(w_T)\rangle=\langle v,w_T\rangle$, it suffices
  to prove (2) for $w=w_N$ ($w_N\ne 0$). Pick a curve $v(s)$ in $T_pM$ with
  $v(0)=v$, $v'(0)=w_N$, $|v(s)|=\text{const}$. Since $\exp_p v$ is defined,
  choose $\varepsilon>0$ so that $\exp_p(tv(s))$ is defined for $0\le t\le 1$,
  $|s|<\varepsilon$, and consider the parametrized surface
  $f(t,s)=\exp_p tv(s)$ whose curves $t\mapsto f(t,s_0)$ are geodesics.
  Then $\langle\frac{\partial f}{\partial s},\frac{\partial f}{\partial t}\rangle(1,0)=\langle(d\exp_p)_v(w_N),(d\exp_p)_v(v)\rangle$.
  Since $\frac{\partial f}{\partial t}$ is the tangent of a geodesic and by symmetry
  of the connection,
  
  $$
  \frac{\partial}{\partial t}\Big\langle\frac{\partial f}{\partial s},\frac{\partial f}{\partial t}\Big\rangle=\Big\langle\frac{D}{\partial t}\frac{\partial f}{\partial s},\frac{\partial f}{\partial t}\Big\rangle=\Big\langle\frac{D}{\partial s}\frac{\partial f}{\partial t},\frac{\partial f}{\partial t}\Big\rangle=\frac12\frac{\partial}{\partial s}\Big\langle\frac{\partial f}{\partial t},\frac{\partial f}{\partial t}\Big\rangle=0,
  $$
  
  so $\langle\frac{\partial f}{\partial s},\frac{\partial f}{\partial t}\rangle$ is
  independent of $t$. As $\lim_{t\to0}\frac{\partial f}{\partial s}(t,0)=\lim_{t\to0}(d\exp_p)_{tv}\,t w_N=0$,
  the inner product is $0$, proving (2).
\end{proof}

\subsection{Normal neighborhoods}
If $\exp_p$ is a diffeomorphism of a neighborhood $V\ni 0$ in $T_pM$ onto
$U=\exp_p(V)$, then $U$ is a \emph{normal neighborhood} of $p$. If
$\overline{B_\varepsilon(0)}\subset V$, then $\exp_p B_\varepsilon(0)=B_\varepsilon(p)$
is the \emph{normal ball} (or \emph{geodesic ball}) of center $p$ and radius $\varepsilon$.
By the Gauss lemma (3.5), its boundary is a hypersurface orthogonal to the
geodesics from $p$, the \emph{geodesic sphere} $S_\varepsilon(p)$; the geodesics in
$B_\varepsilon(p)$ that begin at $p$ are the \emph{radial geodesics}.

\begin{lemma}[smoothed-radius comparison]
  \label{lem:dc-ch3-3-6-radius}
  \lean{Riemannian.Exponential.gauss_radius_comparison, Riemannian.Exponential.exists_gauss_radius_comparison_ball, Riemannian.hasDerivAt_chartMetricInner_const_base, Riemannian.continuousOn_chartMetricInner_along, Riemannian.chartMetricInner_self_nonneg_of_mem_target}
  \uses{lem:dc-ch3-3-5-radial}
  \leanok
  \dcref{ch3:3.6}
  On the Gauss ball $B_\rho(0)\subset T_pM$ of \cref{lem:dc-ch3-3-5-radial}: for
  every differentiable curve $w:[a,b]\to B_\rho(0)$ with continuous derivative,
  write $c=\exp_p\circ w$ for the composed curve (read in the chart at $p$) and
  $r(t)=\sqrt{\langle w(t),w(t)\rangle_p}$ for the radius of the polar lift.
  Then
  $$r(b)-r(a)\;\le\;\int_a^b\big|\dot c(t)\big|_{c(t)}\,dt,$$
  where $\dot c(t)=(d\exp_p)_{w(t)}\,w'(t)$ by the chain rule and the norm is
  taken in the metric at the moving foot.
\end{lemma}
\begin{proof}
  \sketch
  \leanok
  For $\delta>0$ consider the smoothed radius
  $r_\delta(t)=\sqrt{\langle w(t),w(t)\rangle_p+\delta}$. The map
  $t\mapsto\langle w(t),w(t)\rangle_p$ is differentiable with derivative
  $2\langle w(t),w'(t)\rangle_p$ — the base point of the pairing is fixed, so
  this is the bilinear product rule with no Christoffel correction — hence
  $r_\delta$ is differentiable everywhere (the smoothing keeps the square root
  away from $0$) with
  $$r_\delta'(t)=\frac{\langle w(t),w'(t)\rangle_p}
    {\sqrt{\langle w(t),w(t)\rangle_p+\delta}}.$$
  The radial lower bound (\cref{lem:dc-ch3-3-5-radial}) at the pair
  $(w(t),w'(t))$ gives
  $$\langle w,w'\rangle_p^2\;\le\;\langle w,w\rangle_p\,\big|\dot c\big|^2
    \;\le\;\big(\langle w,w\rangle_p+\delta\big)\,\big|\dot c\big|^2,$$
  the second step because $|\dot c|^2\ge 0$ (the metric is positive
  semidefinite at every point of the chart image). Taking square roots and
  dividing, $r_\delta'(t)\le|\dot c(t)|$ for \emph{every} $t\in[a,b]$ —
  including the zeros of $w$, where the radius $r$ has a corner but
  $r_\delta$ does not. All the data are continuous ($w$, $w'$, the Gram
  entries along $c$, and $(d\exp_p)$ on the ball by $C^1$ regularity,
  \cref{lem:dc-ch3-2-9-c1ball}), so by the fundamental theorem of calculus
  and monotonicity of the integral
  $$r_\delta(b)-r_\delta(a)=\int_a^b r_\delta'(t)\,dt
    \le\int_a^b\big|\dot c(t)\big|\,dt.$$
  Letting $\delta\to0^+$ gives the claim.
\end{proof}

\begin{lemma}[reach estimate: length dominates the radius reached]
  \label{lem:dc-ch3-3-6-reach}
  \lean{Riemannian.Exponential.gauss_radius_reach, Riemannian.Exponential.exists_gauss_radius_reach_ball}
  \uses{lem:dc-ch3-3-6-radius}
  \leanok
  \dcref{ch3:3.6}
  On the Gauss ball: a curve $c=\exp_p\circ w$ starting at $p$ (i.e.\
  $w(a)=0$) has chart-read length at least the $g_p$-radius its polar lift
  reaches at \emph{any} time:
  $$\sqrt{\langle w(t_1),w(t_1)\rangle_p}\;\le\;\int_a^b|\dot c(t)|\,dt
    \qquad(t_1\in[a,b]).$$
  This is the inequality behind the escape case in
  \cref{prop:dc-ch3-3-6}: a competitor leaving the normal ball of
  $g_p$-radius $\rho'$ reaches radius $\rho'$ at its first exit, so its
  length is at least $\rho'>\ell(\gamma)$.
\end{lemma}
\begin{proof}
  \sketch
  \leanok
  Apply the smoothed-radius comparison (\cref{lem:dc-ch3-3-6-radius}) on
  $[a,t_1]$ and enlarge the integration interval: the speed integrand is
  pointwise nonnegative, so $\int_a^{t_1}\le\int_a^b$. With $w(a)=0$ the
  radius gain is the radius reached.
\end{proof}

\begin{lemma}[radial geodesics have unit chart speed]
  \label{lem:dc-ch3-3-6-rayspeed}
  \lean{Riemannian.Exponential.exists_expMap_ray_speed_ball}
  \uses{lem:dc-ch3-3-5-ball, lem:dc-ch3-3-5-rayode, def:dc-ch3-expmap}
  \leanok
  \dcref{ch3:3.6}
  There is $\rho>0$ such that for every $v$ with $\|v\|<\rho$ and every
  $t\in[0,1]$, read in the chart at $p$,
  $$\big\langle(d\exp_p)_{tv}(v),\,(d\exp_p)_{tv}(v)\big\rangle_{\exp_p(tv)}
    =\langle v,v\rangle_p.$$
  The radial geodesic $\gamma(t)=\exp_p(tv)$, $0\le t\le1$, therefore has
  constant speed $|v|_p$ and length
  $\ell(\gamma)=\int_0^1|v|_p\,dt=|v|_p$ (the normalized identity
  $|\partial f/\partial r|=1$).
\end{lemma}
\begin{proof}
  \sketch
  \leanok
  For $t\ne0$ apply the Gauss identity (\cref{lem:dc-ch3-3-5-ball}) at the
  pair $(tv,v)$:
  $$\big\langle(d\exp_p)_{tv}(tv),\,(d\exp_p)_{tv}(v)\big\rangle_{\exp_p(tv)}
    =\langle tv,v\rangle_p.$$
  Pulling the scalar $t$ out of both sides — by linearity of $(d\exp_p)_{tv}$
  on the left, by bilinearity of the metric on the right — and cancelling
  $t\ne0$ gives the claim. At $t=0$ the ray package
  (\cref{lem:dc-ch3-3-5-rayode}) supplies $(d\exp_p)_0=\mathrm{id}$ and
  $\exp_p(0)=p$, so the left-hand side is $\langle v,v\rangle_p$ directly.

\end{proof}

\begin{lemma}[radial geodesics minimize, polar form]
  \label{lem:dc-ch3-3-6-ball}
  \lean{Riemannian.Exponential.exists_minimizing_geodesic_ball}
  \uses{lem:dc-ch3-3-6-radius, lem:dc-ch3-3-6-rayspeed}
  \leanok
  \dcref{ch3:3.6}
  There is $\rho>0$ (a Gauss ball for $\exp_p$) such that for every $v$ with
  $\|v\|<\rho$ and every competing curve given in polar form through the
  exponential chart — a differentiable $w:[0,1]\to B_\rho(0)\subset T_pM$ with
  continuous derivative, $w(0)=0$, $w(1)=v$ — the chart-read length of the
  radial geodesic $\gamma(t)=\exp_p(tv)$ is at most that of
  $c(t)=\exp_p(w(t))$:
  $$\ell(\gamma)=\int_0^1\sqrt{\langle\dot\gamma,\dot\gamma\rangle}\,dt
    \;\le\;\int_0^1\sqrt{\langle\dot c,\dot c\rangle}\,dt=\ell(c).$$
\end{lemma}
\begin{proof}
  \sketch
  \leanok
  By \cref{lem:dc-ch3-3-6-rayspeed}, $\ell(\gamma)=|v|_p$. By the radius
  comparison (\cref{lem:dc-ch3-3-6-radius}),
  $$\ell(c)\;\ge\;r(1)-r(0)=|w(1)|_p-|w(0)|_p=|v|_p-0=\ell(\gamma).
  $$
\end{proof}

\begin{lemma}[radial geodesics minimize in a normal ball]
  \label{lem:dc-ch3-3-6-normal}
  \lean{Riemannian.Exponential.exists_minimizing_geodesic_normal_ball}
  \uses{lem:dc-ch3-3-6-ball, lem:dc-ch3-2-9-c1diffeo}
  \leanok
  \dcref{ch3:3.6}
  There is $\varepsilon>0$ such that $\exp_p$ is injective on
  $B_\varepsilon(0)\subset T_pM$ with open image — so
  $B=\exp_p(B_\varepsilon(0))$ is a normal ball — and for every $v$ with
  $\|v\|<\varepsilon$ and every curve $c:[0,1]\to B$ with $c(0)=p$,
  $c(1)=\exp_p(v)$, whose chart reading $t\mapsto\varphi_p(c(t))$ is
  differentiable with continuous derivative, the radial geodesic
  $\gamma(t)=\exp_p(tv)$ satisfies $\ell(\gamma)\le\ell(c)$ (chart-read
  lengths).
\end{lemma}
\begin{proof}
  \sketch
  \leanok
  Put the competing curve in polar form: $w=\exp_p^{-1}\circ c$ through the
  local $C^1$ inverse of $\exp_p$ (\cref{lem:dc-ch3-2-9-c1diffeo}), so
  $c=\exp_p\circ w$ with $w$ differentiable into $B_\varepsilon(0)$,
  $w(0)=0$, $w(1)=v$. The image $\varphi_p(\exp_p(B_\varepsilon(0)))$ is open
  — at every point of the ball the derivative of the chart reading of
  $\exp_p$ is a linear equivalence (\cref{lem:dc-ch3-2-9-invertible}), so the
  map is open — hence the chain rule applies to $w=\exp_p^{-1}\circ c$ and
  $w'$ is continuous. The polar minimizing property
  (\cref{lem:dc-ch3-3-6-ball}) gives $\ell(\gamma)\le\ell(\exp_p\circ w)$,
  and $\ell(\exp_p\circ w)=\ell(c)$: the two length integrands agree, the
  feet by $\exp_p(w(t))=c(t)$ and the velocities
  $(d\exp_p)_{w(t)}w'(t)=\frac{d}{dt}\varphi_p(c(t))$ by uniqueness of
  one-sided derivatives on $[0,1]$ (both compute the derivative of
  $u=\varphi_p\circ c=f\circ w$ within $[0,1]$).
\end{proof}

\begin{lemma}[radial geodesics minimize among piecewise competitors in a normal ball]
  \label{lem:dc-ch3-3-6-piecewise}
  \lean{Riemannian.Exponential.gauss_radius_comparison_piecewise, Riemannian.Exponential.gauss_radius_reach_piecewise, Riemannian.Exponential.exists_minimizing_geodesic_normal_ball_piecewise}
  \uses{def:dc-ch3-3-1, lem:dc-ch3-3-6-normal}
  \leanok
  \dcref{ch3:3.6}
  There is $\varepsilon>0$ such that $\exp_p$ is injective on
  $B_\varepsilon(0)\subset T_pM$ with open image — so
  $B=\exp_p(B_\varepsilon(0))$ is a normal ball — and for every $v$ with
  $\|v\|<\varepsilon$, every partition $0=\tau_0\le\dots\le\tau_n=1$, and
  every continuous curve $c:[0,1]\to B$ from $p$ to $\exp_p(v)$ whose chart
  reading is differentiable on each piece $[\tau_i,\tau_{i+1}]$ in the
  one-sided sense with continuous piece derivatives (a piecewise
  differentiable curve, \cref{def:dc-ch3-3-1}), the chart-read length of the
  radial geodesic $\gamma(t)=\exp_p(tv)$ is at most the total length of $c$:
  $$\ell(\gamma)\;\le\;\sum_{i<n}\int_{\tau_i}^{\tau_{i+1}}
    \big|\dot c(t)\big|_{c(t)}\,dt=\ell(c).$$
  Moreover the underlying radius comparison telescopes: for a continuous,
  piecewise differentiable polar lift $w$ with values in the Gauss ball, both
  the radius gain $|w(\tau_n)|_p-|w(\tau_0)|_p$ and — when $w(\tau_0)=0$ —
  the radius $|w(t_1)|_p$ reached at \emph{any} time $t_1$ are dominated by
  the total length of $\exp_p\circ w$.
\end{lemma}
\begin{proof}
  \sketch
  \leanok
  \uses{lem:dc-ch3-3-6-radius, lem:dc-ch3-3-6-rayspeed, lem:dc-ch3-2-9-c1diffeo, lem:dc-ch3-2-9-invertible}
  Each piece $[\tau_i,\tau_{i+1}]$ satisfies the single-piece smoothed-radius
  comparison (\cref{lem:dc-ch3-3-6-radius}): the smoothed radius
  $r_\delta=\sqrt{\langle w,w\rangle_p+\delta}$ is differentiable on the
  closed piece using only the one-sided derivatives of $w$ at its endpoints —
  exactly what a corner of a piecewise curve provides — and the fundamental
  theorem of calculus on the closed piece gives
  $r_\delta(\tau_{i+1})-r_\delta(\tau_i)\le\int_{\tau_i}^{\tau_{i+1}}|\dot c|$,
  hence the same bound for $r$ in the limit $\delta\to0^+$. Summing over the
  pieces telescopes the radius gains:
  $$|w(\tau_n)|_p-|w(\tau_0)|_p
    =\sum_{i<n}\big(|w(\tau_{i+1})|_p-|w(\tau_i)|_p\big)
    \le\sum_{i<n}\int_{\tau_i}^{\tau_{i+1}}|\dot c|\,dt.$$
  For the reach form at a time $t_1$ in the piece $[\tau_k,\tau_{k+1}]$, apply
  the telescoped comparison to the partition of $[\tau_0,t_1]$ obtained by
  cutting at $t_1$, then enlarge: the dropped terms are integrals of a
  nonnegative integrand. The minimizing statement follows as in
  \cref{lem:dc-ch3-3-6-normal}: the competitor is put in polar form
  $w=\exp_p^{-1}\circ c$ through the local $C^1$ inverse of $\exp_p$
  (\cref{lem:dc-ch3-2-9-c1diffeo}), which is differentiable piece by piece
  with continuous one-sided derivatives by the chain rule
  (\cref{lem:dc-ch3-2-9-invertible}), and the telescoped comparison with
  $w(0)=0$, $w(1)=v$ bounds the total length of $c$ below by
  $|v|_p=\ell(\gamma)$ (\cref{lem:dc-ch3-3-6-rayspeed}).
\end{proof}

\begin{lemma}[the competitor bound: piecewise curves, escape case included]
  \label{lem:dc-ch3-3-6-competitor}
  \lean{Riemannian.Exponential.pathELength_sum_partition, Riemannian.Exponential.exists_le_pathELength, Riemannian.Exponential.exists_le_pathELength_piecewise}
  \uses{def:dc-ch3-3-1, lem:dc-ch3-3-6-piecewise}
  \leanok
  \dcref{ch3:3.6}
  Let $p\in M$. There are $\varepsilon>0$ and a constant $C>0$ (comparing the
  coordinate norm at $p$ with the $g_p$-norm) such that $\exp_p$ is injective
  on $B_\varepsilon(0)\subset T_pM$, open on sub-balls, and — measuring the
  length of a piecewise differentiable curve $\sigma:[0,1]\to M$
  intrinsically, as the sum of the lengths of its differentiable pieces:
  \begin{enumerate}
  \item every piecewise differentiable curve $\sigma$ from $p$ to
    $\exp_p(v)$ with $\|v\|<\varepsilon$ satisfies
    $\ell(\sigma)\ge|v|_p=\ell(\gamma)$, where $\gamma(t)=\exp_p(tv)$ is the
    radial geodesic — the inequality, with competitors allowed to
    leave the normal ball;
  \item every piecewise differentiable curve from $p$ ending outside
    $\exp_p(B_r(0))$, $0<r\le\varepsilon$, has length at least $r/C$ (the
    escape estimate).
  \end{enumerate}
\end{lemma}
\begin{proof}
  \sketch
  \leanok
  \uses{lem:dc-ch3-3-6-piecewise, lem:dc-ch3-2-9-c1diffeo}
  Fix a radius $r'$ so small that on the coordinate ball $B_{r'}(0)$ the
  polar machinery of \cref{lem:dc-ch3-3-6-piecewise} applies and the
  coordinate norm is comparable to the $g_p$-norm, $\|z\|^2\le C^2\,
  \langle z,z\rangle_p$; write $U=\exp_p(B_{r'}(0))$, an open neighborhood of
  $p$. A competitor $\sigma$ either stays in $U$ or leaves it.

  If $\sigma$ stays in $U$, its polar lift $w=\exp_p^{-1}\circ\sigma$ exists
  on all of $[0,1]$, is continuous, and is piecewise differentiable with
  continuous one-sided piece derivatives (\cref{lem:dc-ch3-2-9-c1diffeo});
  the telescoped radius comparison (\cref{lem:dc-ch3-3-6-piecewise}) with
  $w(0)=0$ bounds $\ell(\sigma)$ below by the $g_p$-radius of the lifted
  endpoint — which is $|v|_p$ in case (1), and at least $r/C$ in case (2)
  since the endpoint has coordinate norm $\ge r$.

  If $\sigma$ leaves $U$, the image of the \emph{closed} coordinate ball
  $\overline{B_{r'}(0)}$ under $\exp_p$ is compact, hence closed; the set of
  times at which $\sigma$ lies outside $U$ is closed in $[0,1]$ and avoids
  $t=0$, so it has a smallest element $T>0$, and on $[0,T)$ the curve stays
  in $U$. By continuity $\sigma(T)$ lies in the compact image, and it is not
  in the open image $U$, so its polar coordinate has norm exactly $r'$: the
  curve first exits through the coordinate $r'$-sphere. Truncate the
  partition at $T$ — each piece $[\tau_i,\tau_{i+1}]$ becomes
  $[\min(\tau_i,T),\min(\tau_{i+1},T)]$, which either sits inside the
  original piece or degenerates to the point $T$ — and apply the reach form
  of the telescoped comparison on $[0,T]$: the length accumulated up to $T$
  is at least the $g_p$-radius of the exit point, which the norm comparison
  bounds below by $r'/C$. Choosing $\varepsilon\le r'$ so small that every
  $v$ with $\|v\|<\varepsilon$ has $|v|_p<r'/C$, both cases give (1); (2) is
  the same estimate with $r\le\varepsilon\le r'$.

  Finally, the intrinsic total length of $\sigma$ — the sum of the piece
  lengths — is the manifold path length of $\sigma$ over $[0,1]$, by
  additivity of the length integral over a monotone partition; on each piece
  the intrinsic speed agrees with the chart-read speed used by the polar
  comparison.
\end{proof}

\begin{lemma}[equality forces radial velocity]
  \label{lem:dc-ch3-3-6-equality-radial}
  \lean{Riemannian.Exponential.radial_equality_forcing, Riemannian.Exponential.gauss_radius_equality_ftc, Riemannian.Exponential.exists_gauss_equality_radial_ball}
  \uses{lem:dc-ch3-3-5-ball, lem:dc-ch3-3-6-radius}
  \leanok
  \dcref{ch3:3.6}
  There is $\rho>0$ (a Gauss ball for $\exp_p$) such that for every
  differentiable curve $w:[a,b]\to B_\rho(0)\subset T_pM$ with continuous
  derivative whose composed curve $c=\exp_p\circ w$ realizes the radius
  comparison with \emph{equality},
  $$|w(b)|_p-|w(a)|_p=\int_a^b|\dot c(t)|\,dt,$$
  the following hold at every $t\in(a,b)$ with $w(t)\neq0$: the radius
  $r=|w|_p$ is differentiable at $t$ with $r'(t)=|\dot c(t)|\ge0$, and the
  velocity is radial and outward,
  $$w'(t)=\frac{\langle w(t),w'(t)\rangle_p}{\langle w(t),w(t)\rangle_p}\,
    w(t),\qquad\langle w(t),w'(t)\rangle_p\ge0.$$
\end{lemma}
\begin{proof}
  \sketch
  \leanok
  \uses{lem:dc-ch3-3-5-ball, lem:dc-ch3-3-6-radius, lem:dc-ch3-2-9-invertible}
  The radius comparison holds on every subinterval $[c,d]\subseteq[a,b]$
  (\cref{lem:dc-ch3-3-6-radius}), so the deficit
  $\varphi(t)=\int_a^t|\dot c|-\big(r(t)-r(a)\big)$ is nondecreasing;
  vanishing at both ends, it vanishes identically, giving the FTC form
  $r(t)=r(a)+\int_a^t|\dot c|$ for \emph{all} $t$. The speed $|\dot c|$ is
  continuous, so $r$ is differentiable on $(a,b)$ with $r'=|\dot c|\ge0$.
  Where $w(t)\neq0$, the radius is also differentiable by the product rule
  with $r'=\langle w,w'\rangle_p/r$; identifying the two derivatives gives
  $\langle w,w'\rangle_p=r\,|\dot c|\ge0$, i.e.\ \emph{equality} in the
  radial lower bound
  $\langle w,w'\rangle_p^2\le\langle w,w\rangle_p\,|\dot c|^2$.

  Equality there forces the velocity to be radial. Decompose
  $w'=\lambda w+w_N'$ with
  $\lambda=\langle w,w'\rangle_p/\langle w,w\rangle_p$ and
  $\langle w,w_N'\rangle_p=0$. Expanding
  $|\dot c|^2=\big|(d\exp_p)_w(w')\big|^2$ by bilinearity and applying the
  Gauss identity (\cref{lem:dc-ch3-3-5-ball}) twice kills the cross terms:
  $$\big|(d\exp_p)_w(w')\big|^2
    =\lambda^2\langle w,w\rangle_p+\big|(d\exp_p)_w(w_N')\big|^2.$$
  Equality in the radial bound then forces
  $\big|(d\exp_p)_w(w_N')\big|^2=0$; the metric is positive definite at
  every point of a small chart ball around $p$ (the Gram comparison
  $\|z\|^2\le c\,\langle z,z\rangle_y$ holds for all $y$ near the pole), so
  $(d\exp_p)_w(w_N')=0$, and $(d\exp_p)_w$ is invertible on the ball
  (\cref{lem:dc-ch3-2-9-invertible}), so $w_N'=0$.
\end{proof}

\begin{lemma}[equality case in polar form: the competitor runs along the ray]
  \label{lem:dc-ch3-3-6-equality-ray}
  \lean{Riemannian.Exponential.exists_gauss_equality_ray_ball}
  \uses{lem:dc-ch3-3-6-equality-radial}
  \leanok
  \dcref{ch3:3.6}
  There is $\rho>0$ (a Gauss ball for $\exp_p$) such that every
  differentiable polar competitor $w:[0,1]\to B_\rho(0)\subset T_pM$ with
  continuous derivative and $w(0)=0$ that realizes the radius comparison with
  \emph{equality}, $|w(1)|_p=\int_0^1|\dot c|\,dt$ where $c=\exp_p\circ w$,
  satisfies, for every $t\in[0,1]$,
  $$w(t)=\frac{|w(t)|_p}{|w(1)|_p}\,w(1):$$
  the competitor traces the radial segment from $0$ to $w(1)$, with monotone
  radius $|w(t)|_p$.
\end{lemma}
\begin{proof}
  \sketch
  \leanok
  \uses{lem:dc-ch3-3-6-equality-radial}
  If $|w(1)|_p=0$ the monotone radius (its derivative is the nonnegative
  speed, by the FTC form of \cref{lem:dc-ch3-3-6-equality-radial}) is
  squeezed to $0$, so $w\equiv0$ by positive definiteness of
  $\langle\cdot,\cdot\rangle_p$ and both sides vanish. Otherwise let
  $t_*=\sup\{t:|w(t)|_p=0\}$, attained since the radius is continuous; then
  $t_*<1$ and $r=|w|_p$ is positive on $(t_*,1]$. On $(t_*,1)$ the
  normalized lift $\eta=w/r$ is differentiable with
  $$\eta'=\frac{w'}{r}-\frac{r'}{r^2}\,w=0,$$
  combining $r'=\langle w,w'\rangle_p/r$ (product rule) with the radial
  velocity $w'=(\langle w,w'\rangle_p/\langle w,w\rangle_p)\,w$
  (\cref{lem:dc-ch3-3-6-equality-radial}); a continuous function with
  vanishing right derivative on $[t,1]$ is constant, so $\eta\equiv\eta(1)$
  on $(t_*,1]$, i.e.\ $w(t)=r(t)\,w(1)/r(1)$ there. On $[0,t_*]$ the
  monotone radius satisfies $0\le r(t)\le r(t_*)=0$, so $w(t)=0$ and the
  identity holds with coefficient $0$.
\end{proof}

\begin{lemma}[confinement: a short curve from $p$ stays in the Gauss ball]
  \label{lem:dc-ch3-3-6-confinement}
  \lean{Riemannian.Exponential.exists_forall_mem_expMap_ball_of_pathELength_lt}
  \uses{lem:dc-ch3-3-6-competitor}
  \leanok
  \dcref{ch3:3.6}
  For every target radius $\rho'>0$ there is a length threshold $\delta>0$
  such that every $C^1$ curve $\sigma:[0,1]\to M$ starting at $p$ with
  $\ell(\sigma)<\delta$ satisfies $\sigma(t)\in\exp_p(B_{\rho'}(0))$ for all
  $t\in[0,1]$. This discharges the escape case in the equality analysis of
  \cref{prop:dc-ch3-3-6}: an equality competitor is automatically confined
  to the Gauss ball.
\end{lemma}
\begin{proof}
  \sketch
  \leanok
  \uses{lem:dc-ch3-3-6-competitor}
  Truncate $\sigma$ at time $t$ and reparametrize the truncation to $[0,1]$
  by the monotone affine map $\tau\mapsto t\tau$, which leaves the length
  unchanged. If the truncated endpoint $\sigma(t)$ escaped the sublevel set
  $\exp_p(\{z:\ \|z\|<\varepsilon,\ |z|_p<\delta\})$, the escape estimate of
  \cref{lem:dc-ch3-3-6-competitor} would force
  $\delta\le\ell(\sigma|_{[0,t]})\le\ell(\sigma)<\delta$, a contradiction.
  The Gram comparison $\|z\|^2\le c\,\langle z,z\rangle_p$ places the
  sublevel set inside the norm ball of radius $\sqrt{c}\,\delta$, which is
  inside $B_{\rho'}(0)$ for $\delta=\rho''/(\sqrt{c}+1)$,
  $\rho''=\min(\rho',\varepsilon)$.
\end{proof}

\begin{lemma}[the equality case in the manifold: monotone reparametrization of the radial geodesic]
  \label{lem:dc-ch3-3-6-equality-manifold}
  \lean{Riemannian.Exponential.exists_gauss_equality_manifold_ball, Riemannian.Exponential.exists_gauss_equality_manifold}
  \uses{lem:dc-ch3-3-6-equality-ray, lem:dc-ch3-3-6-equality-radial, lem:dc-ch3-2-9-c1diffeo, lem:dc-ch7-2-8-chart-length, lem:dc-ch3-3-6-confinement}
  \leanok
  \dcref{ch3:3.6}
  There is $\rho>0$ (a Gauss ball for $\exp_p$, with $\exp_p$ injective on
  $B_\rho(0)$) such that every $C^1$ competitor $\sigma:[0,1]\to M$ from $p$
  to $\exp_p(v)$, $|v|<\rho$, that stays in $\exp_p(B_\rho(0))$ and realizes
  the radial length $\ell(\sigma)=|v|_p$ is a \emph{monotone reparametrization
  of the radial geodesic}: there is a continuous monotone
  $s:[0,1]\to[0,1]$ with $s(0)=0$, $s(1)=1$,
  $$\sigma(t)=\exp_p\big(s(t)\,v\big)\quad\text{and}\quad
    \ell\big(\sigma|_{[0,t]}\big)=s(t)\,|v|_p\quad\text{for all }t\in[0,1];$$
  in particular $\sigma([0,1])=\exp_p([0,1]\cdot v)$. By the confinement lemma
  (\cref{lem:dc-ch3-3-6-confinement}) the stay-in-ball hypothesis is
  automatic: the statement holds for \emph{every} $C^1$ equality competitor
  from $p$ to $\exp_p(v)$, with no assumption on where it travels.
\end{lemma}
\begin{proof}
  \sketch
  \leanok
  \uses{lem:dc-ch3-3-6-equality-ray, lem:dc-ch3-3-6-equality-radial, lem:dc-ch3-2-9-c1diffeo, lem:dc-ch7-2-8-chart-length, lem:dc-ch3-3-6-confinement}
  Lift $\sigma$ through the $C^1$ local inverse of $\exp_p$
  (\cref{lem:dc-ch3-2-9-c1diffeo}): the polar lift
  $w=\exp_p^{-1}\circ\sigma$ is $C^1$ on $[0,1]$ with $w(0)=0$, $w(1)=v$, and
  values in $B_\rho(0)$. Reading the length in the chart at $p$
  (\cref{lem:dc-ch7-2-8-chart-length}) and computing the chart velocity of
  $\sigma$ by the chain rule through $w$, the hypothesis
  $\ell(\sigma)=|v|_p$ becomes the FTC-form equality of the polar lemmas.
  The ray lemma (\cref{lem:dc-ch3-3-6-equality-ray}) then forces
  $w(t)=\big(r(t)/r(1)\big)\,w(1)$ with $r=|w|_p$, and the radial lemma
  (\cref{lem:dc-ch3-3-6-equality-radial}) makes $r$ monotone — its derivative
  on $(0,1)$ is the nonnegative speed — and continuous, so
  $s=r/r(1)$ is the desired reparametrization (when $v=0$ the lift is
  identically $0$ and $s(t)=t$ works). The running-length identity
  $\ell(\sigma|_{[0,t]})=r(t)$ is the fundamental theorem of calculus for
  $r'=|\dot\sigma|$, and the image equality follows from the intermediate
  value theorem applied to the continuous monotone radius.
\end{proof}

\begin{lemma}[an arclength-proportional minimizer is the radial geodesic]
  \label{lem:dc-ch3-3-6-equality-geodesic}
  \lean{Riemannian.Exponential.exists_gauss_equality_geodesic_ball, Riemannian.Exponential.exists_gauss_equality_geodesic}
  \uses{lem:dc-ch3-3-6-equality-manifold, lem:dc-ch7-2-8-ray-geodesic, lem:dc-ch3-3-6-confinement}
  \leanok
  \dcref{ch3:3.6}
  There is $\rho>0$ such that every $C^1$ curve $\sigma:[0,1]\to M$ from $p$
  to $\exp_p(v)$, $|v|<\rho$, staying in $\exp_p(B_\rho(0))$, that is
  \emph{parametrized proportionally to arc length while realizing the radial
  length} — $\ell(\sigma|_{[0,t]})=t\,|v|_p$ for all $t\in[0,1]$ — coincides
  with the radial geodesic with matching parametrization:
  $\sigma(t)=\exp_p(t\,v)$ for all $t\in[0,1]$, and $\sigma$ satisfies the
  intrinsic geodesic equation on $(0,1)$. This is the normal-ball core of
  \cref{cor:dc-ch3-3-9}: the equality case of the minimizing property turns
  an arclength-proportional minimizer into a geodesic. By the confinement
  lemma (\cref{lem:dc-ch3-3-6-confinement}) the stay-in-ball hypothesis is
  again automatic.
\end{lemma}
\begin{proof}
  \sketch
  \leanok
  \uses{lem:dc-ch3-3-6-equality-manifold, lem:dc-ch7-2-8-ray-geodesic, lem:dc-ch3-3-6-confinement}
  The running-length hypothesis at $t=1$ realizes the radial length, so
  \cref{lem:dc-ch3-3-6-equality-manifold} yields the monotone
  reparametrization $s$ with $\sigma(t)=\exp_p(s(t)v)$ and
  $\ell(\sigma|_{[0,t]})=s(t)|v|_p$. Comparing with the hypothesis
  $\ell(\sigma|_{[0,t]})=t\,|v|_p$ and cancelling $|v|_p>0$ (for $v\ne0$;
  the case $v=0$ is the constant curve) pins $s(t)=t$. The geodesic equation
  transfers from the radial geodesic (\cref{lem:dc-ch7-2-8-ray-geodesic}) at
  every interior time by locality of the moving-foot equation, since
  $\sigma$ agrees with the ray on a neighborhood of each $t\in(0,1)$.

\end{proof}

\begin{lemma}[equality case in polar form, piecewise: the competitor runs along the ray]
  \label{lem:dc-ch3-3-6-equality-ray-piecewise}
  \lean{Riemannian.Exponential.exists_gauss_equality_ray_interval, Riemannian.Exponential.exists_gauss_equality_ray_piecewise, Riemannian.Exponential.monotoneOn_of_partition}
  \uses{lem:dc-ch3-3-6-equality-radial, lem:dc-ch3-3-6-equality-ray, lem:dc-ch3-3-6-radius}
  \leanok
  \dcref{ch3:3.6}
  There is $\rho>0$ (a Gauss ball for $\exp_p$) such that:
  \begin{itemize}
  \item (\emph{general interval, no anchoring}) every differentiable polar
    competitor $w:[a,b]\to B_\rho(0)\subset T_pM$ with continuous derivative
    that realizes the radius comparison with \emph{equality},
    $|w(b)|_p-|w(a)|_p=\int_a^b|\dot c|\,dt$ where $c=\exp_p\circ w$, has
    monotone radius $r=|w|_p$ and satisfies
    $w(t)=\bigl(r(t)/r(b)\bigr)\,w(b)$ for every $t\in[a,b]$ — the lift is
    \emph{not} required to start at the origin;
  \item (\emph{piecewise}) every polar competitor on a partition
    $\tau_0\le\dots\le\tau_n$ — continuous on the span, $C^1$ on each piece —
    realizing the total radius comparison with equality has monotone radius
    on the span, satisfies the equality in FTC form on every piece,
    $r(t)=r(\tau_i)+\int_{\tau_i}^t|\dot c|$, and traces the ray of its
    final point: $w(t)=\bigl(r(t)/r(\tau_n)\bigr)\,w(\tau_n)$ for every $t$
    in the span.
  \end{itemize}
\end{lemma}
\begin{proof}
  \sketch
  \leanok
  \uses{lem:dc-ch3-3-6-equality-radial, lem:dc-ch3-3-6-equality-ray, lem:dc-ch3-3-6-radius}
  For the general interval: the argument of
  \cref{lem:dc-ch3-3-6-equality-ray} never used the anchoring $w(a)=0$
  beyond producing a zero of the radius. If the radius never vanishes on
  $[a,b]$, the normalized lift $\eta=w/r$ has vanishing derivative on the
  interior — combining $r'=\langle w,w'\rangle_p/r$ with the radial velocity
  of \cref{lem:dc-ch3-3-6-equality-radial} — hence is constant on $(a,b]$,
  and takes the same value in the limit at $a$ by continuity of $w$ and $r$;
  so $w(t)=r(t)\,\eta(b)=\bigl(r(t)/r(b)\bigr)w(b)$ throughout. If it
  vanishes somewhere, let $t_*$ be its last zero: on $[a,t_*]$ the monotone
  radius is squeezed to $0$, so $w=0$ and the identity holds with
  coefficient $0$; on $(t_*,b]$ the previous argument applies verbatim.

  For the piecewise form: the radius comparison of
  \cref{lem:dc-ch3-3-6-radius} holds on every piece, and the per-piece
  radius gains telescope to $r(\tau_n)-r(\tau_0)$; total equality therefore
  forces \emph{equality on every piece} (a sum of nonnegative deficits
  vanishes only if each does). Each piece then has monotone radius and
  traces the ray of its right endpoint by the general-interval case, the
  per-piece monotone radii glue along the partition, and a backward
  induction over the pieces composes the scalar factors:
  $$\frac{r(t)}{r(\tau_{i+1})}\cdot\frac{r(\tau_{i+1})}{r(\tau_n)}
    =\frac{r(t)}{r(\tau_n)},$$
  the degenerate case $r(\tau_{i+1})=0$ being absorbed by monotonicity
  ($r(t)=0$ forces $w(t)=0$). No gluing of directions across the corners is
  needed.
\end{proof}

\begin{lemma}[confinement, piecewise: a short piecewise curve from $p$ stays in the Gauss ball]
  \label{lem:dc-ch3-3-6-confinement-piecewise}
  \lean{Riemannian.Exponential.exists_forall_mem_expMap_ball_of_pathELength_lt_piecewise}
  \uses{lem:dc-ch3-3-6-competitor, lem:dc-ch3-3-6-confinement}
  \leanok
  \dcref{ch3:3.6}
  For every target radius $\rho'>0$ there is a length threshold $\delta>0$
  such that every piecewise-$C^1$ curve $\sigma:[0,1]\to M$ starting at $p$
  with $\ell(\sigma)<\delta$ satisfies $\sigma(t)\in\exp_p(B_{\rho'}(0))$
  for all $t\in[0,1]$: the piecewise analogue of
  \cref{lem:dc-ch3-3-6-confinement}, discharging the escape case of the
  piecewise equality analysis.
\end{lemma}
\begin{proof}
  \sketch
  \leanok
  \uses{lem:dc-ch3-3-6-competitor}
  Truncate $\sigma$ at time $t$ along the truncated partition
  $\min(\tau_j,t)$ and reparametrize each truncated piece affinely to a
  piecewise-$C^1$ curve on $[0,1]$ (collapsed pieces are constants); by
  additivity of the length over the partition and invariance of each
  piece's length under its monotone affine reparametrization, the
  reparametrized truncation has length
  $\ell(\sigma|_{[0,t]})\le\ell(\sigma)$. If the endpoint $\sigma(t)$
  escaped $\exp_p(B_{\rho''}(0))$ with $\rho''=\min(\rho',\varepsilon)$, the
  piecewise escape estimate of \cref{lem:dc-ch3-3-6-competitor} would force
  $\rho''/\sqrt{c}\le\ell(\sigma|_{[0,t]})<\delta=\rho''/\sqrt{c}$, a
  contradiction.
\end{proof}

\begin{lemma}[the equality case in the manifold, piecewise competitors]
  \label{lem:dc-ch3-3-6-equality-manifold-piecewise}
  \lean{Riemannian.Exponential.exists_gauss_equality_manifold_piecewise_ball, Riemannian.Exponential.exists_gauss_equality_manifold_piecewise}
  \uses{lem:dc-ch3-3-6-equality-ray-piecewise, lem:dc-ch3-3-6-confinement-piecewise, lem:dc-ch3-2-9-c1diffeo, lem:dc-ch7-2-8-chart-length, lem:dc-ch3-3-6-equality-manifold}
  \leanok
  \dcref{ch3:3.6}
  There is $\rho>0$ (a Gauss ball for $\exp_p$, with $\exp_p$ injective on
  $B_\rho(0)$) such that every \emph{piecewise-$C^1$} competitor
  $\sigma:[0,1]\to M$ from $p$ to $\exp_p(v)$, $|v|<\rho$ — continuous on
  $[0,1]$, $C^1$ on each piece of a partition $0=\tau_0\le\dots\le\tau_n=1$ —
  that realizes the radial length $\ell(\sigma)=|v|_p$ is a \emph{monotone
  reparametrization of the radial geodesic}: there is a continuous monotone
  $s:[0,1]\to[0,1]$ with $s(0)=0$, $s(1)=1$,
  $$\sigma(t)=\exp_p\big(s(t)\,v\big)\quad\text{and}\quad
    \ell\big(\sigma|_{[0,t]}\big)=s(t)\,|v|_p\quad\text{for all }t\in[0,1];$$
  in particular $\sigma([0,1])=\exp_p([0,1]\cdot v)$ in its literal regularity
  class. By the piecewise
  confinement lemma (\cref{lem:dc-ch3-3-6-confinement-piecewise}) no
  stay-in-ball hypothesis is needed.
\end{lemma}
\begin{proof}
  \sketch
  \leanok
  \uses{lem:dc-ch3-3-6-equality-ray-piecewise, lem:dc-ch3-3-6-confinement-piecewise, lem:dc-ch3-2-9-c1diffeo, lem:dc-ch7-2-8-chart-length}
  As in the $C^1$ case (\cref{lem:dc-ch3-3-6-equality-manifold}), lift
  $\sigma$ through the $C^1$ local inverse of $\exp_p$
  (\cref{lem:dc-ch3-2-9-c1diffeo}): the polar lift $w=\exp_p^{-1}\circ\sigma$
  is continuous on $[0,1]$ and $C^1$ on each piece, with $w(0)=0$,
  $w(1)=v$. Reading the length in the chart at $p$
  (\cref{lem:dc-ch7-2-8-chart-length}) piece by piece and summing over the
  partition, the hypothesis $\ell(\sigma)=|v|_p$ becomes the total FTC-form
  equality of the piecewise polar lemma
  (\cref{lem:dc-ch3-3-6-equality-ray-piecewise}), which forces
  $w(t)=\bigl(r(t)/r(1)\bigr)\,v$ with $r=|w|_p$ monotone; so $s=r/r(1)$ is
  the desired reparametrization (for $v=0$ the lift is identically $0$ and
  $s(t)=t$ works). The running-length identity
  $\ell(\sigma|_{[0,t]})=r(t)$ telescopes over the truncated partition
  $\min(\tau_j,t)$, evaluating each truncated piece by the per-piece FTC
  clause; the image equality is the intermediate value theorem for the
  continuous monotone radius.
\end{proof}

\begin{proposition}[geodesics locally minimize]
  \label{prop:dc-ch3-3-6}
  \lean{Riemannian.Exponential.exists_le_pathELength_piecewise, Riemannian.Exponential.exists_gauss_equality_manifold_piecewise}
  \uses{lem:dc-ch3-3-6-normal, lem:dc-ch3-3-6-piecewise, lem:dc-ch3-3-6-competitor, lem:dc-ch3-3-6-equality-radial, lem:dc-ch3-3-6-equality-ray, lem:dc-ch3-3-6-equality-manifold, lem:dc-ch3-3-6-equality-geodesic, lem:dc-ch3-3-6-confinement, lem:dc-ch3-3-6-equality-ray-piecewise, lem:dc-ch3-3-6-confinement-piecewise, lem:dc-ch3-3-6-equality-manifold-piecewise, prop:dc-ch3-2-9}
  \leanok
  \dcref{ch3:3.6}
  Let $p\in M$, $U$ a normal neighborhood of $p$, and $B\subset U$ a normal ball of
  center $p$. Let $\gamma:[0,1]\to B$ be a geodesic segment with $\gamma(0)=p$. If
  $c:[0,1]\to M$ is any piecewise differentiable curve joining $\gamma(0)$ to
  $\gamma(1)$, then $\ell(\gamma)\le\ell(c)$, and if equality holds then
  $\gamma([0,1])=c([0,1])$.
\end{proposition}
\begin{proof}
  \sketch
  \leanok
  \uses{lem:dc-ch3-3-6-competitor, lem:dc-ch3-3-6-equality-manifold-piecewise, lem:dc-ch3-3-6-piecewise, lem:dc-ch3-3-6-radius, lem:dc-ch3-3-6-rayspeed, lem:dc-ch3-3-6-equality-radial, lem:dc-ch3-3-6-equality-ray, lem:dc-ch3-3-6-equality-manifold, lem:dc-ch3-3-6-equality-geodesic, lem:dc-ch3-3-6-confinement, lem:dc-ch3-3-6-equality-ray-piecewise, lem:dc-ch3-3-6-confinement-piecewise}
  Suppose first $c([0,1])\subset B$. For $t\ne 0$ write uniquely
  $c(t)=\exp_p(r(t)\,v(t))=f(r(t),t)$, $|v(t)|=1$, $r:(0,1]\to\mathbb{R}$ positive
  piecewise differentiable. Then except at finitely many points
  $\frac{dc}{dt}=\frac{\partial f}{\partial r}r'(t)+\frac{\partial f}{\partial t}$.
  By the Gauss lemma $\langle\frac{\partial f}{\partial r},\frac{\partial f}{\partial t}\rangle=0$
  and $|\frac{\partial f}{\partial r}|=1$, so
  
  $$
  \Big|\frac{dc}{dt}\Big|^2=|r'(t)|^2+\Big|\frac{\partial f}{\partial t}\Big|^2\ge|r'(t)|^2,\qquad\text{(1)}
  $$
  
  $$
  \int_\varepsilon^1\Big|\frac{dc}{dt}\Big|dt\ge\int_\varepsilon^1|r'(t)|dt\ge r(1)-r(\varepsilon).\qquad\text{(2)}
  $$
  
  Letting $\varepsilon\to0$ gives $\ell(c)\ge r(1)=\ell(\gamma)$. If equality holds
  then $|\frac{\partial f}{\partial t}|=0$ and $r'(t)>0$, so $c$ is a monotone
  reparametrization of $\gamma$ and $c([0,1])=\gamma([0,1])$. If $c([0,1])\not\subset B$,
  let $t_1$ be the first time $c(t_1)\in\partial B$; with $\rho$ the radius of $B$,
  $\ell(c)\ge\ell_{[0,t_1]}(c)\ge\rho>\ell(\gamma)$.
\end{proof}

\begin{lemma}[geodesic segments of the chart flow, from any base point]
  \label{lem:dc-ch3-3-7-descent}
  \lean{Riemannian.Exponential.isGeodesicOn_uniform_flow_segment}
  \uses{def:dc-ch3-2-1, lem:dc-ch3-2-7-vuniform, lem:dc-ch7-2-8-chart-independence}
  \leanok
  \dcref{ch3:3.7}
  Let $p\in M$, let $\mathbf{x}$ be the chart at $p$, and let $Z$ be a local flow
  of the coordinate geodesic equation as in \cref{lem:dc-ch3-2-7-vuniform}: for
  every initial condition $\zeta$ in a ball around $(\mathbf{x}(p),0)$, the curve
  $Z(\zeta)$ solves $z'=(z_2,\,-\Gamma_p(z_2,z_2)(z_1))$ on
  $[-\varepsilon,\varepsilon]$ with $Z(\zeta)(0)=\zeta$ and position component in
  the chart image. Let $0<T<\varepsilon$ and let $(y,w/T)$ be an initial
  condition in the ball. Then the curve
  $\gamma(s)=\mathbf{x}^{-1}\bigl(Z(y,w/T)(sT)_1\bigr)$, $s\in[0,1]$,

  \begin{itemize}
  \item starts at $\gamma(0)=\mathbf{x}^{-1}(y)$ and is continuous;
  \item stays in the chart domain at $p$, where its $\mathbf{x}$-reading is the
    position component of the flow line;
  \item has $\mathbf{x}$-velocity $w$ at $s=0$;
  \item is a geodesic on $[0,1]$: it satisfies the geodesic equation of
    \cref{def:dc-ch3-2-1} at every $s$, read in the chart at its own foot
    $\gamma(s)$;
  \item has $\mathbf{x}$-acceleration
    $x''(0)=-\Gamma_p(w,w)(y)=\bigl(\mathrm{spray}(y,w)\bigr)_2$ at $s=0$ — the
    fixed-chart-at-$p$ reading of the geodesic equation $x''+\Gamma(x',x')=0$,
    read algebraically off the spray with no time integration.
  \end{itemize}

  The base point $\mathbf{x}^{-1}(y)$ ranges over the whole chart-image ball —
  the initial condition is not required to lie over $p$ — so a single flow
  provides geodesic segments with a uniform velocity ball at every base point
  near $p$.
\end{lemma}
\begin{proof}
  \sketch
  \leanok
  Write $u(s)=Z(y,w/T)(sT)_1$ and $v(s)=Z(y,w/T)(sT)_2$ for the position and
  velocity components of the reparametrized flow line. By the chain rule
  $u'(s)=T\,v(s)$ and $v'(s)=-T\,\Gamma_p(v(s),v(s))(u(s))$ on the open window
  $|s|<\varepsilon/T\supset[0,1]$, so, the Christoffel contraction being
  quadratic in its velocity slot,
  $$
  u''(s)=T\,v'(s)\cdot T/T=-\Gamma_p\bigl(T v(s),T v(s)\bigr)(u(s))
    =-\Gamma_p\bigl(u'(s),u'(s)\bigr)(u(s)):
  $$
  the reading $u$ solves the second-order geodesic equation in the fixed chart
  at $p$. The initial conditions give $u(0)=y$ and $u'(0)=T\cdot w/T=w$; hence
  $u''(0)=-\Gamma_p(w,w)(y)=\bigl(\mathrm{spray}(y,w)\bigr)_2$. Every
  foot $\gamma(s)=\mathbf{x}^{-1}(u(s))$ lies in the chart domain at $p$, so by
  the chart-independence of the geodesic equation
  (\cref{lem:dc-ch7-2-8-chart-independence}) $\gamma$ satisfies the geodesic
  equation read in the chart at $\gamma(s)$ as well. Continuity of $\gamma$
  follows since $u$ is differentiable and $\mathbf{x}^{-1}$ is continuous on the
  chart image.
\end{proof}

\begin{lemma}[interior chart velocity of the chart-flow segment]
  \label{lem:dc-ch3-3-7-descent-velocity}
  \lean{Riemannian.Geodesic.isGeodesicOn_uniform_flow_segment_Ioo}
  \uses{lem:dc-ch3-3-7-descent}
  \leanok
  \dcref{ch3:3.7}
  In the setting of \cref{lem:dc-ch3-3-7-descent}, the flow-geodesic segment
  $\gamma(s)=\mathbf{x}^{-1}\bigl(Z(y,w/T)(sT)_1\bigr)$ has a well-defined
  $\mathbf{x}$-velocity \emph{at every} interior time, not only at $s=0$: on the
  open window $|s|<\varepsilon/T\supset[0,1]$ one has
  $$
  (\mathbf{x}\circ\gamma)'(s)=T\,Z(y,w/T)(sT)_2,
  $$
  the rescaled velocity component of the flow line. In particular the joining
  geodesic of \cref{thm:dc-ch3-3-7} carries at each interior time $t_0\in(0,1)$ the
  chart-velocity datum $\mathbf{x}$-$\gamma'(t_0)=T\,Z(\cdot)(t_0T)_2$ that the
  convex-neighborhood admissibility \cref{lem:dc-ch3-4-2-interior} consumes —
  closing the interface gap between \cref{thm:dc-ch3-3-7} (whose $s=0$ velocity
  clause is \cref{lem:dc-ch3-3-7-descent}) and the interior deduction.
\end{lemma}
\begin{proof}
  \sketch
  \leanok
  Differentiate the position component $u(s)=Z(y,w/T)(sT)_1$ by the chain rule:
  the flow solves $z'(t)=\bigl(z_2(t),-\Gamma_p(z_2,z_2)(z_1)\bigr)$
  (\cref{lem:dc-ch3-3-7-descent}), so the time-derivative of its first component
  is its second component; the reparametrisation $t=sT$ multiplies the derivative
  by $T$, giving $u'(s)=T\,Z(y,w/T)(sT)_2$ on the whole open window where the flow
  derivative is interior to $[-\varepsilon,\varepsilon]$. This is the chart velocity
  of $\gamma$ since $\mathbf{x}\circ\gamma=u$ on the chart-image window.
\end{proof}

\begin{lemma}[differential of the pair map]
  \label{lem:dc-ch3-3-7-pairmap}
  \lean{Riemannian.Exponential.exists_pairMap_hasStrictFDerivAt}
  \uses{lem:dc-ch3-2-7-vuniform, lem:dc-ch3-2-2-c1dep, lem:dc-ch3-2-9-strict, lem:dc-ch3-3-7-descent}
  \leanok
  \dcref{ch3:3.7}
  In the setting of \cref{lem:dc-ch3-3-7-descent}, the pair map
  $$
  G(y,w)=\bigl(y,\;Z(y,w/T)(T)_1\bigr)
  $$
  — the $\mathbf{x}$-coordinate expression of $F(q,v)=(q,\exp_q v)$, since by
  \cref{lem:dc-ch3-3-7-descent} the point $Z(y,w/T)(T)_1$ is the time-$1$ value
  of the geodesic through $q=\mathbf{x}^{-1}(y)$ with $\mathbf{x}$-velocity $w$ —
  is strictly differentiable at $(\mathbf{x}(p),0)$ with differential
  $$
  dG_{(\mathbf{x}(p),0)}(a,b)=(a,\;a+b),
  $$
  the block matrix $\begin{pmatrix}I & 0\\ I & I\end{pmatrix}$. In
  particular the differential is an invertible linear map.
\end{lemma}
\begin{proof}
  \sketch
  \leanok
  The point $(\mathbf{x}(p),0)$ is an equilibrium of the geodesic spray
  $F(x,u)=(u,\,-\Gamma_p(u,u)(x))$, and $F$ is $C^1$ near it. By the
  $C^1$-dependence theorem for solution families
  (\cref{lem:dc-ch3-2-2-c1dep}, applied at the equilibrium exactly as in
  \cref{lem:dc-ch3-2-9-strict}), the map sending an initial condition $\zeta$
  near the equilibrium to the solution curve $Z(\zeta)|_{[0,T]}$ is strictly
  differentiable there, with differential $D$ the solution of the variational
  equation $Dv(t)-\int_0^t A\,(Dv)(s)\,ds=v$ along the constant operator curve
  $A=dF_{(\mathbf{x}(p),0)}$. Since $\Gamma_p$ is quadratic in the velocity,
  $A(a,b)=(b,0)$; hence $A^2=0$ and $Dv(t)=v+tAv$ solves the variational
  equation. The pair map is the composition of $\zeta\mapsto Z(\zeta)|_{[0,T]}$
  with the linear rescaling $(a,b)\mapsto(a,b/T)$, the evaluation at time $T$
  and the projection to the position component, so
  $$
  dG_{(\mathbf{x}(p),0)}(a,b)
    =\Bigl(a,\;\bigl[(a,b/T)+T\,A(a,b/T)\bigr]_1\Bigr)
    =\bigl(a,\;a+T\cdot b/T\bigr)=(a,\;a+b).
  $$
  The inverse of the differential is $(a,c)\mapsto(a,c-a)$.
\end{proof}

\begin{lemma}[the pair map is $C^1$ on a ball]
  \label{lem:dc-ch3-3-7-c1pair}
  \lean{Riemannian.Exponential.exists_pairMap_contDiffOn}
  \uses{lem:dc-ch3-2-7-c1flow, lem:dc-ch3-3-7-pairmap, lem:dc-ch3-2-9-c1}
  \leanok
  \dcref{ch3:3.7}
  In the setting of \cref{lem:dc-ch3-3-7-pairmap}, the pair map
  $G(y,w)=(y,\,Z(y,w/T)(T)_1)$ is $C^1$ on the open set of admissible initial
  conditions $\{(y,w):\ (y,w/T)\in B((\mathbf{x}(p),0),r)\}$. In particular,
  slicing at each base point: for every $y$ with $|y-\mathbf{x}(p)|<r$, the
  chart exponential $w\mapsto Z(y,w/T)(T)_1$ at the base point
  $\mathbf{x}^{-1}(y)$ is $C^1$ on the uniform velocity ball $|w|<Tr$.
\end{lemma}
\begin{proof}
  \sketch
  \leanok
  By the $C^1$-dependence of the flow on its full initial condition, at an
  arbitrary admissible base solution (\cref{lem:dc-ch3-2-7-c1flow}), the map
  $\zeta\mapsto Z(\zeta)|_{[0,T]}$ is strictly differentiable at every $\zeta$
  in the ball around $(\mathbf{x}(p),0)$ — the differential being the solution
  of the variational equation along the flow line through $\zeta$, no longer
  the explicit shear of \cref{lem:dc-ch3-3-7-pairmap} away from the
  equilibrium. Composing with the linear rescaling, the evaluation at $T$ and
  the projection, $G$ is strictly differentiable at every point of the stated
  open set, and a map strictly differentiable at every point of an open set is
  $C^1$ there (as in \cref{lem:dc-ch3-2-9-c1}). The slice statement follows by
  composing with the affine embedding $w\mapsto(y,w)$, which maps the ball
  $|w|<Tr$ into the admissible set.
\end{proof}

\begin{lemma}[the geodesic segment determines the pair map]
  \label{lem:dc-ch3-3-7-flow-agree}
  \lean{Riemannian.Exponential.uniform_flow_pairMap_agree}
  \uses{lem:dc-ch3-3-7-descent}
  \leanok
  \dcref{ch3:3.7}
  Let $Z$ and $Z'$ be two local flows of the chart-$p$ geodesic system as in
  \cref{lem:dc-ch3-3-7-descent}, with time scales $0<T<\varepsilon$ and
  $0<T'<\varepsilon'$ and admissibility balls of radii $r$, $r'$ around
  $(\mathbf{x}(p),0)$. For any initial data $(y,w)$ with $(y,w/T)$ admissible
  for $Z$ and $(y,w/T')$ admissible for $Z'$,
  $$Z(y,w/T)(T)_1=Z'(y,w/T')(T')_1:$$
  the chart reading of the time-$1$ endpoint of the geodesic from
  $\mathbf{x}^{-1}(y)$ with chart velocity $w$ does not depend on the flow
  used to compute it. In particular the pair maps
  $G(y,w)=(y,\,Z(y,w/T)(T)_1)$ built from any two such flows agree on the
  common admissible set.
\end{lemma}
\begin{proof}
  \sketch
  \leanok
  \uses{lem:dc-ch3-3-7-descent, lem:dc-ch7-2-8-uniqueness}
  By \cref{lem:dc-ch3-3-7-descent} — whose proof gives the geodesic property
  on the open window $|s|<\varepsilon/T\supset[0,1]$, and likewise
  $|s|<\varepsilon'/T'$ for $Z'$ — the two rescaled segments
  $s\mapsto\mathbf{x}^{-1}(Z(y,w/T)(sT)_1)$ and
  $s\mapsto\mathbf{x}^{-1}(Z'(y,w/T')(sT')_1)$ are continuous geodesics
  passing through $\mathbf{x}^{-1}(y)$ at $s=0$ with the same chart velocity
  $w$. By uniqueness of geodesics with given initial data
  (\cref{lem:dc-ch7-2-8-uniqueness}) they agree on the common open window,
  which contains $s=1$ since $T<\varepsilon$ and $T'<\varepsilon'$. Both
  endpoint readings lie in the chart image, where $\mathbf{x}^{-1}$ is
  injective, so evaluating the identity at $s=1$ equates them.
\end{proof}

\begin{lemma}[the derivative of the pair map is invertible near the zero section]
  \label{lem:dc-ch3-3-7-invertible}
  \lean{Riemannian.Exponential.exists_pairMap_hasStrictFDerivAt_equiv_ball}
  \uses{lem:dc-ch3-3-7-pairmap, lem:dc-ch3-3-7-c1pair}
  \leanok
  \dcref{ch3:3.7}
  There are a flow package $(r,\varepsilon,T,Z)$ as in
  \cref{lem:dc-ch3-3-7-descent} and a radius $\rho>0$ such that the pair map
  $G(y,w)=(y,\,Z(y,w/T)(T)_1)$ fixes the center
  ($G(\mathbf{x}(p),0)=(\mathbf{x}(p),\mathbf{x}(p))$), is $C^1$ on the open
  set of admissible initial conditions, is differentiable at the center with
  derivative the shear $(a,b)\mapsto(a,a+b)$ of
  \cref{lem:dc-ch3-3-7-pairmap}, and has an \emph{invertible} differential at
  every point of the ball $B((\mathbf{x}(p),0),\rho)$.
\end{lemma}
\begin{proof}
  \sketch
  \leanok
  \uses{lem:dc-ch3-3-7-pairmap, lem:dc-ch3-3-7-c1pair, lem:dc-ch3-3-7-flow-agree}
  Take the flow package of \cref{lem:dc-ch3-3-7-c1pair}, so that $G$ is $C^1$
  on the open admissible set and its differential $dG$ is continuous there.
  The derivative at the center is the shear: \cref{lem:dc-ch3-3-7-pairmap}
  computes it for \emph{its} flow package, and by
  \cref{lem:dc-ch3-3-7-flow-agree} the two pair maps agree on a neighborhood
  of $(\mathbf{x}(p),0)$, so the derivative transports; the center is fixed
  because the zero-velocity flow line is constant. The shear $S(a,b)=(a,a+b)$
  is invertible with inverse $(a,c)\mapsto(a,c-a)$, and invertibility is
  stable under small perturbation: choose $\rho>0$ so that
  $\|dG_x-S\|<(\|S^{-1}\|+1)^{-1}$ for $|x-(\mathbf{x}(p),0)|<\rho$; then
  $t:=S^{-1}\circ(S-dG_x)$ has $\|t\|\le\|S^{-1}\|\,\|S-dG_x\|<1$, the
  Neumann series inverts $1-t$, and $dG_x=S\circ(1-t)$ is invertible — as in
  \cref{lem:dc-ch3-2-9-invertible}.
\end{proof}

\begin{theorem}[totally normal neighborhood]
  \label{thm:dc-ch3-3-7}
  \lean{Riemannian.Exponential.exists_totallyNormal_c1_diffeo, Riemannian.Exponential.exists_totallyNormal_neighborhood}
  \uses{prop:dc-ch3-2-7, lem:dc-ch3-3-7-descent, lem:dc-ch3-3-7-descent-velocity, lem:dc-ch3-3-7-pairmap, lem:dc-ch3-3-7-c1pair}
  \leanok
  \dcref{ch3:3.7}
  For any $p\in M$ there exist a neighborhood $W$ of $p$ and a number $\delta>0$
  such that for every $q\in W$, $\exp_q$ is a diffeomorphism on
  $B_\delta(0)\subset T_qM$ and $\exp_q(B_\delta(0))\supset W$; that is, $W$ is a
  normal neighborhood of each of its points.
\end{theorem}
\begin{proof}
  \sketch
  \leanok
  \uses{lem:dc-ch3-3-7-descent, lem:dc-ch3-3-7-descent-velocity, lem:dc-ch3-3-7-pairmap, lem:dc-ch3-3-7-c1pair, lem:dc-ch3-3-7-invertible}
  Let $\varepsilon$, $V$, $\mathcal{U}$ be as in \cref{prop:dc-ch3-2-7}, with
  $\mathcal{U}\subset TU$, $V\subset\mathbf{x}(U)$. Define $F:\mathcal{U}\to M\times M$
  by $F(q,v)=(q,\exp_q v)$. Around $F(p,0)=(p,p)$ take coordinates
  $(U\times U;\mathbf{x},\mathbf{x})$; then since $(d\exp_p)_0=I$,

  $$
  dF_{(p,0)}=\begin{pmatrix}I & I\\ 0 & I\end{pmatrix},
  $$

  so $F$ is a local diffeomorphism near $(p,0)$, mapping a neighborhood
  $\mathcal{U}'=\{(q,v);\ q\in V',\ v\in T_qM,\ |v|<\delta\}$ diffeomorphically onto
  a neighborhood $W'$ of $(p,p)$. Choose $W\ni p$ with $W\times W\subset W'$. For
  $q\in W$ and $B_\delta(0)\subset T_qM$, $F(\{q\}\times B_\delta(0))\supset\{q\}\times W$,
  so by definition of $F$, $\exp_q(B_\delta(0))\supset W$.
\end{proof}

\begin{remark}[uniqueness and differentiable dependence]
  \label{rem:dc-ch3-3-8}
  \uses{thm:dc-ch3-3-7}
  \dcref{ch3:3.8}
  From \cref{thm:dc-ch3-3-7} and the minimizing property of geodesics, any two points
  $q_1,q_2\in W$ are joined by a unique minimizing geodesic $\gamma$ of length
  $<\delta$; moreover $\gamma$ depends differentiably on $(q_1,q_2)$: there is a
  unique $v\in T_{q_1}M$ with $F^{-1}(q_1,q_2)=(q_1,v)$, depending differentiably on
  $(q_1,q_2)$, with $\gamma'(0)=v$. Such a $W$ is a \emph{totally normal neighborhood} of
  $p\in M$.
\end{remark}

\begin{lemma}[a minimizing broken pair of radial geodesics has no corner]
  \label{lem:dc-ch3-3-9-corner}
  \lean{Riemannian.Exponential.eq_neg_of_forall_edist_expMap_eq}
  \uses{def:dc-ch3-expmap, def:dc-ch7-2-4, lem:dc-ch3-2-9-c1ball, lem:dc-ch7-2-8-gram-bound, lem:dc-ch7-2-8-chart-length}
  \leanok
  \dcref{ch3:3.9}
  Let $M$ be a Riemannian manifold whose distance $d$ is the Riemannian
  distance, $x\in M$, and let $u_1,u_2\in T_xM$ be $g$-unit vectors. If the
  two radial legs through $x$ jointly realize the distance between their
  endpoints for every small radius — there is $\eta_0>0$ with
  $d\big(\exp_x(\eta u_1),\exp_x(\eta u_2)\big)=2\eta$ for all
  $0<\eta<\eta_0$ — then $u_2=-u_1$: the broken curve
  $\exp_x(\eta u_1)\rightsquigarrow x\rightsquigarrow\exp_x(\eta u_2)$ has no
  corner at $x$.
\end{lemma}
\begin{proof}
  \sketch
  \leanok
  The engine is the \emph{chord upper bound}: for every $\theta>1$ there
  is $\rho>0$ with $d(\exp_x v,\exp_x w)\le\theta\,|w-v|_x$ for
  $\|v\|,\|w\|<\rho$. Its witness is the exponential image
  $s\mapsto\exp_x\big(v+s\,(w-v)\big)$ of the straight segment: since
  $\exp_x$ is $C^1$ near $0$ with $d(\exp_x)_0=\operatorname{id}$
  (\cref{lem:dc-ch3-2-9-c1ball}) and the Gram form is jointly continuous
  and positive definite at the pole (\cref{lem:dc-ch7-2-8-gram-bound}), the
  $g$-speed of this curve is at most $\theta\,|w-v|_x$ on a small enough
  ball, uniformly in the direction; integrating
  (\cref{lem:dc-ch7-2-8-chart-length}) bounds its length, and the distance
  is at most the length of any such competitor (\cref{def:dc-ch7-2-4}).
  Applying the bound to $v=\eta u_1$, $w=\eta u_2$ for small $\eta$ gives
  $2\eta\le\theta\,\eta\,|u_2-u_1|_x$, i.e. $2\le\theta\,|u_2-u_1|_x$ for
  every $\theta>1$, hence $|u_2-u_1|_x\ge2$. Expanding with
  $|u_1|_x=|u_2|_x=1$:
  $4\le|u_2-u_1|_x^2=2-2\langle u_1,u_2\rangle_x$, so
  $\langle u_1,u_2\rangle_x\le-1$, whence
  $|u_1+u_2|_x^2=2+2\langle u_1,u_2\rangle_x\le0$; positive definiteness of
  $g_x$ gives $u_1+u_2=0$.
\end{proof}

\begin{lemma}[anchored radial identification of a distance-realizing segment]
  \label{lem:dc-ch3-3-9-anchored}
  \lean{Riemannian.Exponential.exists_radial_of_anchored}
  \uses{lem:dc-ch3-3-6-equality-geodesic, lem:dc-ch7-2-8-normal-ball-dist, lem:dc-ch7-2-8-gram-bound}
  \leanok
  \dcref{ch3:3.9}
  Let $\gamma$ be $C^1$ on $[t_0,B]$, $t_0<B$, with speed $\ell>0$ in the
  sense that $\ell(\gamma|_{[t_0,t]})=\ell\,(t-t_0)$ and
  $d\big(\gamma(t_0),\gamma(t)\big)=\ell\,(t-t_0)$ for all $t\in[t_0,B]$.
  Then, with $q=\gamma(t_0)$, there are $t_1\in(t_0,B]$ and $v\in T_qM$ with
  $|v|_q=\ell\,(t_1-t_0)$ such that
  $$\gamma(r)=\exp_q\Big(\tfrac{r-t_0}{t_1-t_0}\,v\Big)
    \qquad\text{for all } r\in[t_0,t_1]:$$
  some initial subsegment of $\gamma$ is exactly the radial geodesic of the
  normal ball at its own starting point.
\end{lemma}
\begin{proof}
  \sketch
  \leanok
  Choose $t_1\in(t_0,B]$ with $\ell\,(t_1-t_0)$ smaller than both the
  metric-ball threshold $\delta$ of the normal ball at $q$
  (\cref{lem:dc-ch7-2-8-normal-ball-dist}) and, after conversion by the Gram
  comparison $\|w\|^2\le c\,\langle w,w\rangle_q$
  (\cref{lem:dc-ch7-2-8-gram-bound}, as packaged in
  \cref{lem:dc-ch7-2-8-normal-ball-dist}), the Gauss radius of the equality
  case at $q$. Since $d\big(q,\gamma(t_1)\big)=\ell\,(t_1-t_0)<\delta$, the point
  $\gamma(t_1)$ lies in the normal ball: $\gamma(t_1)=\exp_q v$ with
  $|v|_q=d\big(q,\gamma(t_1)\big)=\ell\,(t_1-t_0)$ by the distance formula
  of \cref{lem:dc-ch7-2-8-normal-ball-dist}, and the Gram bound puts the
  coordinate norm of $v$ inside the Gauss radius. The affine
  reparametrization $\sigma(s)=\gamma\big(t_0+s\,(t_1-t_0)\big)$ is $C^1$ on
  $[0,1]$, runs from $q$ to $\exp_q v$, and has running length
  $\ell(\sigma|_{[0,s]})=\ell\cdot s\,(t_1-t_0)=s\,|v|_q$ — exactly the
  arclength-proportional hypothesis of the equality case
  \cref{lem:dc-ch3-3-6-equality-geodesic}, which yields
  $\sigma(s)=\exp_q(s\,v)$ on $[0,1]$; undoing the reparametrization gives
  the claim.
\end{proof}

\begin{lemma}[localization at moving centers: the two-sided corner argument]
  \label{lem:dc-ch3-3-9-localization}
  \lean{Riemannian.Exponential.hasGeodesicEquationAt_of_arclength_edist}
  \uses{lem:dc-ch3-3-9-anchored, lem:dc-ch3-3-9-corner, lem:dc-ch7-2-8-ray-geodesic, lem:dc-ch3-2-6-reparam, def:dc-ch3-2-1}
  \leanok
  \dcref{ch3:3.9}
  Let $\gamma$ be parametrized proportionally to arc length with speed
  $\ell>0$ and distance-realizing on $[A,B]$ — i.e.
  $\ell(\gamma|_{[s,t]})=d\big(\gamma(s),\gamma(t)\big)=\ell\,(t-s)$ for all
  $A\le s\le t\le B$ — and let $A<u<B$ be an interior time such that
  $\gamma$ is $C^1$ on $[A,u]$ and on $[u,B]$ separately ($u$ may be a corner
  of a piecewise curve). Then $\gamma$ satisfies the geodesic equation of
  \cref{def:dc-ch3-2-1} at $u$.
\end{lemma}
\begin{proof}
  \sketch
  \leanok
  Anchor \cref{lem:dc-ch3-3-9-anchored} at $q=\gamma(u)$ twice: applied to
  $\gamma|_{[u,B]}$ it exhibits a forward subsegment
  $\gamma(r)=\exp_q\big(\tfrac{r-u}{t_1-u}v_2\big)$ on $[u,t_1]$; applied to
  the time reversal $s\mapsto\gamma(2u-s)$ — which is again
  arclength-proportional and distance-realizing, with the same speed — it
  exhibits a backward subsegment
  $\gamma(r)=\exp_q\big(\tfrac{u-r}{u-t_0'}v_1\big)$ on $[t_0',u]$. Let
  $u_i=v_i/|v_i|_q$. For $0<\eta<\ell\min(u-t_0',t_1-u)$ the points
  $\exp_q(\eta u_1)=\gamma(u-\eta/\ell)$ and
  $\exp_q(\eta u_2)=\gamma(u+\eta/\ell)$ realize the distance
  $d=\ell\cdot(2\eta/\ell)=2\eta$, so \cref{lem:dc-ch3-3-9-corner} gives
  $u_2=-u_1$. Substituting back, both subsegments trace the \emph{single}
  ray $\gamma(r)=\exp_q\big(\ell\,(r-u)\,u_2\big)$ for
  $r\in[t_0',t_1]$. Exponential rays are geodesics on an open interval
  around time $0$ (\cref{lem:dc-ch7-2-8-ray-geodesic}), affine
  reparametrization preserves the geodesic equation
  (\cref{lem:dc-ch3-2-6-reparam}), and the geodesic equation at $u$ only
  depends on the curve near $u$, where $\gamma$ coincides with this ray.

\end{proof}

\begin{corollary}[minimizing curve is a geodesic]
  \label{cor:dc-ch3-3-9}
  \lean{Riemannian.Exponential.isGeodesicOn_piecewise_of_arclength_forall_le, Riemannian.Exponential.isGeodesicOn_piecewise_of_arclength_edist, Riemannian.Exponential.isGeodesicOn_of_arclength_edist, Riemannian.Exponential.edist_le_pathELength_piecewise_partition}
  \uses{prop:dc-ch3-3-6, lem:dc-ch3-3-6-equality-geodesic, lem:dc-ch3-3-9-localization, lem:dc-ch3-3-9-corner, lem:dc-ch3-3-9-anchored, lem:dc-ch7-2-8-chart-length, def:dc-ch3-3-1, def:dc-ch3-2-1}
  \leanok
  \dcref{ch3:3.9}
  If a piecewise differentiable curve $\gamma:[a,b]\to M$, parametrized
  proportionally to arc length, has length $\le$ that of any other piecewise
  differentiable curve joining $\gamma(a)$ to $\gamma(b)$, then $\gamma$ is a
  geodesic; in particular it is regular.
\end{corollary}
\begin{proof}
  \sketch
  \leanok
  For $t\in[a,b]$ let $W$ be a totally normal neighborhood of $\gamma(t)$
  and $I\subset[a,b]$ a closed interval with $t\in I$ and $\gamma(I)\subset W$.
  Then $\gamma_I$ joins two points of a normal ball; by the hypothesis and
  \cref{prop:dc-ch3-3-6}, $\ell(\gamma_I)$ equals the length of the radial geodesic joining
  them, so $\gamma_I$ coincides with that geodesic. As $\gamma_I$ is parametrized
  proportionally to arc length it is a geodesic on $I$, hence at $t$.
\end{proof}

\begin{example}[the Lobatchevski plane]
  \label{ex:dc-ch3-3-10}
  \uses{cor:dc-ch3-3-9}
  \notready
  \dcref{ch3:3.10}
  \notready
  The isometries of a Riemannian manifold take geodesics into geodesics.
  Let $G=\{(x,y)\in\mathbb{R}^2;\ y>0\}$ with metric $g_{11}=g_{22}=\frac{1}{y^2}$,
  $g_{12}=g_{21}=0$. The $y$-axis segment $\gamma(t)=(0,t)$, $a\le t\le b$ ($a>0$),
  is the image of a geodesic: for any arc $c(t)=(x(t),y(t))$ with $c(a)=(0,a)$,
  $c(b)=(0,b)$,
  $$
  \ell(c)=\int_a^b\sqrt{(x')^2+(y')^2}\,\frac{dt}{y}\ge\int_a^b\frac{|y'|}{y}\,dt\ge\int_a^b\frac{dy}{y}=\ell(\gamma),
  $$
  so $\gamma$ minimizes, and by \cref{cor:dc-ch3-3-9} its image is a geodesic. The
  isometries $z\mapsto\frac{az+b}{cz+d}$, $z=x+iy$, $ad-bc=1$ carry the $y$-axis
  into upper semicircles or rays $x=x_0$, $y>0$, which
  are therefore geodesics; these are all the geodesics, since through each $p\in G$
  and each direction there passes such a circle centered on the $x$-axis (a normal
  direction giving a vertical ray).
\end{example}

\section{Convex neighborhoods}

By \cref{thm:dc-ch3-3-7} (cf. \cref{rem:dc-ch3-3-8}) every $p\in M$ has a totally normal neighborhood
$W$ with $\delta>0$ so that any two $q_1,q_2\in W$ are joined by a minimizing
geodesic of length $<\delta$; but such a geodesic need not lie entirely in $W$. A
subset $S\subset M$ is \emph{strongly convex} if for any two points $q_1,q_2$ in the
closure $\overline{S}$ there is a unique minimizing geodesic $\gamma$ joining them
whose interior lies in $S$. We show the radius of a totally normal ball can be
chosen so the ball is strongly convex.

The analytic heart of \cref{lem:dc-ch3-4-1} is the behaviour of the squared radial
distance $F(t)=|u(t)|_p^2$, $u(t)=\exp_p^{-1}(\gamma(t))$, along a geodesic $\gamma$
read in normal coordinates at $p$. We isolate its fibrewise (fixed base point $q$
and velocity $v$) calculus in the following lemmas; the metric
$\langle\cdot,\cdot\rangle_p$ is the fixed chart Gram inner product at $p$.

\begin{lemma}[Leibniz rule for the fixed metric at $p$]
  \label{lem:dc-ch3-4-1-leibniz}
  \lean{Riemannian.hasDerivAt_chartMetricInner}
  \leanok
  \dcref{ch3:4.1}
  For curves $f,h:\mathbb{R}\to T_pM$ differentiable at $t$, read in the chart at
  $p$,
  $$\frac{d}{dt}\langle f(t),h(t)\rangle_p
    = \langle f'(t),h(t)\rangle_p + \langle f(t),h'(t)\rangle_p.$$
\end{lemma}
\begin{proof}
  \sketch
  \leanok
  The form is the finite sum $\sum_{i,j}G_{ij}\,f^i(t)\,h^j(t)$ with constant
  coefficients $G_{ij}$ (the base point is fixed); each coordinate $f^i,h^j$ is a
  continuous linear functional of the curve, so the product rule applies termwise.
\end{proof}

\begin{lemma}[first time-derivative of the squared radial distance]
  \label{lem:dc-ch3-4-1-dF}
  \lean{Riemannian.hasDerivAt_chartMetricInner_diag}
  \uses{lem:dc-ch3-4-1-leibniz}
  \leanok
  \dcref{ch3:4.1}
  For $u:\mathbb{R}\to T_pM$ differentiable at $t$, the squared radial distance
  $F(t)=\langle u(t),u(t)\rangle_p$ satisfies
  $$\frac{\partial F}{\partial t}=2\Big\langle\frac{\partial u}{\partial t},u\Big\rangle_p.$$
\end{lemma}
\begin{proof}
  \sketch
  \leanok
  Apply \cref{lem:dc-ch3-4-1-leibniz} with $f=h=u$ and use the symmetry of
  $\langle\cdot,\cdot\rangle_p$.
\end{proof}

\begin{lemma}[second time-derivative of the squared radial distance]
  \label{lem:dc-ch3-4-1-ddF}
  \lean{Riemannian.hasDerivAt_deriv_chartMetricInner_diag}
  \uses{lem:dc-ch3-4-1-leibniz}
  \leanok
  \dcref{ch3:4.1}
  If $u$ is twice differentiable at $t$, then the first-derivative field
  $s\mapsto 2\langle u'(s),u(s)\rangle_p$ is differentiable at $t$ with
  $$\frac{\partial^2 F}{\partial t^2}
    = 2\Big\langle\frac{\partial^2 u}{\partial t^2},u\Big\rangle_p
      + 2\Big|\frac{\partial u}{\partial t}\Big|_p^2,$$
  the displayed second time-derivative of $F=|u|_p^2$.
\end{lemma}
\begin{proof}
  \sketch
  \leanok
  The field $s\mapsto 2\langle u'(s),u(s)\rangle_p$ is twice the bilinear pairing
  of $u'$ and $u$; differentiate it by \cref{lem:dc-ch3-4-1-leibniz} with
  $f=u'$, $h=u$.
\end{proof}

\begin{lemma}[normal coordinate along the radial geodesic]
  \label{lem:dc-ch3-4-1-radialcoord}
  \lean{Riemannian.Exponential.exists_expMap_inv_smul_eq}
  \uses{def:dc-ch3-expmap, lem:dc-ch3-2-9-c1diffeo}
  \leanok
  \dcref{ch3:4.1}
  There is $\varepsilon>0$ and a local inverse $\exp_p^{-1}$ of $\exp_p$ on
  $B_\varepsilon(0)$ such that for every $v$ and every $t$ with
  $\|tv\|<\varepsilon$, $\exp_p^{-1}(\exp_p(tv))=tv$. Hence along the radial
  geodesic $t\mapsto\exp_p(tv)$ the normal coordinate $u(t)=\exp_p^{-1}(\gamma(t,p,v))$
  is $tv$ itself, and $F(t)=|u(t)|_p^2=t^2|v|_p^2$.
\end{lemma}
\begin{proof}
  \sketch
  \leanok
  This is the left-inverse clause of the $C^1$ diffeomorphism
  \cref{lem:dc-ch3-2-9-c1diffeo}, instantiated at $w=tv$.
\end{proof}

\begin{lemma}[center Hessian: $\partial^2 F/\partial t^2(0,p,v)=2|v|^2$]
  \label{lem:dc-ch3-4-1-hessian-center}
  \lean{Riemannian.Exponential.hasDerivAt_secondDeriv_expMap_inv_sqNorm_radial}
  \uses{lem:dc-ch3-4-1-ddF, lem:dc-ch3-4-1-radialcoord}
  \leanok
  \dcref{ch3:4.1}
  For $q=p$ the geodesic through $p$ with velocity $v$ is the radial ray, so
  $u(t)=tv$ (\cref{lem:dc-ch3-4-1-radialcoord}), $u'=v$, $u''=0$, and by
  \cref{lem:dc-ch3-4-1-ddF}
  $$\frac{\partial^2 F}{\partial t^2}(0,p,v)=2|v|_p^2.$$
  In particular for a unit $v$ this equals $2>0$, the strict-minimum seed.
\end{lemma}
\begin{proof}
  \sketch
  \leanok
  With $u(t)=tv$ the first-derivative field is
  $s\mapsto 2\langle v,sv\rangle_p=2s\,|v|_p^2$, whose derivative at $0$ is
  $2|v|_p^2$.
\end{proof}

\begin{lemma}[the metric at $p$ is positive definite]
  \label{lem:dc-ch3-4-1-posdef}
  \lean{Riemannian.Exponential.chartMetricInner_extChartAt_self_pos}
  \leanok
  \dcref{ch3:4.1}
  For $v\neq0$, $|v|_p^2=\langle v,v\rangle_p>0$. Hence the center Hessian
  $\frac{\partial^2 F}{\partial t^2}(0,p,v)=2|v|_p^2$
  (\cref{lem:dc-ch3-4-1-hessian-center}) is strictly positive, giving the
  strict minimum of $F$ at $q=p$.
\end{lemma}
\begin{proof}
  \sketch
  \leanok
  Positive-definiteness of the Riemannian metric $g$ at $p$, read in the chart
  at $p$.
\end{proof}

\begin{lemma}[a tangent geodesic with positive Hessian stays outside]
  \label{lem:dc-ch3-4-1-outside}
  \lean{Riemannian.eventually_ge_of_deriv_deriv_pos}
  \leanok
  \dcref{ch3:4.1}
  If $F:\mathbb{R}\to\mathbb{R}$ is continuous at $0$ with $F'(0)=0$ and
  $F''(0)>0$, then $F(t)\ge F(0)$ for all $t$ near $0$. With $F=|u|_p^2$ and
  $F(0)=r^2$, a geodesic tangent to $S_r(p)$ at $\gamma(0)$ keeps
  $|\exp_p^{-1}\gamma(t)|^2\ge r^2$, i.e. does not re-enter the open ball
  $B_r(p)$.
\end{lemma}
\begin{proof}
  \sketch
  \leanok
  The second-derivative test: $F'(0)=0$ and $F''(0)>0$ give a local minimum of
  $F$ at $0$, so $F(t)\ge F(0)$ near $0$.
\end{proof}

\begin{lemma}[a tangent geodesic with positive Hessian stays \emph{strictly} outside]
  \label{lem:dc-ch3-4-1-strict-outside}
  \lean{Riemannian.eventually_lt_of_deriv_deriv_pos}
  \leanok
  \dcref{ch3:4.1}
  If $F:\mathbb{R}\to\mathbb{R}$ is continuous at $0$ with $F'(0)=0$ and
  $F''(0)>0$, then $F(0)<F(t)$ for \emph{every} $t\neq0$ near $0$ (a strict
  punctured minimum). With $F=|u|_p^2$ and $F(0)=r^2$, a geodesic tangent to
  $S_r(p)$ at $\gamma(0)$ keeps $|\exp_p^{-1}\gamma(t)|^2>r^2$ for $t\neq0$, i.e.
  stays \emph{strictly} outside the closed ball $\overline{B_r(p)}$ except at the
  point of tangency. This strengthens \cref{lem:dc-ch3-4-1-outside} and is the
  separation the convex-neighborhood argument \cref{prop:dc-ch3-4-2}
  consumes: it rules out the interior distance-maximum whose existence the
  proposition assumes for contradiction.
\end{lemma}
\begin{proof}
  \sketch
  \leanok
  By the sign form of the second-derivative test, $\operatorname{sign} F'(t)$
  agrees with $\operatorname{sign}(t)$ near $0$ (from $F'(0)=0$, $F''(0)>0$): $F'>0$
  on a right interval $(0,c)$ and $F'<0$ on a left interval $(a,0)$. Hence $F$ is
  strictly increasing on $[0,c)$ and strictly decreasing on $(a,0]$, so $F(0)<F(t)$
  on both punctured sides.
\end{proof}

\begin{lemma}[joint continuity of the second time-derivative in the base point]
  \label{lem:dc-ch3-4-1-continuity}
  \lean{Riemannian.Exponential.continuousOn_secondDerivChartForm}
  \uses{lem:dc-ch3-4-1-ddF, lem:dc-ch3-2-9-c2diffeo}
  \leanok
  \dcref{ch3:4.1}
  Write $G(q,v)=\frac{\partial^2 F}{\partial t^2}(0,q,v)$. Read in normal coordinates
  at $p$, with $a=\varphi_p(q)$ the chart position and $w$ the chart velocity, this
  value is the algebraic expression
  $$G(q,v)=2\big\langle D^2(\exp_p^{-1})_a(w,w)+D(\exp_p^{-1})_a\!\big(-\Gamma_a(w,w)\big),\ \exp_p^{-1}(q)\big\rangle_p
    +2\big|D(\exp_p^{-1})_a(w)\big|_p^2,$$
  where $-\Gamma_a(w,w)$ is the geodesic acceleration $x''(0)$ read in the chart at
  $p$ (from the geodesic equation $x''+\Gamma(x',x')=0$). Since $\exp_p^{-1}$ is $C^2$
  and the Christoffel symbols are smooth, $G$ is jointly continuous in $(q,v)$.
\end{lemma}
\begin{proof}
  \sketch
  \leanok
  The value $G(q,v)$ needs no time integration: the geodesic acceleration at $t=0$
  is $x''(0)=-\Gamma_a(w,w)$ directly from the geodesic ODE, so by the second-order
  chain rule $u''(0)=D^2(\exp_p^{-1})_a(w,w)+D(\exp_p^{-1})_a(x''(0))$. The three maps
  $\exp_p^{-1}$, $D(\exp_p^{-1})$, $D^2(\exp_p^{-1})$ are continuous on the normal chart
  image ($\exp_p^{-1}$ is $C^2$, \cref{lem:dc-ch3-2-9-c2diffeo}), $x''(0)=-\Gamma_a(w,w)$
  is continuous in $(a,w)$ (the geodesic spray is smooth), and the chart Gram form
  $\langle\cdot,\cdot\rangle_p$ is a continuous bilinear form; $G$ is a finite algebraic
  combination of these, hence continuous.
\end{proof}

\begin{lemma}[strict positivity of the second time-derivative at $q=p$]
  \label{lem:dc-ch3-4-1-poscenter}
  \lean{Riemannian.Exponential.secondDerivChartForm_extChartAt_self_pos}
  \uses{lem:dc-ch3-4-1-posdef, lem:dc-ch3-2-9-c2diffeo}
  \leanok
  \dcref{ch3:4.1}
  At $q=p$ one has $\exp_p^{-1}(p)=0$, so the $\langle u'',u\rangle_p$ term of $G(p,v)$
  vanishes and $G(p,v)=2\big|D(\exp_p^{-1})_0(w)\big|_p^2$. Since $D(\exp_p^{-1})_0$ is a
  linear isomorphism and $\langle\cdot,\cdot\rangle_p$ is positive definite, $G(p,v)>0$
  for $v\neq0$.
\end{lemma}
\begin{proof}
  \sketch
  \leanok
  $u(0,p,v)=\exp_p^{-1}(p)=0$ kills the first term. The derivative $D(\exp_p^{-1})_0$ is
  the inverse of $D(\varphi_p\circ\exp_p)_0$ (inverse function theorem), hence injective,
  so $D(\exp_p^{-1})_0(w)\neq0$ for $w\neq0$; then
  $\big|D(\exp_p^{-1})_0(w)\big|_p^2>0$ by positive definiteness of the metric at $p$
  (\cref{lem:dc-ch3-4-1-posdef}).
\end{proof}

\begin{lemma}[a neighborhood of $p$ with strictly positive Hessian]
  \label{lem:dc-ch3-4-1-posnbhd}
  \lean{Riemannian.Exponential.exists_secondDerivChartForm_pos_nhds}
  \uses{lem:dc-ch3-4-1-continuity, lem:dc-ch3-4-1-poscenter}
  \leanok
  \dcref{ch3:4.1}
  There is a neighborhood $V$ of $p$ (inside a normal chart) such that
  $\frac{\partial^2 F}{\partial t^2}(0,q,v)>0$ for every $q\in V$ and every unit chart
  velocity $v$.
\end{lemma}
\begin{proof}
  \sketch
  \leanok
  $G(q,v)$ is jointly continuous (\cref{lem:dc-ch3-4-1-continuity}) and $G(p,v)>0$ for
  $v\neq0$ (\cref{lem:dc-ch3-4-1-poscenter}). Since the unit sphere of chart velocities
  is compact, the tube lemma produces an open neighborhood $V$ of $p$ on which
  $G(q,v)>0$ for all $q\in V$ and all unit $v$.
\end{proof}

\begin{lemma}[degree-2 homogeneity: strict positivity for every nonzero velocity]
  \label{lem:dc-ch3-4-1-homog}
  \lean{Riemannian.Exponential.exists_secondDerivChartForm_pos_nhds_ne}
  \uses{lem:dc-ch3-4-1-posnbhd, lem:dc-ch3-2-6}
  \leanok
  \dcref{ch3:4.1}
  The value $G(q,v)=\frac{\partial^2 F}{\partial t^2}(0,q,v)$ is degree-2 homogeneous in
  the chart velocity: $G(q,cv)=c^2\,G(q,v)$. Consequently the neighborhood $V$ of
  \cref{lem:dc-ch3-4-1-posnbhd} satisfies $G(q,v)>0$ for every $q\in V$ and every
  \emph{nonzero} chart velocity $v$, removing the unit-speed restriction.
\end{lemma}
\begin{proof}
  \sketch
  \leanok
  Under $v\mapsto cv$ every building block of $G$ scales:
  $D(\exp_p^{-1})_a(cv)=c\,D(\exp_p^{-1})_a(v)$,
  $D^2(\exp_p^{-1})_a(cv,cv)=c^2\,D^2(\exp_p^{-1})_a(v,v)$, and the geodesic acceleration
  $-\Gamma_a(cv,cv)=c^2\,\big(-\Gamma_a(v,v)\big)$ by the degree-2 homogeneity of the
  Christoffel contraction (the coordinate form of the homogeneity \cref{lem:dc-ch3-2-6});
  since $\langle\cdot,\cdot\rangle_p$ is bilinear both summands acquire the factor $c^2$.
  Writing a nonzero $v=|v|\,(v/|v|)$ with $v/|v|$ of unit chart norm,
  $G(q,v)=|v|^2\,G(q,v/|v|)>0$ by \cref{lem:dc-ch3-4-1-posnbhd}.
\end{proof}

\begin{lemma}[the F-identity: second-order chain rule for the squared radial distance]
  \label{lem:dc-ch3-4-1-fidentity}
  \lean{Riemannian.Exponential.hasDerivAt_and_deriv_deriv_sqNormComp}
  \uses{lem:dc-ch3-4-1-dF, lem:dc-ch3-4-1-ddF}
  \leanok
  \dcref{ch3:4.1}
  Let $\mathrm{finv}$ be $C^2$ on an open set $S\ni a$, and let $x:\mathbb R\to T_pM$ be a
  curve differentiable near $0$ with values eventually in $S$, and with $x(0)=a$, $x'(0)=w$,
  second derivative $x''(0)=\xi$. Then, writing $u=\mathrm{finv}\circ x$, the squared radial
  distance $F(s)=\langle u(s),u(s)\rangle_p$ satisfies
  $$\frac{\partial F}{\partial t}(0)=2\langle D\mathrm{finv}(a)(w),\mathrm{finv}(a)\rangle_p,$$
  $$\frac{\partial^2 F}{\partial t^2}(0)
    =2\big\langle D^2\mathrm{finv}(a)(w,w)+D\mathrm{finv}(a)(\xi),\ \mathrm{finv}(a)\big\rangle_p
      +2\,\big|D\mathrm{finv}(a)(w)\big|_p^2,$$
  the displayed derivatives of $F=|u|_p^2$ follow from the second-order chain rule
  $u''(0)=D^2\mathrm{finv}(a)(w,w)+D\mathrm{finv}(a)(x''(0))$ made explicit.
\end{lemma}
\begin{proof}
  \sketch
  \leanok
  \uses{lem:dc-ch3-4-1-leibniz, lem:dc-ch3-4-1-dF, lem:dc-ch3-4-1-ddF}
  Near $0$, $u$ is differentiable with $u'(s)=D\mathrm{finv}(x(s))\big(x'(s)\big)$ (chain rule:
  $\mathrm{finv}$ is differentiable on the open $S$ and $x(s)\in S$ for $s$ near $0$), so
  $\partial F/\partial t(s)=2\langle u'(s),u(s)\rangle_p$ (\cref{lem:dc-ch3-4-1-dF}); evaluating
  at $0$ gives the first identity. Differentiating the field $s\mapsto D\mathrm{finv}(x(s))$
  against $x'$ at $0$ by the composition and bilinear-application rules gives
  $u''(0)=D^2\mathrm{finv}(a)(w,w)+D\mathrm{finv}(a)(x''(0))$, and \cref{lem:dc-ch3-4-1-ddF}
  turns $\partial^2 F/\partial t^2(0)=2\langle u''(0),u(0)\rangle_p+2|u'(0)|_p^2$ into the
  second identity.
\end{proof}

\begin{lemma}[the F-identity along the moving-base geodesic flow]
  \label{lem:dc-ch3-4-1-flowreading}
  \lean{Riemannian.Exponential.hasDerivAt_and_deriv_deriv_sqNorm_flowReading}
  \uses{lem:dc-ch3-4-1-fidentity, lem:dc-ch3-3-7-descent, lem:dc-ch3-4-1-continuity}
  \leanok
  \dcref{ch3:4.1}
  Let $Z$ be a local flow of the chart-$p$ geodesic spray, $\mathrm{finv}$ a $C^2$ exponential
  inverse on an open $S\ni y$, and $(y,T^{-1}w)$ an admissible initial condition. Along the
  geodesic reading $x(s)=\varphi_p(\gamma(s))$ of the geodesic through $q=\varphi_p^{-1}(y)$
  with chart velocity $w$, the squared radial distance $F(s)=|\mathrm{finv}(x(s))|_p^2$ has
  $\partial F/\partial t(0)=2\langle D\mathrm{finv}(y)(w),\mathrm{finv}(y)\rangle_p$ and
  $$\frac{\partial^2 F}{\partial t^2}(0,q,v)=G(q,v),$$
  the algebraic second-derivative form of \cref{lem:dc-ch3-4-1-continuity}. This is the
  identification of the abstract form $G$ with the genuine second time-derivative of $F$ along
  the geodesic family — the assembly step of the Lemma 4.1.
\end{lemma}
\begin{proof}
  \sketch
  \leanok
  \uses{lem:dc-ch3-4-1-fidentity, lem:dc-ch3-3-7-descent}
  The base descent (\cref{lem:dc-ch3-3-7-descent}) gives the chart reading with $x(0)=y$,
  $x'(0)=w$ and the fixed-chart acceleration $x''(0)=(\mathrm{spray}(y,w))_2$; its open-window
  form gives differentiability of $x$ on a neighborhood of $0$, and continuity of $x$ places
  $x(s)$ in the $C^2$ domain $S$ for $s$ near $0$. Applying the F-identity
  (\cref{lem:dc-ch3-4-1-fidentity}) with $a=y$ and $\xi=(\mathrm{spray}(y,w))_2$ yields exactly
  the algebraic form $G(q,v)$ of \cref{lem:dc-ch3-4-1-continuity}.
\end{proof}

\begin{lemma}[geodesic tangent to a sphere stays outside]
  \label{lem:dc-ch3-4-1}
  \lean{Riemannian.Exponential.exists_forall_geodesic_tangent_stays_outside_ball}
  \uses{lem:dc-ch3-4-1-flowreading, lem:dc-ch3-4-1-homog, lem:dc-ch3-4-1-outside, lem:dc-ch3-4-1-posnbhd, lem:dc-ch3-3-7-descent}
  \leanok
  \dcref{ch3:4.1}
  For any $p\in M$ there exists $c>0$ such that any geodesic of $M$ tangent at
  $q\in M$ to the geodesic sphere $S_r(p)$ of radius $r<c$ stays outside the
  geodesic ball $B_r(p)$ in some neighborhood of $q$.
\end{lemma}
\begin{proof}
  \sketch
  \leanok
  \uses{lem:dc-ch3-4-1-dF, lem:dc-ch3-4-1-ddF, lem:dc-ch3-4-1-hessian-center,
    lem:dc-ch3-4-1-radialcoord, lem:dc-ch3-4-1-posdef, lem:dc-ch3-4-1-outside,
    lem:dc-ch3-4-1-continuity, lem:dc-ch3-4-1-poscenter, lem:dc-ch3-4-1-posnbhd,
    lem:dc-ch3-3-7-descent, lem:dc-ch3-4-1-homog, lem:dc-ch3-3-5-ball,
    lem:dc-ch3-4-1-fidentity, lem:dc-ch3-4-1-flowreading}
  Let $W$ be a totally normal neighborhood of $p$; by the lemma of
  homogeneity restrict to unit-speed geodesics, so to the unit bundle
  $T_1W=\{(q,v);\ q\in W,\ v\in T_qM,\ |v|=1\}$. After shrinking $W$, assume that
  $\overline{W}$ is compactly contained in a normal neighborhood $U$ of $p$; since
  $T_1\overline{W}$ is compact, decrease $I$ so that
  $\gamma(I\times T_1W)\subset U$. Let
  $\gamma:I\times T_1W\to M$, $I=(-\varepsilon,\varepsilon)$, be the differentiable
  map with $t\mapsto\gamma(t,q,v)$ the unit-speed geodesic through $q$ with velocity
  $v$ at $t=0$. Set $u(t,q,v)=\exp_p^{-1}(\gamma(t,q,v))$ and
  $F:I\times T_1W\to\mathbb{R}$, $F(t,q,v)=|u(t,q,v)|^2$, the squared "distance"
  from $p$ to a point moving along $\gamma$. Then
  
  $$
  \frac{\partial F}{\partial t}=2\Big\langle\frac{\partial u}{\partial t},u\Big\rangle,\qquad\frac{\partial^2 F}{\partial t^2}=2\Big\langle\frac{\partial^2 u}{\partial t^2},u\Big\rangle+2\Big|\frac{\partial u}{\partial t}\Big|^2.
  $$
  
  Choose $r>0$ with $\exp_p B_r(0)=B_r(p)\subset W$. If $\gamma$ is tangent to
  $S_r(p)$ at $q=\gamma(0,q,v)$, then by the Gauss lemma
  $\langle\frac{\partial u}{\partial t}(0,q,v),u(0,q,v)\rangle=0$, i.e.
  $\frac{\partial F}{\partial t}(0,q,v)=0$. It suffices to show that for $r$ small
  the critical point $(0,q,v)$ is a strict minimum. For $q=p$, $u(t,p,v)=tv$, so
  $\frac{\partial^2 F}{\partial t^2}(0,p,v)=2|v|^2=2>0$. Hence there is a
  neighborhood $V\subset W$ of $p$ with $\frac{\partial^2 F}{\partial t^2}(0,q,v)>0$
  for all $q\in V$, $|v|=1$. Take $c>0$ with $\exp_p B_c(0)\subset V$: any geodesic
  in $B_c(p)$ tangent to $S_r(p)$, $r<c$, has a strict local minimum of $F$ at
  $(0,q,v)$, so in a neighborhood of $q$ its points stay outside $B_r(p)$.
\end{proof}

\begin{lemma}[strict form of Lemma 4.1: tangent geodesic stays \emph{strictly} outside]
  \label{lem:dc-ch3-4-1-strict}
  \lean{Riemannian.Exponential.exists_forall_geodesic_tangent_stays_strictly_outside_ball}
  \uses{lem:dc-ch3-4-1-flowreading, lem:dc-ch3-4-1-homog, lem:dc-ch3-4-1-strict-outside, lem:dc-ch3-4-1-posnbhd, lem:dc-ch3-3-7-descent}
  \leanok
  \dcref{ch3:4.1}
  Same data as \cref{lem:dc-ch3-4-1}, but with the \emph{strict} conclusion. For any
  $p\in M$ there are a $C^2$ exponential inverse $\mathrm{finv}$, an open neighborhood
  $V$ of $\varphi_p(p)$, and a local geodesic flow package $(r,\varepsilon,T,Z)$ such
  that: for every base position $y\in V$ and every nonzero chart velocity $w$ with
  $(y,T^{-1}w)$ admissible, if the geodesic reading $F(s)=|\exp_p^{-1}(\gamma(s))|_p^2$
  is \emph{tangent} to the geodesic sphere at $s=0$ (i.e. $\partial F/\partial t(0)=0$,
  the Gauss lemma), then $F(0)<F(s)$ for \emph{every} $s\neq0$ near $0$: the tangent
  geodesic stays \emph{strictly} outside the ball in a punctured neighborhood of its base
  point. This is the form the convex-neighborhood proof \cref{prop:dc-ch3-4-2} consumes.
\end{lemma}
\begin{proof}
  \sketch
  \leanok
  \uses{lem:dc-ch3-4-1-flowreading, lem:dc-ch3-4-1-homog, lem:dc-ch3-4-1-strict-outside,
    lem:dc-ch3-4-1-posnbhd, lem:dc-ch3-3-7-descent}
  Verbatim the assembly of \cref{lem:dc-ch3-4-1} — the F-identity
  $\partial^2 F/\partial t^2(0,q,v)=G(q,v)$ (\cref{lem:dc-ch3-4-1-flowreading}), the strict
  positivity of $G$ on $V$ for nonzero velocities (\cref{lem:dc-ch3-4-1-posnbhd},
  \cref{lem:dc-ch3-4-1-homog}), and the tangency $\partial F/\partial t(0)=0$ — but the
  final second-derivative test is the \emph{strict} one (\cref{lem:dc-ch3-4-1-strict-outside},
  in place of the non-strict
  \cref{lem:dc-ch3-4-1-outside}, giving $F(0)<F(s)$ on the punctured neighborhood.
\end{proof}

\begin{lemma}[intrinsic strict separation: an arbitrary geodesic tangent to a sphere stays strictly outside]
  \label{lem:dc-ch3-4-2-rebase}
  \lean{Riemannian.Exponential.exists_forall_intrinsic_geodesic_tangent_strictly_outside_ball}
  \uses{lem:dc-ch3-4-1-strict, lem:dc-ch3-2-5-unique}
  \leanok
  \dcref{ch3:4.1}
  The strict Lemma 4.1 (\cref{lem:dc-ch3-4-1-strict}) is phrased for the geodesic
  \emph{flow reading} $s\mapsto\varphi_p^{-1}((Z(y,T^{-1}w)(sT))_1)$ of its own flow $Z$;
  the convex-neighborhood proof must apply it to an \emph{arbitrary} intrinsic geodesic — the
  joining geodesic of \cref{thm:dc-ch3-3-7}, re-based at its interior distance-maximum. For any
  $p\in M$ there are a $C^2$ exponential inverse $\mathrm{finv}$, an open neighborhood $V$ of
  $\varphi_p(p)$, and radii $0<r$, $0<T<\varepsilon$ such that for \emph{every} continuous
  intrinsic geodesic $\sigma$ on an open interval $(-a,a)$, $a>0$, whose base point $\sigma(0)$
  reads into $V$ (with $\sigma(0)$ in the chart at $p$) and whose chart velocity
  $w=(\varphi_p\circ\sigma)'(0)\neq0$ is admissible ($(\varphi_p(\sigma 0),T^{-1}w)$ in the flow's
  ball of initial conditions), the radial functional $F(s)=|\mathrm{finv}(\varphi_p(\sigma s))|_p^2$
  satisfies: if $\sigma$ is tangent to the geodesic sphere at $0$ ($\partial F/\partial t(0)=0$),
  then $F(0)<F(s)$ for every $s\neq0$ near $0$.
\end{lemma}
\begin{proof}
  \sketch
  \leanok
  \uses{lem:dc-ch3-4-1-strict, lem:dc-ch3-2-5-unique}
  Feed $\sigma$ into the \emph{readback} of \cref{lem:dc-ch3-2-5-unique} for the strict
  lemma's own flow $Z$: since
  the readback consumes only the generic flow clauses of $Z$ (not a particular construction), it
  applies verbatim, and on the overlap window $\sigma(s)=\varphi_p^{-1}((Z(\varphi_p(\sigma 0),
  T^{-1}w)(sT))_1)$. Hence $F$ coincides near $0$ with the strict lemma's radial functional, so
  the tangency hypothesis and the strict separation $F(0)<F(s)$ transfer across the equality of
  functions (using $\varphi_p(\sigma s)$ equals the flow reading's chart position on the window).

\end{proof}

\begin{lemma}[the radial distance has no interior local maximum along a geodesic]
  \label{lem:dc-ch3-4-2-nomax}
  \lean{Riemannian.Exponential.exists_forall_intrinsic_geodesic_not_isLocalMax_radial}
  \uses{lem:dc-ch3-4-2-rebase}
  \leanok
  \dcref{ch3:4.2}
  With the package of \cref{lem:dc-ch3-4-2-rebase}, for every admissible continuous intrinsic
  geodesic $\sigma$ the radial functional $F(s)=|\mathrm{finv}(\varphi_p(\sigma s))|_p^2$ has
  \emph{no} local maximum at the base point $0$. This is the exclusion the
  convex-neighborhood contradiction rests on: at an interior point where the distance from $p$ to
  a joining geodesic attains a maximum, $F$ would have a local maximum, which this rules out.
\end{lemma}
\begin{proof}
  \sketch
  \leanok
  \uses{lem:dc-ch3-4-2-rebase}
  A local maximum of $F$ at $0$ forces $\partial F/\partial t(0)=0$ (Fermat's theorem,
  ), whence \cref{lem:dc-ch3-4-2-rebase} gives the strict
  separation $F(0)<F(s)$ for every $s\neq0$ near $0$. But a local maximum also gives $F(s)\le F(0)$
  nearby; the two contradict on the (nonempty) punctured neighborhood $\mathbb R\setminus\{0\}$.

\end{proof}

\begin{lemma}[metric--radial bridge: the radial functional is the squared distance]
  \label{lem:dc-ch3-4-2-bridge}
  \lean{Riemannian.Exponential.sq_dist_eq_chartMetricInner_expMapInv}
  \uses{lem:dc-ch3-4-2-nomax, prop:dc-ch3-3-6}
  \leanok
  \dcref{ch3:4.2}
  Let $\mathrm{finv}$ be the $C^2$ exponential inverse of \cref{lem:dc-ch3-4-2-nomax} (which now
  exposes its left-inverse clause $\mathrm{finv}(\varphi_p(\exp_p w))=w$ for $\|w\|<\varepsilon_L$),
  and let the ambient distance be the Riemannian distance of $g$. Then there is $\rho>0$ such that
  for every chart velocity $v$ with $\|v\|<\rho$,
  $$ \big(d(p,\exp_p v)\big)^2 = \big\langle \mathrm{finv}(\varphi_p(\exp_p v)),\,
     \mathrm{finv}(\varphi_p(\exp_p v))\big\rangle_p . $$
  This is the identification $F=d^2$ that the convex-neighborhood contradiction consumes:
  the chart radial functional $F(x)=\langle\exp_p^{-1}(x),\exp_p^{-1}(x)\rangle_p$ of
  \cref{lem:dc-ch3-4-2-nomax} equals the squared metric distance from $p$, so a maximum of
  $d(p,\cdot)$ along a geodesic is a maximum of $F$.
\end{lemma}
\begin{proof}
  \sketch
  \leanok
  \uses{lem:dc-ch3-4-2-nomax, prop:dc-ch3-3-6}
  Rewrite $\mathrm{finv}(\varphi_p(\exp_p v))=v$ by the left inverse (valid for $\|v\|<\varepsilon_L$),
  reducing the claim to $\big(d(p,\exp_p v)\big)^2=\langle v,v\rangle_p$. The radial distance
  realization $d(p,\exp_p v)=\sqrt{\langle v,v\rangle_p}$ (\cref{prop:dc-ch3-3-6}, the Gauss-lemma
  consequence for the normal ball) together with $\mathrm{edist}=\mathrm{ofReal}\circ\mathrm{dist}$ in
  the ambient metric space and the positive-semidefiniteness $\langle v,v\rangle_p\ge0$ of the chart
  Gram form gives the identity on the overlap $\|v\|<\rho=\min(\varepsilon_e,\varepsilon_L)$.
\end{proof}

\begin{lemma}[the radial functional is the squared distance on a geodesic ball]
  \label{lem:dc-ch3-4-2-bridge-ball}
  \lean{Riemannian.Exponential.exists_ball_sq_dist_eq_chartMetricInner}
  \uses{lem:dc-ch3-4-2-bridge, prop:dc-ch3-3-6}
  \leanok
  \dcref{ch3:4.2}
  With the same data as \cref{lem:dc-ch3-4-2-bridge}, there is $\delta'>0$ such that on the whole
  metric ball $B_{\delta'}(p)$ the chart radial functional is the squared distance:
  $\big(d(p,x)\big)^2 = \langle \mathrm{finv}(\varphi_p x), \mathrm{finv}(\varphi_p x)\rangle_p$ for
  every $x$ with $d(p,x)<\delta'$. This is the form the max-distance contradiction consumes
  directly: on a small geodesic ball around $p$, a local maximum of $d(p,\cdot)$ along a geodesic is
  a local maximum of $F$, which \cref{lem:dc-ch3-4-2-nomax} forbids.
\end{lemma}
\begin{proof}
  \sketch
  \leanok
  \uses{lem:dc-ch3-4-2-bridge, prop:dc-ch3-3-6}
  Every $x$ with $d(p,x)<\delta'$ lies in the normal $\varepsilon_e$-ball (the escape clause of the
  normal-ball realization \cref{prop:dc-ch3-3-6}: points outside $\exp_p(B_{\varepsilon_e}(0))$ are at
  distance $\ge\delta_e$), so $x=\exp_p v$ with $d(p,x)=\sqrt{\langle v,v\rangle_p}$. The local
  coordinate-norm bound $\|v\|^2\le c\,\langle v,v\rangle_p$ (positive-definiteness of the chart Gram
  form at $p$) forces $\|v\|<\min(\varepsilon_e,\varepsilon_L)$ once
  $\delta'\le\min(\delta_e,\min(\varepsilon_e,\varepsilon_L)/\sqrt c)$, so the left inverse gives
  $\mathrm{finv}(\varphi_p x)=v$ and \cref{lem:dc-ch3-4-2-bridge} applies.
\end{proof}

\begin{lemma}[the max-distance interior deduction: convexity from no interior maximum]
  \label{lem:dc-ch3-4-2-maxdeduction}
  \lean{Riemannian.lt_of_forall_not_isLocalMax_of_le}
  \leanok
  \dcref{ch3:4.2}
  A continuous function $h$ on $[0,1]$ with $h(0)\le\beta$, $h(1)\le\beta$ and \emph{no} interior
  local maximum on $(0,1)$ stays \emph{strictly} below $\beta$ on the open interval: $h(s)<\beta$ for
  every $s\in(0,1)$. This is the convex-neighborhood contradiction stripped of the geometry:
  with $h(s)=d(p,\gamma(s))$ the distance from $p$ to a joining geodesic, the maximum of $h$ over
  $[0,1]$ sits at an endpoint (an interior maximizer would be a local maximum, forbidden by
  \cref{lem:dc-ch3-4-2-nomax} through the bridge $F=d^2$ of \cref{lem:dc-ch3-4-2-bridge-ball}), hence
  is $\le\beta$, and no interior point can equal it.
\end{lemma}
\begin{proof}
  \sketch
  \leanok
  The maximum of $h$ over the compact $[0,1]$ is attained at some $s_0$; if $s_0$ were interior it
  would be a local maximum, contradicting the hypothesis, so $s_0\in\{0,1\}$ and $h(s_0)\le\beta$.
  If some interior $s$ had $h(s)=\beta$, then $\beta\le h(s)\le h(s_0)\le\beta$ forces $h(s)=h(s_0)$,
  so $s$ is itself an interior maximizer — again a contradiction. Hence $h(s)<\beta$.
\end{proof}

\begin{definition}[strongly convex set]
  \label{def:dc-ch3-4-2-stronglyconvex}
  \lean{Riemannian.Exponential.StronglyConvex}
  \leanok
  \dcref{ch3:4.2}
  A set $S\subset M$ is \emph{strongly convex} if for any two points $q_1,q_2\in\overline{S}$
  there is a minimizing geodesic $\gamma:[0,1]\to M$ (constant-speed, realizing the distance
  $d(q_1,q_2)$ proportionally to arc length: $d(\gamma(s),\gamma(t))=|s-t|\,d(q_1,q_2)$) whose
  open-arc interior $\gamma((0,1))$ lies in $S$, and this joining geodesic is unique among such
  competitors (pointwise equality on $[0,1]$).
\end{definition}

\begin{lemma}[intrinsic speed in a fixed chart]
  \label{lem:dc-ch3-4-2-speedchart}
  \lean{Riemannian.Exponential.speedSq_eq_chartMetricInner_extChartAt}
  \uses{lem:dc-ch3-2-1-constspeed}
  \leanok
  \dcref{ch3:4.2}
  For a curve $\gamma$ continuous at $t$ with $\gamma(t)$ in the source of the chart at $\alpha$ and
  coordinate velocity $\xi=(\varphi_\alpha\circ\gamma)'(t)$, the intrinsic squared speed
  $\langle\gamma'(t),\gamma'(t)\rangle_g$ equals the chart-Gram value
  $\langle\xi,\xi\rangle$ read in the \emph{fixed} chart at $\alpha$ (not the moving chart at
  $\gamma(t)$), at the chart position $\varphi_\alpha(\gamma(t))$.
\end{lemma}
\begin{proof}
  \sketch
  \leanok
  The intrinsic velocity is the tangent-coordinate-change readback of $\xi$, and the chart-Gram form
  is the intrinsic inner product of the readbacks; composing the two identities gives the claim.
\end{proof}

\begin{lemma}[chart velocity bounded by the conserved speed]
  \label{lem:dc-ch3-4-2-velbound}
  \lean{Riemannian.Exponential.exists_sq_norm_deriv_le_speedSq}
  \uses{lem:dc-ch3-4-2-speedchart}
  \leanok
  \dcref{ch3:4.2}
  For every $p$ there are a constant $c>0$ and a neighborhood $V$ of $\varphi_p(p)$ such that, whenever
  a curve $\gamma$ reads into $V$ at $t$ with coordinate velocity $\xi=(\varphi_p\circ\gamma)'(t)$, one
  has $\|\xi\|^2\le c\,\langle\gamma'(t),\gamma'(t)\rangle_g$. Along a geodesic the right-hand side is
  constant, so this bounds the chart velocity \emph{uniformly along the geodesic} by its initial speed.
\end{lemma}
\begin{proof}
  \sketch
  \leanok
  The uniform coordinate-norm bound $\|\xi\|^2\le c\,\langle\xi,\xi\rangle_{p,y}$ (positive-definiteness
  of the Gram form, spread over a neighborhood of $\varphi_p(p)$ by the tube lemma) combined with the
  fixed-chart speed reading \cref{lem:dc-ch3-4-2-speedchart}.
\end{proof}

\begin{lemma}[the inverse pair map fixes the center diagonal]
  \label{lem:dc-ch3-4-2-invdiagonal}
  \lean{Riemannian.Exponential.exists_totallyNormal_c1_diffeo}
  \uses{thm:dc-ch3-3-7}
  \leanok
  \dcref{ch3:4.2}
  The inverse pair map $\mathrm{Ginv}$ of \cref{thm:dc-ch3-3-7} fixes the center
  diagonal: $\mathrm{Ginv}(\varphi_p(p),\varphi_p(p))=(\varphi_p(p),0)$. This is
  the $\exp_p^{-1}(p)=0$ read in the chart at $p$ — the joining velocity
  from $p$ to itself is zero.
\end{lemma}
\begin{proof}
  \sketch
  \leanok
  The pair map $G(y,w)=(y,(Z(y,w/T)T)_1)$ sends $(\varphi_p(p),0)$ to
  $(\varphi_p(p),\varphi_p(p))$, since the zero-velocity flow segment is the
  constant geodesic at $p$: $(Z(\varphi_p(p),0)T)_1=\varphi_p(p)$ (the center
  fixed point of the spray flow). As $(\varphi_p(p),0)$ lies in the product ball
  $B$ on which $\mathrm{Ginv}\circ G=\mathrm{id}$, applying the two-sided inverse
  gives $\mathrm{Ginv}(\varphi_p(p),\varphi_p(p))=(\varphi_p(p),0)$.
\end{proof}

\begin{lemma}[uniform smallness of a chart-pair functional near the center]
  \label{lem:dc-ch3-4-2-pairsmall}
  \lean{Riemannian.Exponential.exists_closedBall_forall_pair_norm_lt}
  \leanok
  \dcref{ch3:4.2}
  Let $\Psi$ be a map of chart-coordinate pairs $E\times E\to F$, continuous at the
  diagonal center $(\varphi_p(p),\varphi_p(p))$ and vanishing there. Then for every
  tolerance $r>0$ there is a radius $\beta>0$ such that the closed geodesic ball
  $\overline{B_\beta(p)}$ lies inside the chart source at $p$ and every pair of
  points $q_1,q_2\in\overline{B_\beta(p)}$ reads to a small value:
  $\|\Psi(\varphi_p(q_1),\varphi_p(q_2))\|<r$.
\end{lemma}
\begin{proof}
  \sketch
  \leanok
  By continuity of $\Psi$ at the diagonal center and $\Psi=0$ there, the set
  $\{x:\|\Psi(x)\|<r\}$ is a neighborhood of $(\varphi_p(p),\varphi_p(p))$, hence
  contains a product ball of some radius $\rho$. The chart $\varphi_p$ is
  continuous at $p$ with $\varphi_p(p)$ the center, so
  $\varphi_p^{-1}(B_\rho(\varphi_p(p)))\cap(\text{chart source})$ is a neighborhood
  of $p$ and contains $\overline{B_\beta(p)}$ for some $\beta>0$. For
  $q_1,q_2\in\overline{B_\beta(p)}$ both chart images lie in $B_\rho(\varphi_p(p))$,
  so the pair lies in the product ball and $\|\Psi(\varphi_p(q_1),\varphi_p(q_2))\|<r$.

\end{proof}

\begin{lemma}[the joining velocity is uniformly small on a small ball]
  \label{lem:dc-ch3-4-2-velsmall}
  \lean{Riemannian.Exponential.exists_closedBall_forall_ginvSnd_norm_lt}
  \uses{lem:dc-ch3-4-2-pairsmall, lem:dc-ch3-4-2-invdiagonal, thm:dc-ch3-3-7}
  \leanok
  \dcref{ch3:4.2}
  With the inverse pair map $\mathrm{Ginv}$ of \cref{thm:dc-ch3-3-7}, $C^1$ (hence
  continuous) on its open diffeomorphism image $U\ni(\varphi_p(p),\varphi_p(p))$ and
  fixing the center diagonal (\cref{lem:dc-ch3-4-2-invdiagonal}): for every $r>0$
  there is $\beta>0$ with $\overline{B_\beta(p)}$ inside the chart source and, for all
  $q_1,q_2\in\overline{B_\beta(p)}$, the joining chart velocity
  $v=(\mathrm{Ginv}(\varphi_p(q_1),\varphi_p(q_2)))_2$ obeys $\|v\|<r$.
\end{lemma}
\begin{proof}
  \sketch
  \leanok
  Apply \cref{lem:dc-ch3-4-2-pairsmall} to $\Psi(x)=(\mathrm{Ginv}(x))_2$. This
  $\Psi$ is continuous at the center (restriction of the continuous $\mathrm{Ginv}$
  to the open $U$, followed by the second projection) and vanishes there by
  \cref{lem:dc-ch3-4-2-invdiagonal}.
\end{proof}

\begin{lemma}[the interior of a joining geodesic stays inside the ball]
  \label{lem:dc-ch3-4-2-interior}
  \lean{Riemannian.Exponential.exists_forall_geodesic_dist_lt_of_admissible}
  \uses{lem:dc-ch3-4-2-nomax, lem:dc-ch3-4-2-bridge-ball, lem:dc-ch3-4-2-maxdeduction}
  \leanok
  \dcref{ch3:4.2}
  With the exponential-inverse package of \cref{lem:dc-ch3-4-2-nomax} and the bridge radius $\delta'$ of
  \cref{lem:dc-ch3-4-2-bridge-ball}: for every continuous intrinsic geodesic $\gamma$ on an open interval
  $(l,h)\supsetneq[0,1]$ whose endpoints lie within $\beta\le\delta'$ of $p$, which stays inside
  $B_{\delta'}(p)$ over $[0,1]$, and which is \cref{lem:dc-ch3-4-2-nomax}-admissible at every interior
  time (base reads into $V$, nonzero chart velocity in the flow's ball of initial conditions), the whole
  open arc stays strictly inside $B_\beta(p)$: $d(p,\gamma(t))<\beta$ for all $t\in(0,1)$.
\end{lemma}
\begin{proof}
  \sketch
  \leanok
  The distance $h(t)=d(p,\gamma(t))$ is continuous on $[0,1]$ with $h(0),h(1)\le\beta$. If $h$ had an
  interior local maximum at $t_0$, re-base $\sigma(s)=\gamma(t_0+s)$ (an affine reparametrisation of a
  geodesic is a geodesic): the metric$\leftrightarrow$radial bridge $F=d^2$ of
  \cref{lem:dc-ch3-4-2-bridge-ball} turns the local maximum of $h$ into a local maximum of the chart
  radial functional at $0$, which \cref{lem:dc-ch3-4-2-nomax} forbids. So $h$ has no interior local
  maximum, and \cref{lem:dc-ch3-4-2-maxdeduction} yields $h<\beta$ on $(0,1)$.
\end{proof}

\begin{lemma}[the chart Gram quadratic form is uniformly small near the center]
  \label{lem:dc-ch3-4-2-gramsmall}
  \lean{Riemannian.Exponential.exists_ball_forall_chartMetricInner_lt}
  \leanok
  \dcref{ch3:4.2}
  For every tolerance $\eta>0$ there is a radius $\rho>0$ such that the open chart ball
  $B_\rho(\varphi_p(p))$ lies in the chart target and, for every base position $y$ in it and every
  chart vector $w$ with $\|w\|<\rho$, the chart-Gram value $\langle w,w\rangle_{p,y}<\eta$.
\end{lemma}
\begin{proof}
  \sketch
  \leanok
  The functional $\Psi(y,w)=\langle w,w\rangle_{p,y}$ is jointly continuous (a finite sum of the
  continuous Gram components composed with the base, times the continuous chart coordinates of $w$)
  and vanishes at $(\varphi_p(p),0)$ (both vector slots are $0$), so $\{\Psi<\eta\}$ is a
  neighborhood of the center and contains a product ball.
\end{proof}

\begin{lemma}[interior arc in the ball from the minimizing property and conserved speed]
  \label{lem:dc-ch3-4-2-interior-minimizing}
  \lean{Riemannian.Exponential.exists_forall_minimizing_geodesic_interior_ball}
  \uses{lem:dc-ch3-4-2-interior, lem:dc-ch3-4-2-velbound, lem:dc-ch3-4-2-speedchart, lem:dc-ch3-4-2-gramsmall}
  \leanok
  \dcref{ch3:4.2}
  For every $p$ there are radii $\beta_0,\rho>0$ such that any minimizing geodesic $\gamma$
  (with $d(\gamma(s),\gamma(t))=|s-t|\,d(q_1,q_2)$) joining $q_1,q_2\in\overline{B_\beta(p)}$
  ($\beta\le\beta_0$), extended as a continuous intrinsic geodesic on an open window
  $(l,h)\supsetneq[0,1]$ with nonzero small initial chart velocity $w$ ($w\neq0$, $\|w\|<\rho$) and a
  chart-velocity derivative at every interior time, has its open arc strictly inside $B_\beta(p)$:
  $d(p,\gamma(t))<\beta$ for all $t\in(0,1)$.
\end{lemma}
\begin{proof}
  \sketch
  \leanok
  This discharges the admissibility of \cref{lem:dc-ch3-4-2-interior} without any flow-continuity
  (tube) plumbing. \emph{Base confinement} $d(p,\gamma(t))\le3\beta$ follows from the triangle
  inequality on the minimizing length ($d(\gamma(0),\gamma(t))=t\,d(q_1,q_2)\le2\beta$), which puts
  the base into every prescribed chart neighborhood of $p$ and inside $B_{\delta'}(p)$. The squared
  speed is conserved along the geodesic, and at $t=0$ equals the small chart-Gram value
  $\langle w,w\rangle$ (\cref{lem:dc-ch3-4-2-speedchart}); the coordinate-velocity bound
  \cref{lem:dc-ch3-4-2-velbound} together with the chart-Gram smallness \cref{lem:dc-ch3-4-2-gramsmall}
  then bounds the interior chart velocity $T^{-1}w_0$ inside the flow's ball of initial conditions,
  while the Gram lower bound keeps the conserved speed positive so that $w_0\neq0$.
  \cref{lem:dc-ch3-4-2-interior} then yields the strict interior confinement.
\end{proof}

\begin{lemma}[a geodesic segment has length equal to speed times elapsed parameter]
  \label{lem:dc-ch3-geodesic-pathlength}
  \lean{Riemannian.Geodesic.IsGeodesicOn.pathELength_eq}
  \uses{lem:dc-ch3-2-1-constspeed}
  \leanok
  \dcref{ch3:3.6}
  Let $\gamma$ be a continuous intrinsic geodesic on a connected open set of times containing
  $a$ and $b$. Then the length of the arc over $[a,b]$ is the constant speed times the elapsed
  parameter,
  \[
    \ell\big(\gamma|_{[a,b]}\big)=\sqrt{\langle\gamma'(a),\gamma'(a)\rangle_g}\,(b-a).
  \]
\end{lemma}
\begin{proof}
  \sketch
  \leanok
  The length is the integral of the intrinsic speed $\sqrt{\langle\gamma',\gamma'\rangle_g}$ over
  $[a,b]$. Along a geodesic the squared speed is constant (\cref{lem:dc-ch3-2-1-constspeed}), equal
  to its value at $a$, so the integrand is the constant $\sqrt{\langle\gamma'(a),\gamma'(a)\rangle_g}$
  and the integral is that constant times $b-a$.
\end{proof}

\begin{lemma}[the joining geodesic segment has length $\sqrt{\langle w,w\rangle}$, uniformly in the base]
  \label{lem:dc-ch3-4-2-segment-length}
  \lean{Riemannian.Geodesic.pathELength_uniform_flow_segment_Ioo}
  \uses{lem:dc-ch3-geodesic-pathlength, lem:dc-ch3-3-7-descent, lem:dc-ch3-3-7-descent-velocity, lem:dc-ch3-4-2-speedchart}
  \leanok
  \dcref{ch3:4.2}
  For the rescaled chart-$p$ spray segment $\gamma(s)=\varphi_p^{-1}\big((Z(y,T^{-1}w)(sT))_1\big)$
  joining $q=\varphi_p^{-1}(y)$ to a nearby point (the geodesic of \cref{thm:dc-ch3-3-7}), the length
  over $[0,1]$ equals the chart-Gram norm of the initial velocity read at the base $y$:
  \[
    \ell\big(\gamma|_{[0,1]}\big)=\sqrt{\langle w,w\rangle_{p,y}}.
  \]
  This is uniform in the base point $y$: the same flow $Z$ and time scale $T$ serve every $q$ in the
  neighborhood. It is the metric payload that \cref{thm:dc-ch3-3-7} otherwise omits.
\end{lemma}
\begin{proof}
  \sketch
  \leanok
  The segment is a continuous intrinsic geodesic on the open time window, so
  \cref{lem:dc-ch3-geodesic-pathlength} gives its length as
  $\sqrt{\langle\gamma'(0),\gamma'(0)\rangle_g}$. Its squared speed at $s=0$ is the chart-Gram value
  $\langle w,w\rangle$ read in the fixed chart at $p$ (\cref{lem:dc-ch3-4-2-speedchart}), at the base
  position $\varphi_p(\gamma(0))=y$ (the flow starts at its initial condition).
\end{proof}

\begin{lemma}[uniform radial upper bound]
  \label{lem:dc-ch3-4-2-segment-edist}
  \lean{Riemannian.Geodesic.edist_uniform_flow_segment_le}
  \uses{lem:dc-ch3-4-2-segment-length}
  \leanok
  \dcref{ch3:4.2}
  For the joining geodesic segment of \cref{thm:dc-ch3-3-7}, the distance from the base
  $q=\varphi_p^{-1}(y)$ to its endpoint is at most the chart-Gram norm of the initial velocity,
  \[
    d\big(q,\gamma(1)\big)\le\sqrt{\langle w,w\rangle_{p,y}},
  \]
  uniformly in the base point $y$.
\end{lemma}
\begin{proof}
  \sketch
  \leanok
  The distance between the endpoints of a $C^1$ curve is at most its length; combine this with the
  length identity \cref{lem:dc-ch3-4-2-segment-length}. This is the $\le$ half of the normal-ball
  minimizing property (\cref{prop:dc-ch3-3-6}), now available \emph{uniformly} in the base point $y$
  --- the direction that does not require a base-uniform Gauss estimate.
\end{proof}

\begin{lemma}[the joining geodesic carries the base-uniform metric payload]
  \label{lem:dc-ch3-4-2-join-metric}
  \lean{Riemannian.Exponential.exists_convex_join_ball}
  \uses{thm:dc-ch3-3-7, lem:dc-ch3-3-7-descent, lem:dc-ch3-4-2-segment-length, lem:dc-ch3-4-2-segment-edist, lem:dc-ch3-4-2-velsmall, lem:dc-ch3-geodesic-pathlength}
  \leanok
  \dcref{ch3:4.2}
  For every $p$ and every velocity tolerance $\rho>0$ there is $\beta>0$ such that any two points
  $q_1,q_2\in\overline{B_\beta(p)}$ are joined by a geodesic
  $\gamma(s)=\varphi_p^{-1}\big((Z(\varphi_p q_1,T^{-1}w)(sT))_1\big)$ on an open time window
  $(l,h)\supsetneq[0,1]$, with joining chart velocity $\|w\|<\rho$ (and $w\neq0$ unless $q_1=q_2$),
  whose length is arclength-proportional,
  $\ell(\gamma|_{[0,t]})=\sqrt{\langle w,w\rangle_{p,\varphi_p q_1}}\,t$, and whose endpoints obey the
  radial upper bound $d(q_1,q_2)\le\sqrt{\langle w,w\rangle_{p,\varphi_p q_1}}$. The chart velocity at
  every interior time and the geodesic equation are also available. This is the metric-laden
  totally-normal joining geodesic, uniform over the ball of base points.
\end{lemma}
\begin{proof}
  \sketch
  \leanok
  The totally-normal diffeomorphism \cref{thm:dc-ch3-3-7} joins $q_1,q_2$ by the flow segment and now
  re-exposes the raw spray-flow predicate with the same $(T,Z)$, so the descent
  \cref{lem:dc-ch3-3-7-descent} gives the geodesic on the open window and
  \cref{lem:dc-ch3-4-2-segment-length} its length $\sqrt{\langle w,w\rangle}$ (the constant-speed
  identity \cref{lem:dc-ch3-geodesic-pathlength}); the radial upper bound is
  \cref{lem:dc-ch3-4-2-segment-edist}. The joining velocity is uniformly small by
  \cref{lem:dc-ch3-4-2-velsmall}, and $w=0$ forces zero length hence $q_1=q_2$.
\end{proof}

\begin{lemma}[a short geodesic reaches the radial flow endpoint]
  \label{lem:dc-ch3-4-2-endpoint-id}
  \lean{Riemannian.Exponential.movingBase_geodesic_endpoint_eq_flow_reading}
  \uses{thm:dc-ch3-3-7, lem:dc-ch3-3-7-descent, def:dc-ch3-2-1}
  \leanok
  \dcref{ch3:3.6}
  Let $Z$ be a local flow of the coordinate geodesic spray at $p$ on a ball around the zero section
  (as in \cref{thm:dc-ch3-3-7}). Let $\gamma$ be a continuous intrinsic geodesic on an open window
  $(l,h)\supsetneq[0,1]$ with $\gamma(0)=q_1$ near $p$ and initial chart velocity $w$ small enough that
  $(\varphi_p q_1,T^{-1}w)$ lies in the flow ball. Then $\gamma$ coincides with the rescaled flow line
  $s\mapsto\varphi_p^{-1}\big((Z(\varphi_p q_1,T^{-1}w)(sT))_1\big)$ on the overlap window; in
  particular $\gamma(1)=\varphi_p^{-1}\big((Z(\varphi_p q_1,T^{-1}w)T)_1\big)$.
\end{lemma}
\begin{proof}
  \sketch
  \leanok
  The flow line is a continuous intrinsic geodesic on $(-(\varepsilon/T),\varepsilon/T)$
  (\cref{lem:dc-ch3-3-7-descent}) starting at $\varphi_p^{-1}(\varphi_p q_1)=q_1$ with chart velocity
  $w$ at time $0$ --- the same position and velocity as $\gamma$. On the asymmetric open overlap
  $(l,h)\cap(-(\varepsilon/T),\varepsilon/T)$, which contains $[0,1]$ because $h>1$ and
  $\varepsilon/T>1$, intrinsic geodesic uniqueness identifies the two curves; evaluating at $s=1$ gives
  the endpoint identity. Using the asymmetric overlap avoids the symmetric window $(-a,a)$, so the raw
  $l<0<1<h$ data suffices.
\end{proof}

\begin{lemma}[base-uniform competitor bound with escape]
  \label{lem:dc-ch3-4-2-competitor-lb}
  \lean{Riemannian.Exponential.exists_movingBase_competitor_le_pathELength}
  \uses{thm:dc-ch3-3-7, lem:dc-ch3-3-6-radius, lem:dc-ch3-3-6-competitor, lem:dc-ch3-4-2-gramsmall}
  \leanok
  \dcref{ch3:3.6}
  For every $p$ there are $\beta,\rho'>0$, a time scale $T$ and a spray flow $Z$ such that for every
  base point $q\in\overline{B_\beta(p)}$ and every chart velocity $\|w\|<\rho'$, every $C^1$ curve
  $\sigma$ from $q$ to the radial flow endpoint $\varphi_p^{-1}\big((Z(\varphi_p q,T^{-1}w)T)_1\big)$
  has length at least $\sqrt{\langle w,w\rangle_{p,\varphi_p q}}$. This is the Proposition 3.6
  competitor inequality (\cref{lem:dc-ch3-3-6-competitor}), stated over the whole ball of centres $q$
  at once.
\end{lemma}
\begin{proof}
  \sketch
  \leanok
  Fix $q$ and read the totally-normal diffeomorphism (\cref{thm:dc-ch3-3-7}) at $y=\varphi_p q$: the
  slice $f_y(z)=(Z(y,T^{-1}z)T)_1$ is a $C^1$ diffeomorphism of a ball onto a chart region, with
  inverse $\mathrm{finv}_y=(\mathrm{Ginv}(y,\cdot))_2$. If $\sigma$ stays in the closed $r'$-region
  $\{\varphi_p^{-1}(f_y z):\|z\|\le r'\}$, the (moving-base) Gauss radius comparison
  (\cref{lem:dc-ch3-3-6-radius}) applied to the polar lift $\mathrm{finv}_y\circ\varphi_p\circ\sigma$
  gives $\sqrt{\langle\mathrm{finv}_y(\varphi_p\sigma(1)),\cdot\rangle_y}\le\ell(\sigma)$; since the
  endpoint is $\varphi_p^{-1}(f_y w)$ and $\mathrm{finv}_y(f_y w)=w$, this is
  $\sqrt{\langle w,w\rangle_y}$. If instead $\sigma$ leaves the region, at the first exit time it
  crosses the coordinate $r'$-sphere at some $z_0$ with $\|z_0\|=r'$; the confined comparison up to
  that time plus the coordinate-norm bound $\|z_0\|^2\le c\,\langle z_0,z_0\rangle_y$ gives
  $\ell(\sigma)\ge r'/\sqrt c>\sqrt{\langle w,w\rangle_y}$, the last inequality by the base-uniform
  smallness (\cref{lem:dc-ch3-4-2-gramsmall}) for $\|w\|<\rho'$. The region is open (the slice inverse
  is a homeomorphism) and its closure is compact, so the first exit time is well defined.
\end{proof}

\begin{lemma}[base-uniform lower bound: short geodesics realize the distance]
  \label{lem:dc-ch3-4-2-lowerbound}
  \lean{Riemannian.Exponential.exists_movingBase_prop36_lower_bound}
  \uses{lem:dc-ch3-4-2-competitor-lb, lem:dc-ch3-4-2-endpoint-id, lem:dc-ch3-geodesic-pathlength, lem:dc-ch3-4-2-speedchart, prop:dc-ch3-3-6}
  \leanok
  \dcref{ch3:3.6}
  For every $p$ there are $\beta_H,\rho_H>0$ such that every geodesic $\gamma$ on an open window
  $(l,h)\supsetneq[0,1]$ joining $q_1,q_2$ within $\beta_H$ of $p$, with small initial chart velocity
  $\|w\|<\rho_H$ and constant-speed length rate $\ell$ (i.e.\ $\ell(\gamma|_{[0,t]})=\ell t$), realizes
  the distance between its endpoints: $\ell\le d(q_1,q_2)$. This is the
  base-uniform form of the distance-realization result in
  \cref{prop:dc-ch3-3-6}.
\end{lemma}
\begin{proof}
  \sketch
  \leanok
  Constant speed (\cref{lem:dc-ch3-geodesic-pathlength}), read in the fixed chart at $p$
  (\cref{lem:dc-ch3-4-2-speedchart}), forces $\ell=\sqrt{\langle w,w\rangle_{p,\varphi_p q_1}}$. The
  flow-reading endpoint identification (\cref{lem:dc-ch3-4-2-endpoint-id}) pins $q_2$ as the radial
  endpoint of $w$, so the base-uniform competitor bound (\cref{lem:dc-ch3-4-2-competitor-lb}) applies:
  every $C^1$ curve from $q_1$ to $q_2$ has length at least $\sqrt{\langle w,w\rangle}=\ell$. Since the
  distance is the infimum of $C^1$-curve lengths (the metric is the Riemannian distance),
  $d(q_1,q_2)\ge\ell$.
\end{proof}

\begin{lemma}[minimizing geodesics with interior in the ball, modulo the lower bound]
  \label{lem:dc-ch3-4-2-minimizing-interior}
  \lean{Riemannian.Exponential.exists_minimizing_interior_ball_of_lower_bound}
  \uses{lem:dc-ch3-4-2-join-metric, lem:dc-ch3-4-2-interior-minimizing, prop:dc-ch3-3-6, cor:dc-ch3-3-9, def:dc-ch3-4-2-stronglyconvex, lem:dc-ch3-4-2-lowerbound}
  \leanok
  \dcref{ch3:4.2}
  Assume the following base-uniform lower-bound property: there are $\beta_H,\rho_H>0$ such that every
  \emph{geodesic} $\gamma$ (on an open window $\supsetneq[0,1]$) joining $q_1,q_2$ within $\beta_H$ of
  $p$, with \emph{small} initial chart velocity ($\|w\|<\rho_H$) and constant-speed length
  $\ell(\gamma|_{[0,t]})=\ell t$, satisfies $\ell\le d(q_1,q_2)$ (a \emph{short} geodesic near $p$
  \emph{realizes} the distance between its endpoints --- \cref{prop:dc-ch3-3-6}, phrased
  base-uniformly). The geodesic and smallness hypotheses are essential: for an \emph{arbitrary}
  constant-speed curve one always has $d(q_1,q_2)\le\ell$, so a bare $\ell\le d(q_1,q_2)$ would force
  every constant-speed curve minimizing --- false on any manifold of dimension $\ge2$ (a constant-speed
  loop through $q_1=q_2$ has $\ell>0=d$). Then there is
  $\beta>0$ such that any $q_1,q_2\in\overline{B_\beta(p)}$ are joined by a \emph{minimizing} geodesic
  ($d(\gamma(s),\gamma(t))=|s-t|\,d(q_1,q_2)$) whose open arc lies in $\overline{B_\beta(p)}$. This is
  the strong convexity (\cref{def:dc-ch3-4-2-stronglyconvex}) \emph{minus} the uniqueness
  clause, under this lower-bound hypothesis.
\end{lemma}
\begin{proof}
  \sketch
  \leanok
  The joining geodesic and its length payload $\ell=\sqrt{\langle w,w\rangle}$ come from
  \cref{lem:dc-ch3-4-2-join-metric}; the lower-bound hypothesis supplies the matching lower bound, so
  $d(q_1,q_2)=\ell$,
  and the arclength distance-realization of \cref{cor:dc-ch3-3-9}'s segment machinery propagates
  $d(\gamma(s),\gamma(t))=|s-t|\,d(q_1,q_2)$ from the endpoints. The interior-arc clause is
  \cref{lem:dc-ch3-4-2-interior-minimizing}; the degenerate $q_1=q_2$ case is the constant geodesic.

\end{proof}

\begin{lemma}[strong convexity of the closed ball, modulo the lower bound and uniqueness]
  \label{lem:dc-ch3-4-2-stronglyconvex-reduction}
  \lean{Riemannian.Exponential.exists_stronglyConvex_closedBall_of_lower_bound}
  \uses{lem:dc-ch3-4-2-minimizing-interior, def:dc-ch3-4-2-stronglyconvex, thm:dc-ch3-3-7, prop:dc-ch3-3-6}
  \leanok
  \dcref{ch3:4.2}
  Assume the base-uniform lower-bound property from
  \cref{lem:dc-ch3-4-2-minimizing-interior} and local uniqueness of minimizing
  geodesics near $p$: there is $\beta_U>0$ such that
  any two constant-speed distance-realizing geodesics joining the same $q_1,q_2$ within $\beta_U$ of
  $p$ coincide on $[0,1]$ (the reading of the injectivity of $\exp_{q_1}$). Then there is
  $\beta>0$ such that the \emph{closed} ball $\overline{B_\beta(p)}$ is strongly convex
  (\cref{def:dc-ch3-4-2-stronglyconvex}). This is the whole of the Proposition 4.2 for the
  closed ball.
\end{lemma}
\begin{proof}
  \sketch
  \leanok
  The existence, minimizing and interior clauses are \cref{lem:dc-ch3-4-2-minimizing-interior}
  (shrinkable in the radius, so the strong-convexity radius $\beta=\min(B_0,\beta_U)$ couples
  correctly); the closure of a closed ball is itself, so the quantifier over $\overline{S}$ reads on
  $\overline{B_\beta(p)}$; local uniqueness supplies the uniqueness clause. The closed ball is used
  because the literal open-ball statement fails at a boundary diagonal $q_1=q_2\in\partial B_\beta(p)$
  (constant geodesic, interior $\{q_1\}\not\subset$ open ball).
\end{proof}

\begin{lemma}[minimizing geodesics with interior in the ball --- unconditional]
  \label{lem:dc-ch3-4-2-minimizing-interior-uncond}
  \lean{Riemannian.Exponential.exists_minimizing_interior_ball}
  \uses{lem:dc-ch3-4-2-minimizing-interior, lem:dc-ch3-4-2-lowerbound}
  \leanok
  \dcref{ch3:4.2}
  There is $\beta>0$ such that any $q_1,q_2\in\overline{B_\beta(p)}$ are joined by a \emph{minimizing}
  geodesic ($d(\gamma(s),\gamma(t))=|s-t|\,d(q_1,q_2)$) whose open arc lies in $\overline{B_\beta(p)}$.
  The lower-bound hypothesis in \cref{lem:dc-ch3-4-2-minimizing-interior} is supplied by
  \cref{lem:dc-ch3-4-2-lowerbound}; hence this gives the existence, minimizing, and interior
  clauses of strong convexity (\cref{def:dc-ch3-4-2-stronglyconvex}).
\end{lemma}
\begin{proof}
  \sketch
  \leanok
  Apply the lower-bound lemma \cref{lem:dc-ch3-4-2-lowerbound} to the reduction
  \cref{lem:dc-ch3-4-2-minimizing-interior}.
\end{proof}

\begin{lemma}[strong convexity of the closed ball, reduced to uniqueness alone]
  \label{lem:dc-ch3-4-2-stronglyconvex-uniq}
  \lean{Riemannian.Exponential.exists_stronglyConvex_closedBall_of_uniq}
  \uses{lem:dc-ch3-4-2-stronglyconvex-reduction, lem:dc-ch3-4-2-lowerbound, lem:dc-ch3-4-2-minimizing-interior-uncond}
  \leanok
  \dcref{ch3:4.2}
  Assume only local uniqueness of minimizing geodesics near $p$ (any two
  constant-speed distance-realizing geodesics joining the same $q_1,q_2$ within $\beta_U$ of $p$
  coincide on $[0,1]$). Then there is $\beta>0$ such that the \emph{closed} ball
  $\overline{B_\beta(p)}$ is strongly convex (\cref{def:dc-ch3-4-2-stronglyconvex}). This is
  \cref{lem:dc-ch3-4-2-stronglyconvex-reduction} together with
  \cref{lem:dc-ch3-4-2-lowerbound}.
\end{lemma}
\begin{proof}
  \sketch
  \leanok
  Apply \cref{lem:dc-ch3-4-2-lowerbound} to
  \cref{lem:dc-ch3-4-2-stronglyconvex-reduction}.
\end{proof}

\begin{lemma}[moving-base geodesic coincides with its flow reading on the overlap window]
  \label{lem:dc-ch3-4-2-eqon-flow-reading}
  \lean{Riemannian.Exponential.movingBase_geodesic_eqOn_flow_reading}
  \uses{lem:dc-ch3-4-2-endpoint-id, thm:dc-ch3-3-7, lem:dc-ch3-3-4}
  \leanok
  \dcref{ch3:4.2}
  Let $Z$ be the geodesic-spray flow of \cref{thm:dc-ch3-3-7}. A continuous intrinsic geodesic
  $\gamma$ on an open window $(l,h)\ni0$ whose foot $\gamma(0)=q_1$ lies in the chart at $p$ and whose
  chart-$p$ coordinate velocity at $0$ is $w$ \emph{coincides} with the flow-reading geodesic
  $s\mapsto\varphi_p^{-1}\bigl((Z(\varphi_p q_1,T^{-1}w)(sT))_1\bigr)$ on the overlap window
  $(l,h)\cap(-(\varepsilon/T),\varepsilon/T)$.
\end{lemma}
\begin{proof}
  \sketch
  \leanok
  Both curves are geodesics on the overlap window and share the position $q_1$ and the chart-$p$
  velocity $w$ at $0$; intrinsic uniqueness (\cref{lem:dc-ch3-3-4}) identifies them on the whole
  window. This is the \emph{EqOn} strengthening of the single-endpoint identification
  \cref{lem:dc-ch3-4-2-endpoint-id}, freed of the $1<h$ restriction (only $l<0<h$ is used).
\end{proof}

\begin{lemma}[base-uniform geodesic-uniqueness engine]
  \label{lem:dc-ch3-4-2-uniqueness-engine}
  \lean{Riemannian.Exponential.exists_movingBase_geodesic_uniqueness}
  \uses{lem:dc-ch3-4-2-eqon-flow-reading, thm:dc-ch3-3-7}
  \leanok
  \dcref{ch3:4.2}
  For any $p$ there are radii $\beta,\rho_w>0$ such that any two continuous intrinsic geodesics
  $\gamma_1,\gamma_2$ on open windows $\supsetneq[0,1]$, both joining $q_1\in\overline{B_\beta(p)}$ to
  a common $q_2\in\overline{B_\beta(p)}$, with initial chart-$p$ coordinate velocities $w_1,w_2$ of
  norm $<\rho_w$, coincide on $[0,1]$. Thus sufficiently short geodesics with
  the same endpoints and initial chart velocity are unique in this neighborhood.
\end{lemma}
\begin{proof}
  \sketch
  \leanok
  Each $\gamma_i$ coincides with the flow-reading geodesic $g_{w_i}$ on a window containing $[0,1]$
  (\cref{lem:dc-ch3-4-2-eqon-flow-reading}), so $g_{w_i}(1)=\gamma_i(1)=q_2$. The unique-velocity
  clause of the totally-normal $C^1$ diffeomorphism (\cref{thm:dc-ch3-3-7}) --- the injectivity
  of $\exp_{q_1}$ --- forces $w_1=w_2$, whence $g_{w_1}=g_{w_2}$ and therefore $\gamma_1=\gamma_2$ on
  $[0,1]$.
\end{proof}

\begin{lemma}[a closed-interval geodesic extends to an open time window]
  \label{lem:dc-ch3-4-2-window}
  \lean{Riemannian.Exponential.exists_geodesic_window_of_isGeodesicOn_Icc}
  \uses{prop:dc-ch3-2-5, lem:dc-ch3-3-4}
  \leanok
  \dcref{ch3:4.2}
  A continuous intrinsic geodesic $\alpha$ on the \emph{closed} interval $[0,1]$ extends to a
  continuous intrinsic geodesic on an \emph{open} window $(l,h)$ with $l<0<1<h$, agreeing with
  $\alpha$ on $[0,1]$.
\end{lemma}
\begin{proof}
  \sketch
  \leanok
  The endpoint values $\alpha(0),\alpha(1)$ are known points, not Cauchy limits, so a
  completeness-free forward prolongation of a geodesic on $(a,b)$ --- given the (known) endpoint
  limit --- applies at each end: once forward past $1$ with limit $\alpha(1)$, and once, after the
  time reversal $t\mapsto\alpha(-t)$, forward past $0$ with limit $\alpha(0)$. Each prolongation
  solves the geodesic ODE at the known endpoint by local existence (\cref{prop:dc-ch3-2-5}). The two
  one-sided extensions are glued at the interior time $1/2$, where they both equal $\alpha$; they
  agree there because a geodesic is locally determined by its position and velocity
  (\cref{lem:dc-ch3-3-4}).
\end{proof}

\begin{lemma}[a minimizing geodesic has speed equal to the distance it realizes]
  \label{lem:dc-ch3-4-2-speed-dist}
  \lean{Riemannian.Exponential.sqrt_speedSq_eq_dist_of_minimizing}
  \uses{prop:dc-ch3-3-6, lem:dc-ch3-3-4}
  \leanok
  \dcref{ch3:4.2}
  Let $\gamma$ be a continuous geodesic on an open window $(l,h)$ with $l<0<1<h$, joining
  $q_1=\gamma(0)$ to $q_2=\gamma(1)$ and realizing the distance on $[0,1]$
  (i.e. $d(\gamma(s),\gamma(t))=|s-t|\,d(q_1,q_2)$). Then its constant speed equals that distance:
  $\sqrt{\langle\gamma'(0),\gamma'(0)\rangle_g}=d(q_1,q_2)$.
\end{lemma}
\begin{proof}
  \sketch
  \leanok
  The inequality $d(q_1,q_2)\le\sqrt{S}$ (with $S$ the squared speed) is the length bound: a geodesic
  is at least as long as the chord between its endpoints. For the reverse, rescale the velocity by a
  small $\kappa>0$ so that the exponential ray $\eta(s)=\exp_{q_1}(s\,\kappa w)$ is defined; it is a
  geodesic through $q_1$ with the same chart velocity at $0$ as the rescaled curve
  $s\mapsto\gamma(\kappa s)$, so the two coincide near $0$ by intrinsic uniqueness
  (\cref{lem:dc-ch3-3-4}). The radial isometry of the normal ball (\cref{prop:dc-ch3-3-6}) then reads
  $d(q_1,\gamma(\kappa s))=s\,\kappa\sqrt{S}$ for small $s>0$, while the distance-realizing hypothesis
  reads the same distance as $\kappa s\,d(q_1,q_2)$; cancelling $\kappa s>0$ gives
  $\sqrt{S}=d(q_1,q_2)$.
\end{proof}

\begin{proposition}[convex neighborhoods]
  \label{prop:dc-ch3-4-2}
  \lean{Riemannian.Exponential.exists_stronglyConvex_closedBall}
  \uses{thm:dc-ch3-3-7, lem:dc-ch3-4-1, lem:dc-ch3-4-1-strict, lem:dc-ch3-4-2-rebase, lem:dc-ch3-4-2-window, lem:dc-ch3-4-2-speed-dist, lem:dc-ch3-4-2-nomax, lem:dc-ch3-4-2-bridge, lem:dc-ch3-4-2-bridge-ball, lem:dc-ch3-4-2-maxdeduction, lem:dc-ch3-4-2-speedchart, lem:dc-ch3-4-2-velbound, lem:dc-ch3-4-2-interior, lem:dc-ch3-4-2-interior-minimizing, lem:dc-ch3-4-2-gramsmall, lem:dc-ch3-4-2-invdiagonal, lem:dc-ch3-4-2-pairsmall, lem:dc-ch3-4-2-velsmall, lem:dc-ch3-4-2-segment-length, lem:dc-ch3-4-2-segment-edist, lem:dc-ch3-4-2-join-metric, lem:dc-ch3-4-2-minimizing-interior, lem:dc-ch3-4-2-stronglyconvex-reduction, lem:dc-ch3-4-2-lowerbound, lem:dc-ch3-4-2-minimizing-interior-uncond, lem:dc-ch3-4-2-stronglyconvex-uniq, lem:dc-ch3-4-2-endpoint-id, lem:dc-ch3-4-2-competitor-lb, lem:dc-ch3-4-2-eqon-flow-reading, lem:dc-ch3-4-2-uniqueness-engine, def:dc-ch3-4-2-stronglyconvex}
  \leanok
  \dcref{ch3:4.2}
  For any $p\in M$ there exists $\beta>0$ such that the closed geodesic ball $\overline{B_\beta(p)}$ is
  strongly convex (\cref{def:dc-ch3-4-2-stronglyconvex}). (The closed ball is used deliberately: the
  literal open-ball form is unsatisfiable at a boundary diagonal, see below.)
\end{proposition}
\begin{proof}
  \sketch
  \leanok
  Let $c$ be as in \cref{lem:dc-ch3-4-1}. Choose $\delta>0$ and $W$ as in \cref{thm:dc-ch3-3-7}
  with $\delta<\frac{c}{2}$, and take $\beta<\delta$ with $\overline{B_\beta(p)}\subset W$. Let
  $q_1,q_2\in\overline{B_\beta(p)}$ and let $\gamma$ be the unique geodesic of length
  $<2\delta<c$ joining them; then $\gamma\subset B_c(p)$. If the interior of
  $\gamma$ is not contained in $B_\beta(p)$, there is an interior point $m$ where the
  maximum distance $r$ from $p$ to $\gamma$ is attained; near $m$ the points of
  $\gamma$ remain in $\overline{B_r(p)}$. Since $m\in B_c(p)$ this contradicts
  \cref{lem:dc-ch3-4-1-strict}: at $m$ the tangency $\partial F/\partial t(0)=0$ holds
  (Gauss lemma) and the strict separation $F(0)<F(s)$ for $s\neq0$ forces points of
  $\gamma$ near $m$ \emph{strictly} outside $\overline{B_r(p)}$, contradicting the
  maximum. Hence the interior of $\gamma$ lies in $B_\beta(p)$, proving the proposition.
\end{proof}

\begin{definition}[piecewise parallel transport]
  \label{def:dc-ch3-3-1-piecewise-transport}
  \dcref{ch3:3.1}
  \uses{def:dc-ch3-3-1}
  \lean{Riemannian.Geodesic.IsPiecewiseParallelAlong, Riemannian.Geodesic.piecewiseTransport,
    Riemannian.Geodesic.Segmentation.ofLocal,
    Riemannian.Geodesic.local_of_segmentation,
    Riemannian.Geodesic.exists_isPiecewiseParallelAlong_of_local}
  \leanok
  Parallel transport along a piecewise differentiable curve is obtained by
  transporting a single own-foot field in succession on each differentiable
  segment of a finite partition, with the values agreeing at every vertex.  The
  endpoint transport is the relation between the initial and terminal values of
  such a field.
\end{definition}

\begin{definition}[vertex angle]
  \label{def:dc-ch3-3-1-vertex-angle}
  \dcref{ch3:3.1}
  \uses{def:dc-ch3-3-1}
  \lean{Riemannian.Geodesic.leftVelocity, Riemannian.Geodesic.rightVelocity, Riemannian.Geodesic.vertexAngle,
    Riemannian.Geodesic.HasLeftVelocity, Riemannian.Geodesic.HasRightVelocity,
    Riemannian.Geodesic.hasLeftVelocity_leftVelocity,
    Riemannian.Geodesic.hasRightVelocity_rightVelocity,
    Riemannian.Geodesic.hasLeftVelocity_of_segmentation,
    Riemannian.Geodesic.hasRightVelocity_of_segmentation,
    Riemannian.Geodesic.hasVelocities_of_segmentation,
    Riemannian.RiemannianMetric.angle_eq_pi_iff}
  \leanok
  At an interior vertex $c(t_i)$, the \emph{vertex angle} is the angle between the outgoing
  one-sided velocity $\lim_{t\to t_i^+}c'(t)$ and the incoming one-sided velocity
  $\lim_{t\to t_i^-}c'(t)$, measured by the Riemannian metric.
\end{definition}

\begin{definition}[formal parametrized-surface interface]
  \label{def:dc-ch3-3-3-formal}
  \dcref{ch3:3.3}
  \lean{Riemannian.Geodesic.IsParametrizedSurface,
    Riemannian.Geodesic.IsParametrizedSurfaceOfOrder,
    Riemannian.Geodesic.IsSmoothParametrizedSurface,
    Riemannian.Geodesic.IsSmoothExtendedParametrizedSurface,
    Riemannian.Geodesic.IsExtendedParametrizedSurfaceOfOrder.regularOn,
    Riemannian.Geodesic.IsExtendedParametrizedSurfaceOfOrder.mdifferentiableAt,
    Riemannian.Geodesic.ParametrizedSurface.IsVectorFieldAlong,
    Riemannian.Geodesic.ParametrizedSurface.partialU,
    Riemannian.Geodesic.ParametrizedSurface.partialV,
    Riemannian.Geodesic.IsParametrizedSurfaceOfOrder.surfaceSliceU_contMDiffOn,
    Riemannian.Geodesic.IsParametrizedSurfaceOfOrder.surfaceSliceV_contMDiffOn,
    Riemannian.Geodesic.IsExtendedParametrizedSurfaceOfOrder.partialU_eq_slice_mfderiv,
    Riemannian.Geodesic.IsExtendedParametrizedSurfaceOfOrder.partialV_eq_slice_mfderiv,
    Riemannian.Geodesic.ParametrizedSurface.partialU_eq_slice_mfderiv,
    Riemannian.Geodesic.ParametrizedSurface.partialV_eq_slice_mfderiv,
    Riemannian.Geodesic.IsSmoothParametrizedSurfaceWithBoundary,
    Riemannian.Geodesic.IsParametrizedSurfaceWithBoundary.surfaceDomain,
    Riemannian.Geodesic.PiecewiseC1BoundaryParametrization.interior_vertex_nonPi,
    Riemannian.Geodesic.PiecewiseC1BoundaryParametrization.closing_vertex_nonPi,
    Riemannian.Geodesic.IsNonPiAngle}
  \leanok
  The formal interface consists of the order-graded map predicate on $A$,
  the connected/open-core domain record, and the open-neighborhood extension
  record.  It also supplies the $C^r$ slice restrictions and identifies the
  ambient coordinate partials with the corresponding slice derivatives.  The
  The full boundary-data refinement is supplied by
  `IsParametrizedSurfaceWithBoundary`; its `Nonempty` boundary component keeps
  this predicate proof-irrelevant while exposing the finite parametrization as
  a separate structure.
\end{definition}

\begin{lemma}[smooth dependence of a coordinate flow]
  \label{lem:dc-ch3-2-2-cinfty}
  \dcref{ch3:2.2}
  \uses{lem:dc-ch3-2-2-flow,lem:dc-ch3-2-2-c1dep}
  \lean{Riemannian.GenericFlow.exists_local_flow_contDiffAt,
    Riemannian.GenericFlow.exists_local_flow_joint_contDiff}
  \leanok
  Let $F$ be a finite-dimensional Banach space and let $f:F\to F$ be a globally
  $C^\infty$ autonomous vector field.  At every $z_0\in F$ there are a uniform
  initial-value ball, a fixed time interval, and a family of solutions of
  $z'=f(z)$ which is Lipschitz in the initial value.  The curve-space family and,
  by the autonomous time-state augmentation, its endpoint map are $C^\infty$;
  in particular the solution map is jointly smooth in initial value and time on
  a smaller window.  This is the coordinate-space smooth-dependence core of the
  local-flow theorem; transporting it through a manifold chart remains a
  separate step.
\end{lemma}
\begin{proof}
  \leanok
  Apply Picard--Lindelöf on a ball around $z_0$ to obtain the uniform solution
  family and its Lipschitz estimate.  A compact-ball bound on $\|f'\|$ permits
  choosing $T$ with $T\sup\|f'\|<1$.  The Picard residual is then $C^\infty$ and
  its curve-variable derivative is invertible by the Neumann series, so the
  smooth implicit-function theorem gives the stated $C^\infty$ dependence.
\end{proof}

\begin{lemma}[chart-local uniform geodesic family]
  \label{lem:dc-ch3-2-5-chart-family}
  \dcref{ch3:2.5}
  \lean{Riemannian.GeodesicLocal.geodesic_local_existence,
    Riemannian.GeodesicLocal.geodesic_local_existence_tangentBundle,
    Riemannian.GeodesicLocal.geodesic_local_uniqueness}
  \leanok
  Fix $p\in M$ and a model-space representative $v\in E$ of an initial
  tangent vector in the chart at $p$.  There are $\varepsilon>0$, neighborhoods
  $V_1\in\mathcal N(p)$ and $V_2\in\mathcal N(v)$, and a family
  $c:M\to E\to\mathbb R\to M$ such that, for $q\in V_1$ and $w\in V_2$,
  the curve $c(q,w,\cdot)$ stays in the fixed chart at $p$, solves the
  first-order chart geodesic system on $(-\varepsilon,\varepsilon)$, and has
  $c(q,w,0)=q$ and chart velocity $w$.  Its chart reading is jointly
  $C^\infty$ in $(q,w,t)$.  Two such chart geodesics on open
  order-connected time sets with the same position and chart velocity at one
  common time agree on their overlap.  At the zero tangent vector over $p$,
  the foot and fiber restrictions assemble into a single neighborhood
  $\mathcal U\subset TM$; every $z\in\mathcal U$ seeds one of these curves at
  its foot with the fiber coordinate of $z$ as initial velocity.

  The further refinement of $\mathcal U$ to the explicit base-uniform metric
  disc in $\cref{prop:dc-ch3-2-5}$ is supplied by
  $\cref{lem:dc-ch3-2-5-metric-ball,lem:dc-ch3-2-5-metric-family}$.
\end{lemma}

\begin{lemma}[uniform metric ball in the chart fiber]
  \label{lem:dc-ch3-2-5-metric-ball}
  \dcref{ch3:2.5}
  \uses{lem:dc-ch7-2-8-gram-bound}
  \lean{Riemannian.GeodesicLocal.exists_chartFiberCoord_mem_of_metricInner_lt}
  \leanok
  Let $g$ be a Riemannian metric, let $p\in M$, and let $W$ be a neighborhood of
  $0$ in the model fiber of a chart at $p$.  There are a neighborhood $V$ of $p$
  and a number $\rho>0$ such that, for every $q\in V$ and
  $v\in T_qM$ with $\langle v,v\rangle_q<\rho$, the chart-$p$ coordinate of $v$
  belongs to $W$.
\end{lemma}
\begin{proof}
  \leanok
  Choose $\varepsilon>0$ with the model ball $B_\varepsilon(0)\subset W$.
  The positive-definite Gram form of $g$ in the chart at $p$ dominates a fixed
  positive multiple of the squared model norm on a neighborhood $Y$ of the
  chart image of $p$: there is $c>0$ with
  $\lVert w\rVert^2\le c\,G_y(w,w)$ for $y\in Y$.  Shrink the base neighborhood
  to $V$ so that every $q\in V$ has chart image in $Y$, and put
  $\rho=\varepsilon^2/c$.  If $v\in T_qM$ has squared norm below $\rho$, its
  chart coordinate $w$ satisfies
  $\lVert w\rVert^2\le c\langle v,v\rangle_q<\varepsilon^2$, hence
  $w\in B_\varepsilon(0)\subset W$.
\end{proof}

\begin{lemma}[uniform local geodesic family on a metric velocity ball]
  \label{lem:dc-ch3-2-5-metric-family}
  \dcref{ch3:2.5}
  \uses{lem:dc-ch3-2-5-chart-family,lem:dc-ch3-2-5-metric-ball}
  \lean{Riemannian.GeodesicLocal.geodesic_local_existence_metricBall}
  \leanok
  For every $p\in M$ there are $\delta>0$, a neighborhood $V$ of $p$,
  a model-fiber neighborhood $W$ of $0$, a number $\rho>0$, and a family
  $c:M\times E\times\mathbb R\to M$ such that, whenever $q\in V$ and
  $v\in T_qM$ has $\langle v,v\rangle_q<\rho$, the curve obtained from
  $c(q,\operatorname{coord}_p(v),t)$ is a geodesic on $(-\delta,\delta)$,
  starts at $q$ with the prescribed chart velocity, and is the unique chart
  geodesic with those initial data.  In chart coordinates the family is
  $C^\infty$ on $\bigl(\varphi_p(V)\times W\bigr)\times(-\delta,\delta)$.
\end{lemma}
\begin{proof}
  \leanok
  Apply the chart-local uniform family of
  \cref{lem:dc-ch3-2-5-chart-family} at the zero velocity.  Its model velocity
  neighborhood is a neighborhood of $0$, so
  \cref{lem:dc-ch3-2-5-metric-ball} supplies a smaller base neighborhood and a
  fixed $\rho$ whose metric ball is carried into that model neighborhood.
  Restricting the family and its smoothness statement to the smaller product
  gives the asserted existence and regularity.  If another chart geodesic has
  the same initial position and chart velocity, the phase lifts solve the same
  first-order geodesic system; uniqueness on the connected interval
  $(-\delta,\delta)$ makes the two curves equal.
\end{proof}

\begin{lemma}[$\exp_p$ is a $C^\infty$ diffeomorphism near the origin]
  \label{lem:dc-ch3-2-9-cinfty}
  \dcref{ch3:2.9}
  \uses{lem:dc-ch3-2-9-c2diffeo, lem:dc-ch3-2-9-invertible,
    lem:dc-ch3-2-9-c1diffeo}
  \lean{Riemannian.Exponential.exists_infty_local_diffeomorphism_expMap}
  \leanok
  There is $\varepsilon>0$ such that $\exp_p$ is injective on
  $B_\varepsilon(0)\subset T_pM$, has open image, and its chart reading and
  inverse are $C^\infty$ on the corresponding open images. In particular the
  exponential is a diffeomorphism of a neighbourhood of the origin onto an
  open neighbourhood of $p$.
\end{lemma}
\begin{proof}
  \leanok
  The $C^\infty$ flow-dependence theorem gives a $C^\infty$ chart reading of
  $\exp_p$ on a ball. Shrink this ball inside the strict-derivative and
  injectivity balls. The derivative is then an equivalence at every point, so
  the inverse function theorem makes the chart image open. Define the inverse
  by the unique preimage; at each image point it agrees locally with the
  $C^\infty$ inverse supplied by the inverse function theorem, hence is
  $C^\infty$ on the whole image.
\end{proof}

\begin{lemma}[piecewise curves have length at least distance]
  \label{lem:dc-ch3-3-1-edist}
  \dcref{ch3:3.1}
  \uses{def:dc-ch3-3-1}
  \lean{Riemannian.Geodesic.IsPiecewiseDifferentiableCurve.edist_le_pathELength}
  \leanok
  If $c:[a,b]\to M$ is piecewise differentiable and the ambient metric is the
  Riemannian distance, then
  $$
    d(c(a),c(b))\leq \ell(c|_{[a,b]}).
  $$
  The estimate is applied on each $C^1$ piece and combined across the vertices
  using the triangle inequality and additivity of path length.
\end{lemma}
